% Family and social life — Anglais Tle Tech — Cours signé OnBuch+
% Programme officiel MINESEC. Compiler avec : tectonic N01-family-and-social-life.tex
\newcommand{\DOCMATIERE}{Anglais}
\newcommand{\DOCNIVEAU}{Tle Tech}
\newcommand{\DOCMODULE}{Anglais – classe de Terminale technique}
\newcommand{\DOCLECON}{N01}
\newcommand{\DOCTITRE}{Family and social life}
% OnBuch+ — préambule « cours signés » ENRICHI (compilé avec tectonic / XeLaTeX)
% Système visuel « Pop » d'OnBuch+ : fond crème (jamais blanc), bordures encre
% épaisses, ombres « dures » décalées, titres Archivo Black en capitales.
% Version MATHÉMATIQUES : graphes (pgfplots/tikz), boîtes pédagogiques
% (définition, propriété, méthode, attention, le savais-tu, exercice résolu,
% exercice, TP, bilan), page de garde et signature OnBuch+ ancrée en bas.
%
% Macros attendues dans main.tex AVANT \input{preamble} :
%   \DOCMATIERE  \DOCNIVEAU  \DOCMODULE  \DOCLECON (numéro)  \DOCTITRE
\documentclass[11pt]{article}

\usepackage{fontspec}
\usepackage[french]{babel}
\usepackage[a4paper,margin=2cm,top=3.4cm,bottom=2.8cm,headheight=1.4cm,footskip=1.3cm]{geometry}
\usepackage{amsmath,amssymb,amsfonts}
\usepackage{mathtools} % \xmapsto, \coloneqq... souvent utilisés par les modèles
\usepackage{cancel} % simplifications barrées (bilans de forces)
% \abs{z} (module, valeur absolue) et \norm{u} (norme), utilisés par les modèles
\providecommand{\abs}{}\let\abs\relax\DeclarePairedDelimiter{\abs}{\lvert}{\rvert}
\providecommand{\norm}{}\let\norm\relax\DeclarePairedDelimiter{\norm}{\lVert}{\rVert}
\usepackage[table]{xcolor} % [table] : fournit aussi \rowcolors (alternance de lignes)
\usepackage{pagecolor}
\usepackage{tikz}
\usetikzlibrary{arrows.meta,positioning,calc,shapes.geometric,shapes.misc,shapes.arrows,decorations.pathmorphing,decorations.pathreplacing,decorations.markings,patterns,shadows,mindmap,trees,fit,backgrounds,matrix,intersections,angles,quotes}
\usepackage{pgfplots}
\usepackage[version=4]{mhchem} % équations nucléaires \ce{^{235}_{92}U}
\usepackage[european,siunitx]{circuitikz} % schémas électriques
\pgfplotsset{compat=1.18}
\usepgfplotslibrary{fillbetween,groupplots} % aires entre courbes (calcul intégral)
\usepackage{enumitem}
\usepackage{fancyhdr}
% [most] charge toutes les bibliothèques tcolorbox (listings, minted, raster,
% documentation...) et gonfle inutilement la mémoire TeX sur un document long
% (~300 boîtes) : on ne charge que ce qui est réellement utilisé.
\usepackage[skins,breakable]{tcolorbox}
\usepackage{graphicx}
\usepackage{colortbl}
\usepackage{booktabs}
\usepackage{tabularx}
\usepackage{multirow}
\usepackage{diagbox} % case d'en-tête diagonale des tableaux à double entrée
% Colonnes extensibles centrée / gauche / droite (C, L, R), souvent utilisées par les modèles
\newcolumntype{C}{>{\centering\arraybackslash}X}
\newcolumntype{L}{>{\raggedright\arraybackslash}X}
\newcolumntype{R}{>{\raggedleft\arraybackslash}X}
\usepackage{array}
\usepackage{multicol}
\usepackage{siunitx}
% parse-numbers=false : accepte aussi \SI{2,5 \times 10^{-3}}{...}
\sisetup{locale=FR,output-decimal-marker={,},group-minimum-digits=4,parse-numbers=false}
\DeclareSIUnit\ohm{\ensuremath{\symup{\Omega}}} % U+2126 absent de la police
\DeclareSIUnit\division{div} % réglages d'oscilloscope (ms/div, V/div)
\DeclareSIUnit\tr{tr} % tours (vitesse de rotation, ex. \SI{1500}{\tr\per\minute})
\DeclareSIUnit\molar{M} % concentration molaire (ex. \SI{3}{\molar}, \SI{1}{\micro\molar})
\DeclareSIUnit\an{an} % année (le modèle confond parfois avec \year, un compteur TeX) — ex. \SI{5}{\kilo\meter\per\an}
\DeclareSIUnit\FCFA{FCFA} % franc CFA (ex. \SI{500}{\FCFA\per\kilo\gram})
\DeclareSIUnit\degreeBrix{\textdegree Brix} % teneur en sucre d'un sirop/jus (réfractométrie)
\DeclareSIUnit\month{mois} % le modèle utilise parfois \month, un compteur TeX, comme unité de durée
\DeclareSIUnit\microliter{\micro\liter} % le modèle écrit parfois \microliter en un seul token
\DeclareSIUnit\million{millions} % dénombrement (ex. \SI{15}{\million\per\mL} spermatozoïdes)
\usepackage{qrcode}
% \degree : commande fréquemment utilisée par les modèles pour un angle ou une
% température en dehors de \SI{}{} (non fournie par les paquets chargés).
\providecommand{\degree}{^{\circ}}
% \qed (fin de démonstration), utilisé par les modèles sans amsthm
\providecommand{\qed}{\hfill$\square$}
% \barre : double barre des valeurs interdites dans un tableau de variations (tkz-tab)
\providecommand{\barre}{\ensuremath{\|}}
% environnement proof (amsthm non chargé) utilisé par les modèles
\ifdefined\proof\else\newenvironment{proof}[1][Démonstration]{\par\noindent\textit{#1.}\ }{\qed\par}\fi
\usepackage{lastpage}
\usepackage{hyperref}
\hypersetup{hidelinks}

% ============================================================
% Palette « Pop » OnBuch+ — mêmes hex que la classe P (plus_ui.dart)
% ============================================================
\definecolor{poporange}{HTML}{F59321}
\definecolor{popdark}{HTML}{E07A0C}
\definecolor{popcream}{HTML}{FFF6E8}
\definecolor{popink}{HTML}{1C1714}
\definecolor{poppurple}{HTML}{7A5AE0}
\definecolor{popgreen}{HTML}{2E9E6A}
\definecolor{poppink}{HTML}{E0457B}
\definecolor{popblue}{HTML}{2F7DE1}
\definecolor{popgold}{HTML}{C98A12}
\definecolor{popmuted}{HTML}{8A7F76}
\definecolor{popline}{HTML}{EADFD2}
\definecolor{popgray}{HTML}{8A7F76}
\colorlet{popgrey}{popgray}
% Alias tolérants (le modèle nomme parfois les couleurs différemment)
\colorlet{popwhite}{white}
\colorlet{popblack}{popink}
\colorlet{popred}{poppink}
\colorlet{popyellow}{popgold}
% Teintes claires (fonds de tableaux/encadrés) — le modèle nomme parfois
% les couleurs avec un suffixe L (ex. popblueL) sans les avoir définies.
\colorlet{poporangeL}{poporange!20}
\colorlet{popdarkL}{popdark!20}
\colorlet{poppurpleL}{poppurple!20}
\colorlet{popgreenL}{popgreen!20}
\colorlet{poppinkL}{poppink!20}
\colorlet{popblueL}{popblue!20}
\colorlet{popgoldL}{popgold!20}
\colorlet{popredL}{popred!20}
\colorlet{popgrayL}{popgray!20}
% Teintes claires pour fonds de tableaux / zones
\colorlet{popblueL}{popblue!10!white}
\colorlet{poporangeL}{poporange!14!white}
\colorlet{popgreenL}{popgreen!12!white}
\colorlet{poppurpleL}{poppurple!12!white}
\colorlet{poppinkL}{poppink!12!white}
\colorlet{popgoldL}{popgold!14!white}

\pagecolor{popcream}

% ============================================================
% Typographie
% ============================================================
\defaultfontfeatures{Ligatures=TeX}
\setmainfont{PlusJakartaSans-Medium.ttf}[Path=../../../fonts/,BoldFont=PlusJakartaSans-Bold.ttf]
% Maths en sans-serif (Fira Math) avec chiffres et lettres droites de Plus
% Jakarta, pour que les formules se fondent dans le texte.
\usepackage[math-style=ISO,bold-style=ISO]{unicode-math}
\setmathfont{FiraMath-Regular.otf}[Path=../../../fonts/]
\setmathfont{PlusJakartaSans-Medium.ttf}[Path=../../../fonts/,range={up/num,up/{Latin,latin},\mathup}]
\usepackage{icomma} % virgule décimale sans espace en mode math
% Caractères mathématiques Unicode absents des polices de texte (Plus Jakarta,
% Archivo Black) : rendus par leur équivalent mathématique, en texte comme en
% formule — sinon ils s'affichent en carré vide (ex. « Arithmétique dans ℤ »).
\usepackage{newunicodechar}
\newunicodechar{ℝ}{\ensuremath{\mathbb{R}}}
\newunicodechar{ℤ}{\ensuremath{\mathbb{Z}}}
\newunicodechar{ℂ}{\ensuremath{\mathbb{C}}}
\newunicodechar{ℕ}{\ensuremath{\mathbb{N}}}
\newunicodechar{ℚ}{\ensuremath{\mathbb{Q}}}
\newunicodechar{𝔻}{\ensuremath{\mathbb{D}}}
\newunicodechar{α}{\ensuremath{\alpha}}
\newunicodechar{β}{\ensuremath{\beta}}
\newunicodechar{γ}{\ensuremath{\gamma}}
\newunicodechar{δ}{\ensuremath{\delta}}
\newunicodechar{ε}{\ensuremath{\varepsilon}}
\newunicodechar{θ}{\ensuremath{\theta}}
\newunicodechar{λ}{\ensuremath{\lambda}}
\newunicodechar{μ}{\ensuremath{\mu}}
\newunicodechar{σ}{\ensuremath{\sigma}}
\newunicodechar{τ}{\ensuremath{\tau}}
\newunicodechar{φ}{\ensuremath{\varphi}}
\newunicodechar{ψ}{\ensuremath{\psi}}
\newunicodechar{ω}{\ensuremath{\omega}}
\newunicodechar{Σ}{\ensuremath{\Sigma}}
% majuscules produites par \MakeUppercase dans les titres (α-aminés -> Α)
\newunicodechar{Α}{\ensuremath{\alpha}}
\newunicodechar{Β}{\ensuremath{\beta}}
\newunicodechar{Γ}{\ensuremath{\Gamma}}
\newunicodechar{Δ}{\ensuremath{\Delta}}
\newunicodechar{Ε}{\ensuremath{\varepsilon}}
\newunicodechar{Θ}{\ensuremath{\Theta}}
\newunicodechar{Λ}{\ensuremath{\Lambda}}
\newunicodechar{Μ}{\ensuremath{\mu}}
\newunicodechar{Τ}{\ensuremath{\tau}}
\newunicodechar{Φ}{\ensuremath{\Phi}}
\newunicodechar{Ψ}{\ensuremath{\Psi}}
\newunicodechar{Ω}{\ensuremath{\Omega}}
\newunicodechar{↦}{\ensuremath{\mapsto}}
\newunicodechar{∈}{\ensuremath{\in}}
\newunicodechar{∉}{\ensuremath{\notin}}
\newunicodechar{∧}{\ensuremath{\wedge}}
\newunicodechar{⇔}{\ensuremath{\Leftrightarrow}}
\newunicodechar{⟺}{\ensuremath{\Longleftrightarrow}}
\newunicodechar{⇒}{\ensuremath{\Rightarrow}}
\newunicodechar{⟹}{\ensuremath{\Longrightarrow}}
\newunicodechar{∘}{\ensuremath{\circ}}
\newunicodechar{∪}{\ensuremath{\cup}}
\newunicodechar{∩}{\ensuremath{\cap}}
\newunicodechar{≡}{\ensuremath{\equiv}}
\newunicodechar{⊂}{\ensuremath{\subset}}
\newunicodechar{⊥}{\ensuremath{\perp}}
\newunicodechar{∃}{\ensuremath{\exists}}
\newunicodechar{∀}{\ensuremath{\forall}}
\newunicodechar{∛}{\ensuremath{\sqrt[3]{\,}}}
\newunicodechar{∜}{\ensuremath{\sqrt[4]{\,}}}
\newunicodechar{⃗}{\ensuremath{{}^{\to}}}
\newunicodechar{ⁿ}{\ifmmode^{n}\else\textsuperscript{\ensuremath{n}}\fi}
\newunicodechar{ˣ}{\ifmmode^{x}\else\textsuperscript{\ensuremath{x}}\fi}
\newunicodechar{ᵇ}{\ifmmode^{b}\else\textsuperscript{\ensuremath{b}}\fi}
\newunicodechar{ʳ}{\ifmmode^{r}\else\textsuperscript{\ensuremath{r}}\fi}
\newunicodechar{ᵘ}{\ifmmode^{u}\else\textsuperscript{\ensuremath{u}}\fi}
\newunicodechar{ʸ}{\ifmmode^{y}\else\textsuperscript{\ensuremath{y}}\fi}
\newunicodechar{ᵃ}{\ifmmode^{a}\else\textsuperscript{\ensuremath{a}}\fi}
\newunicodechar{ᵉ}{\ifmmode^{e}\else\textsuperscript{\ensuremath{e}}\fi}
\newunicodechar{ᵏ}{\ifmmode^{k}\else\textsuperscript{\ensuremath{k}}\fi}
\newunicodechar{ᵖ}{\ifmmode^{p}\else\textsuperscript{\ensuremath{p}}\fi}
\newunicodechar{ᴺ}{\ifmmode^{N}\else\textsuperscript{\ensuremath{N}}\fi}
\newunicodechar{ᶜ}{\ifmmode^{c}\else\textsuperscript{\ensuremath{c}}\fi}
\newunicodechar{ʰ}{\ifmmode^{h}\else\textsuperscript{\ensuremath{h}}\fi}
\newunicodechar{ᵠ}{\ifmmode^{q}\else\textsuperscript{\ensuremath{q}}\fi}
\newunicodechar{ₙ}{\ifmmode_{n}\else\textsubscript{\ensuremath{n}}\fi}
\newunicodechar{ₐ}{\ifmmode_{a}\else\textsubscript{\ensuremath{a}}\fi}
\newunicodechar{ᵢ}{\ifmmode_{i}\else\textsubscript{\ensuremath{i}}\fi}
\newunicodechar{ᵣ}{\ifmmode_{r}\else\textsubscript{\ensuremath{r}}\fi}
\newunicodechar{ₖ}{\ifmmode_{k}\else\textsubscript{\ensuremath{k}}\fi}
\newunicodechar{ᵦ}{\ifmmode_{\beta}\else\textsubscript{\ensuremath{\beta}}\fi}
\newunicodechar{ₓ}{\ifmmode_{x}\else\textsubscript{\ensuremath{x}}\fi}
\newfontfamily\popdisplay{ArchivoBlack-Regular.ttf}[Path=../../../fonts/]
\newfontfamily\poplogo{SpaceGrotesk-Bold.ttf}[Path=../../../fonts/]
\newfontfamily\popsemi{PlusJakartaSans-SemiBold.ttf}[Path=../../../fonts/]
\newfontfamily\popxbold{PlusJakartaSans-ExtraBold.ttf}[Path=../../../fonts/]
\newfontfamily\popxboldls{PlusJakartaSans-ExtraBold.ttf}[Path=../../../fonts/,LetterSpace=3.0]

\setlist[enumerate]{leftmargin=1.7em,itemsep=5pt,topsep=4pt}
\setlist[itemize]{leftmargin=1.7em,itemsep=4pt,topsep=4pt,label={\color{poporange}\textbullet}}
\setlength{\parindent}{0pt}
\setlength{\parskip}{5pt}
\renewcommand{\arraystretch}{1.25}

% pgfplots : style Pop par défaut
\pgfplotsset{
  every axis/.append style={axis line style={popink,line width=1pt},tick style={popink},
    grid style={popline,line width=0.5pt},label style={font=\small},tick label style={font=\footnotesize}},
  popcurve/.style={poporange,line width=1.8pt,no markers,smooth},
}
% Même style disponible dans un \draw TikZ simple (hors pgfplots)
\tikzset{popcurve/.style={poporange,line width=1.8pt}}
\tikzset{popgrid/.style={popline,very thin}}
\colorlet{popcurve}{poporange} % le nom du style est parfois utilisé comme couleur

% ============================================================
% Pill / squiggle
% ============================================================
\newcommand{\poppill}[3]{%
  \tikz[baseline=-0.55ex]\node[draw=popink,line width=1.2pt,rounded corners=2.6pt,%
    fill=#2,inner xsep=6.5pt,inner ysep=3pt]{%
    \popxboldls\fontsize{7}{7}\selectfont\color{#3}\MakeUppercase{#1}};%
}
\newcommand{\popsquiggle}[1]{%
  \tikz[baseline=-0.4ex]{
    \draw[#1,line width=2.6pt,line cap=round]
      plot [smooth] coordinates {(0,0.09) (0.34,0.01) (0.68,0.21) (1.02,0.01) (1.36,0.10)};
  }%
}
% Mot-clé surligné (marqueur orange sous le texte)
\newcommand{\cle}[1]{{\popxbold\color{popdark}#1}}

% ============================================================
% En-tête / pied
% ============================================================
\pagestyle{fancy}
\fancyhf{}
\renewcommand{\headrulewidth}{0pt}
\renewcommand{\footrulewidth}{0pt}
\lhead{\raisebox{-0.15ex}{\poplogo\fontsize{13}{13}\selectfont\color{popink}OnBuch\color{poporange}+}}
\rhead{\poppill{\DOCMATIERE\ \textbullet\ \DOCNIVEAU\ \textbullet\ Leçon \DOCLECON}{white}{popink}}
\lfoot{\popxbold\fontsize{7.6}{7.6}\selectfont\color{popmuted}\MakeUppercase{Cours signé OnBuch+}\ \textbullet\ {\popsemi\DOCTITRE}}
\rfoot{\tikz[baseline=-0.6ex]\node[rounded corners=8.5pt,fill=popink,minimum height=17pt,minimum width=17pt,inner xsep=5pt,inner sep=0]{\popxbold\fontsize{8}{8}\selectfont\color{popcream}\thepage\,/\,\pageref*{LastPage}};}

% ============================================================
% Titres
% ============================================================
\newcommand{\chaptitre}[2]{%
  \begin{tcolorbox}[enhanced,colback=poporange,colframe=popink,boxrule=2.6pt,arc=11pt,
      drop shadow=popink,left=15pt,right=15pt,top=12pt,bottom=11pt,before skip=3pt,after skip=18pt]
    \poppill{#2}{popink}{popcream}\par\vspace{9pt}
    {\raggedright\hyphenpenalty=10000\exhyphenpenalty=10000\popdisplay\fontsize{22}{24}\selectfont\color{popcream}\MakeUppercase{#1}\par}
  \end{tcolorbox}
}
% Section numérotée : pastille orange avec le numéro + titre Archivo
\newcounter{popsec}
\newcommand{\coursec}[1]{%
  \stepcounter{popsec}\setcounter{popsub}{0}%
  \par\vspace{16pt}\noindent
  \begin{minipage}{\linewidth}
  \tikz[baseline=-0.7ex]\node[draw=popink,line width=1.6pt,fill=poporange,rounded corners=4pt,minimum size=22pt,inner sep=0]{\popdisplay\fontsize{12}{12}\selectfont\color{popcream}\thepopsec};%
  \hspace{8pt}{\popdisplay\fontsize{15}{17}\selectfont\color{popink}\MakeUppercase{#1}}\par
  \vspace{4pt}\noindent{\color{poporange}\rule{36pt}{3.6pt}}
  \end{minipage}\par\nopagebreak\vspace{8pt}
}
\newcounter{popsub}
\newcommand{\courssub}[1]{%
  \stepcounter{popsub}%
  \par\vspace{9pt}\noindent{\popxbold\fontsize{11.5}{13}\selectfont\color{popink}{\color{poporange}\thepopsec.\thepopsub}\hspace{6pt}#1}\par\nopagebreak\vspace{4pt}
}

% ============================================================
% Boîtes pédagogiques — même recette : fond blanc, bordure épaisse colorée,
% ombre dure encre, badge plein. Titre optionnel [#1].
% ============================================================
\tcbset{popbase/.style={breakable,enhanced,colback=white,boxrule=2.3pt,arc=9pt,
  shadow={3pt}{-3pt}{0pt}{popink},left=14pt,right=14pt,top=11pt,bottom=11pt,before skip=11pt,after skip=15pt,
  fontupper=\popsemi\fontsize{10.6}{14.5}\selectfont\color{popink},
  fontlower=\popsemi\fontsize{10.6}{14.5}\selectfont\color{popink}}}
% Coche dessinée (le glyphe ✓ n'existe pas dans les polices du document)
\newcommand{\popcheck}[1]{\tikz[baseline=-0.5ex]\draw[#1,line width=1.6pt,line cap=round,line join=round](-0.13,0.0)--(-0.04,-0.09)--(0.13,0.1);}
\providecommand{\coche}{\popcheck{popgreen}}
\providecommand{\resultat}[1]{\par\smallskip\noindent\fcolorbox{popgreen}{popgreenL}{\parbox{\dimexpr\linewidth-2\fboxsep-2\fboxrule\relax}{\textbf{Résultat.} #1}}\par\smallskip}
\newcommand{\popboxhead}[3]{\poppill{#1}{#2}{white}\ifx\relax#3\relax\else\hspace{6pt}{\popxbold\fontsize{10.6}{12}\selectfont\color{popink}#3}\fi\par\vspace{7pt}}

\newtcolorbox{aretenir}[1][]{popbase,colframe=popgreen,before upper={\popboxhead{À retenir}{popgreen}{#1}}}
\newtcolorbox{exemplebox}[1][]{popbase,colframe=poporange,before upper={\popboxhead{Exemple}{poporange}{#1}}}
\newtcolorbox{definition}[1][]{popbase,colframe=popblue,before upper={\popboxhead{Définition}{popblue}{#1}}}
\newtcolorbox{propriete}[1][]{popbase,colframe=poppurple,before upper={\popboxhead{Propriété}{poppurple}{#1}}}
\newtcolorbox{methode}[1][]{popbase,colframe=popgold,before upper={\popboxhead{Méthode}{popgold}{#1}}}
\newtcolorbox{attention}[1][]{popbase,colframe=poppink,before upper={\popboxhead{Attention}{poppink}{#1}}}
\newtcolorbox{experience}[1][]{popbase,colframe=popblue,colback=popblueL,before upper={\popboxhead{Expérience}{popblue}{#1}}}
% « Le savais-tu ? » : carte violette pleine (contexte, culture, Cameroun)
\newtcolorbox{savaistu}[1][]{popbase,colback=poppurple,colframe=popink,
  fontupper=\popsemi\fontsize{10.4}{14.2}\selectfont\color{white},
  before upper={\poppill{Le savais-tu\,?}{white}{poppurple}\ifx\relax#1\relax\else\hspace{6pt}{\popxbold\color{white}#1}\fi\par\vspace{7pt}}}
% Exercice résolu : énoncé en haut, \tcblower puis la solution sur fond crème
\newcounter{popexo}\newcounter{popexr}
\newtcolorbox{exoresolu}[1][]{popbase,colframe=poporange,colbacklower=popcream,
  segmentation style={popink,line width=1.2pt,dashed},
  before upper={\stepcounter{popexr}\popboxhead{Exercice résolu \thepopexr}{poporange}{#1}},
  before lower={\poppill{Solution}{popink}{popcream}\par\vspace{6pt}}}
% Exercice (non résolu en place) — la correction est dans la partie Corrigés
\newtcolorbox{exercice}[1][]{popbase,colframe=popink,
  before upper={\stepcounter{popexo}\popboxhead{Exercice \thepopexo}{popink}{#1}}}
% Corrigé
\newtcolorbox{corrige}[1][]{popbase,colframe=popgreen,colback=popgreenL,
  before upper={\popboxhead{Corrigé}{popgreen}{#1}}}
% Objectifs / prérequis / bilan
\newtcolorbox{objectifs}{popbase,colframe=popink,colback=poporangeL,
  before upper={\popboxhead{Objectifs de la leçon}{popink}{}}}
\newtcolorbox{prerequis}{popbase,colframe=popink,colback=popblueL,
  before upper={\popboxhead{Prérequis}{popblue}{}}}
\newtcolorbox{bilan}{popbase,colframe=popink,colback=white,boxrule=2.8pt,
  before upper={\popboxhead{Fiche bilan}{poporange}{L'essentiel à connaître pour l'examen}}}
% Figure encadrée avec légende
\newtcolorbox{popfigure}[1][]{enhanced,breakable=false,colback=white,colframe=popline,boxrule=1.4pt,arc=8pt,
  left=10pt,right=10pt,top=10pt,bottom=8pt,before skip=10pt,after skip=12pt,halign=center,
  lower separated=false,halign lower=center,
  fontlower=\popsemi\fontsize{9}{11.5}\selectfont\color{popmuted},
  #1}
\newcommand{\legende}[1]{\par\vspace{6pt}{\popsemi\fontsize{9}{11.5}\selectfont\color{popmuted}#1}}

% ============================================================
% Signature OnBuch+
% ============================================================
\newcommand{\poplogoinline}[1]{{\poplogo\fontsize{#1}{#1}\selectfont\color{popink}OnBuch\color{poporange}+}}
\newcommand{\signatureonbuch}{%
  \begin{tcolorbox}[enhanced,colback=popink,colframe=popink,boxrule=2.2pt,arc=10pt,
      drop shadow=poporange,left=14pt,right=14pt,top=10pt,bottom=10pt,before skip=0pt,after skip=0pt]
    \tikz[baseline=-0.7ex]\node[circle,fill=popgreen,draw=popcream,line width=1.3pt,minimum size=22pt,inner sep=0]{\popcheck{white}};%
    \hspace{9pt}\begin{minipage}[c]{0.62\linewidth}
      {\popxbold\fontsize{12}{13}\selectfont\color{popcream}Signé\ \textbullet\ {\poplogo OnBuch\color{poporange}+}}\par\vspace{2pt}
      {\raggedright\popsemi\fontsize{8}{10.5}\selectfont\color{popline}\MakeUppercase{Équipe pédagogique OnBuch+}\\\MakeUppercase{Conforme au programme officiel MINESEC}\par}
    \end{minipage}\hfill
    {\poplogo\fontsize{22}{22}\selectfont\color{popcream}OnBuch\color{poporange}+}
  \end{tcolorbox}%
}

% ============================================================
% Page de garde
% #1 titre · #2 matière · #3 niveau · #4 module · #5 n° leçon · #6 durée ·
% #7 description / objectifs (texte ou itemize)
% ============================================================
\newcommand{\pagegarde}[7]{%
  \thispagestyle{empty}
  \newgeometry{a4paper,margin=1.7cm,top=1.6cm,bottom=1.4cm}%
  \begin{tikzpicture}[remember picture,overlay]
    \fill[poporange!18] ([xshift=-3.2cm,yshift=-3.6cm]current page.north east) circle (4.6cm);
    \fill[poppurple!14] ([xshift=2.4cm,yshift=7.6cm]current page.south west) circle (3.8cm);
  \end{tikzpicture}%
  {\poplogo\fontsize{30}{30}\selectfont\color{popink}OnBuch\color{poporange}+}\hfill
  \raisebox{0.6ex}{\poppill{Cours signé \textbullet\ Premium}{popink}{popcream}}\par
  \vspace{1pt}\popsquiggle{poppurple}\par
  \vspace{8pt}
  \begin{tcolorbox}[enhanced,colback=poporange,colframe=popink,boxrule=2.8pt,arc=13pt,
      drop shadow=popink,left=18pt,right=18pt,top=12pt,bottom=13pt,after skip=10pt]
    \poppill{#2 \textbullet\ #3}{popink}{popcream}\hspace{5pt}\poppill{#4}{white}{popink}\par\vspace{8pt}
    {\popdisplay\fontsize{14}{14}\selectfont\color{popink}LEÇON #5}\par\vspace{4pt}
    {\raggedright\hyphenpenalty=10000\exhyphenpenalty=10000\popdisplay\fontsize{25}{27}\selectfont\color{popcream}\MakeUppercase{#1}\par}
    \vspace{8pt}
    {\popxbold\fontsize{9.5}{9.5}\selectfont\color{popink}\MakeUppercase{Durée conseillée : #6}}
  \end{tcolorbox}
  \begin{tcolorbox}[enhanced,colback=white,colframe=popink,boxrule=2.2pt,arc=10pt,
      drop shadow=popink,left=16pt,right=16pt,top=10pt,bottom=10pt,after skip=8pt]
    \poppill{Dans cette leçon}{poppurple}{white}\par\vspace{5pt}
    \popsemi\fontsize{9.3}{12.6}\selectfont\color{popink}#7
  \end{tcolorbox}
  \noindent\begin{minipage}[t]{0.487\linewidth}
    \begin{tcolorbox}[enhanced,colback=popgreen,colframe=popink,boxrule=2.2pt,arc=10pt,
        drop shadow=popink,left=11pt,right=11pt,top=10pt,bottom=10pt,height=3.1cm,valign=center]
      \tcbox[colback=white,colframe=popink,boxrule=1.3pt,arc=5pt,boxsep=0pt,nobeforeafter,
        left=3pt,right=3pt,top=3pt,bottom=3pt,valign=center]{\qrcode[height=1.9cm]{https://play.google.com/store/apps/details?id=com.luvvix.onbuch&hl=fr}}%
      \hspace{8pt}\begin{minipage}[c]{0.5\linewidth}
        {\popdisplay\fontsize{11}{12.5}\selectfont\color{white}\MakeUppercase{Télécharge OnBuch}}\par\vspace{4pt}
        {\raggedright\popsemi\fontsize{8.4}{11}\selectfont\color{white}Gratuit \textbullet\ Google Play\\Tuteur IA, annales, concours, résultats\par}
      \end{minipage}
    \end{tcolorbox}
  \end{minipage}\hfill
  \begin{minipage}[t]{0.487\linewidth}
    \begin{tcolorbox}[enhanced,colback=poppurple,colframe=popink,boxrule=2.2pt,arc=10pt,
        drop shadow=popink,left=11pt,right=11pt,top=10pt,bottom=10pt,height=3.1cm,valign=center]
      \tikz[baseline=-0.7ex]\node[circle,fill=white,draw=popink,line width=1.3pt,minimum size=1.6cm,inner sep=0]{\popdisplay\fontsize{24}{24}\selectfont\color{poppurple}+};%
      \hspace{8pt}\begin{minipage}[c]{0.6\linewidth}
        {\popdisplay\fontsize{11}{12.5}\selectfont\color{white}\MakeUppercase{Passe à OnBuch+}}\par\vspace{4pt}
        {\raggedright\popsemi\fontsize{8.4}{11}\selectfont\color{white}Tous les cours signés, Tuteur IA illimité, fiches PDF — onglet OnBuch+ de l'app\par}
      \end{minipage}
    \end{tcolorbox}
  \end{minipage}
  \vspace{8pt}
  \popcontenu
  \vfill
  \signatureonbuch
  \restoregeometry
  \newpage
}

% Rangée « Ce cours contient » (page de garde)
\newcommand{\poptile}[3]{% #1 couleur #2 gros texte #3 légende
  \begin{tcolorbox}[enhanced,colback=#1,colframe=popink,boxrule=2pt,arc=9pt,shadow={2.5pt}{-2.5pt}{0pt}{popink},
      left=6pt,right=6pt,top=9pt,bottom=9pt,halign=center,valign=center,height=1.75cm,width=\linewidth,nobeforeafter]
    {\popdisplay\fontsize{12.5}{13}\selectfont\color{white}\MakeUppercase{#2}}\par\vspace{3pt}
    {\centering\popsemi\fontsize{7.8}{9.5}\selectfont\color{white}#3\par}
  \end{tcolorbox}}
\newcommand{\popcontenu}{%
  \par\noindent{\popxboldls\fontsize{7.5}{7.5}\selectfont\color{popmuted}CE COURS SIGNÉ CONTIENT}\par\vspace{7pt}
  \noindent\begin{minipage}[t]{0.235\linewidth}\poptile{popblue}{Cours}{complet, illustré, pas à pas}\end{minipage}\hfill
  \begin{minipage}[t]{0.235\linewidth}\poptile{poporange}{Méthodes}{et exercices résolus}\end{minipage}\hfill
  \begin{minipage}[t]{0.235\linewidth}\poptile{poppink}{Exercices}{type examen + corrigés}\end{minipage}\hfill
  \begin{minipage}[t]{0.235\linewidth}\poptile{popgreen}{Fiche bilan}{l'essentiel à retenir}\end{minipage}\par}

% Sommaire
\newtcolorbox{sommaire}{enhanced,colback=white,colframe=popink,boxrule=2.2pt,arc=10pt,
  drop shadow=popink,left=18pt,right=18pt,top=13pt,bottom=13pt,before skip=6pt,after skip=20pt,
  before upper={\poppill{Sommaire}{poppurple}{white}\par\vspace{9pt}\popsemi\fontsize{10.6}{14.5}\selectfont\color{popink}}}

% Signature de fin de cours — poussée tout en bas de la dernière page
\newcommand{\findecours}{%
  \par\vfill\nopagebreak
  \begin{center}\popsquiggle{poporange}\end{center}\vspace{4pt}
  \signatureonbuch
}
\AtBeginDocument{\def\llbracket{\mathopen{[\mkern-3.5mu[}}\def\rrbracket{\mathclose{]\mkern-3.5mu]}}} % intervalles d'entiers (glyphe ⟦ absent de la police)
% Glyphes absents de Fira Math : redéfinis à partir de glyphes disponibles
\AtBeginDocument{%
  \def\setminus{\mathbin{\backslash}}\let\smallsetminus\setminus
  \def\vdots{\mathord{\vbox{\baselineskip3.2pt\lineskiplimit0pt\kern4pt\hbox{.}\hbox{.}\hbox{.}}}}%
  \def\ddots{\mathinner{\mkern1mu\raise6.5pt\hbox{.}\mkern2mu\raise3.6pt\hbox{.}\mkern2mu\raise0.7pt\hbox{.}\mkern1mu}}%
  \def\triangle{\mathord{\scalebox{0.9}{$\symup{\Delta}$}}}%
  \def\nabla{\mathord{\raisebox{\depth}{\scalebox{1}[-1]{$\symup{\Delta}$}}}}%
  \def\checkmark{\popcheck{popdark}}%
  \def\star{\mathbin{\ast}}\def\bigstar{\mathord{\ast}}%
  \def\diamond{\mathbin{\scalebox{0.6}{$\Diamond$}}}%
  \def\bigcup{\mathop{\mathchoice{\raisebox{-0.3em}{\scalebox{1.7}{$\cup$}}}{\scalebox{1.25}{$\cup$}}{\cup}{\cup}}\displaylimits}%
  \def\bigcap{\mathop{\mathchoice{\raisebox{-0.3em}{\scalebox{1.7}{$\cap$}}}{\scalebox{1.25}{$\cap$}}{\cap}{\cap}}\displaylimits}%
  \def\lesseqqgtr{\mathrel{\lessgtr}}\def\bigoplus{\mathop{\oplus}\displaylimits}%
}
\providecommand{\ssi}{si et seulement si} % abréviation fréquente
\AtBeginDocument{\providecommand{\wideparen}{\overparen}} % arc de cercle

%% FIGFIX-BEGIN v5 — correctifs globaux des figures (docs/onbuch-generation/tools/figfix)
\usepackage{adjustbox}\usepackage{etoolbox}\tcbuselibrary{raster}
% toute figure TikZ/pgfplots plus large que la ligne est réduite à la largeur disponible
\BeforeBeginEnvironment{tikzpicture}{\begin{adjustbox}{max width=\linewidth}}
\AfterEndEnvironment{tikzpicture}{\end{adjustbox}}
% légendes pgfplots lisibles par-dessus les courbes
\pgfplotsset{every axis/.append style={legend style={fill=white,fill opacity=0.92,text opacity=1,draw=black!15,inner sep=3pt}}}
% étiquettes d'arêtes (syntaxe edge["..."]) avec fond blanc
\tikzset{every edge quotes/.append style={fill=white,inner sep=1.2pt}}
%% FIGFIX-END
\begin{document}

\pagegarde{Family and social life}{Anglais}{Tle Tech}{Anglais – classe de Terminale technique}{N01}{13 séances (fiche de progression harmonisée)}{Cette leçon prépare les élèves de Terminale technique à réussir un entretien d'embauche en anglais, à parler de leurs activités de loisirs et à anticiper leurs besoins en langue anglaise dans le monde professionnel camerounais. Elle mobilise des outils grammaticaux clés (present perfect continuous, future continuous, embedded questions, question tags, phrasal verbs) à travers des situations concrètes de la vie courante.
\vspace{4pt}
\begin{multicols}{2}\small
\begin{itemize}
\item À la fin de cette leçon, je sais donner et recevoir des instructions pour me préparer à un entretien d'embauche en anglais.
\item À la fin de cette leçon, je sais utiliser le vocabulaire lié aux entretiens d'embauche (CV, qualifications, strengths, weaknesses...).
\item À la fin de cette leçon, je sais comprendre un texte oral sur les activités de loisirs et utiliser le vocabulaire approprié.
\item À la fin de cette leçon, je sais employer correctement le present perfect continuous et le future continuous.
\item À la fin de cette leçon, je sais construire des phrases et questions indirectes (embedded statements/questions).
\item À la fin de cette leçon, je sais utiliser les question tags pour confirmer une information.
\end{itemize}\end{multicols}}

\begin{sommaire}
\begin{multicols}{2}
\begin{enumerate}
\item Undergoing a job interview (Speaking \& Vocabulary)
\item Recreational activities (Listening \& Vocabulary)
\item Tenses: present perfect continuous \& future continuous
\item Embedded questions, question tags \& phrasal verbs
\item Anticipating second-language needs at work (Reading \& Vocabulary)
\item Writing: composition on second-language needs at work
\item Activité d'intégration
\item Méthodes et exercices résolus
\item Exercices
\item Corrigés des exercices
\item Fiche bilan
\end{enumerate}
\end{multicols}
\end{sommaire}

\begin{objectifs}
\begin{itemize}
\item À la fin de cette leçon, je sais donner et recevoir des instructions pour me préparer à un entretien d'embauche en anglais.
\item À la fin de cette leçon, je sais utiliser le vocabulaire lié aux entretiens d'embauche (CV, qualifications, strengths, weaknesses...).
\item À la fin de cette leçon, je sais comprendre un texte oral sur les activités de loisirs et utiliser le vocabulaire approprié.
\item À la fin de cette leçon, je sais employer correctement le present perfect continuous et le future continuous.
\item À la fin de cette leçon, je sais construire des phrases et questions indirectes (embedded statements/questions).
\item À la fin de cette leçon, je sais utiliser les question tags pour confirmer une information.
\item À la fin de cette leçon, je sais employer les phrasal verbs et les prepositional phrases dans des situations réelles.
\item À la fin de cette leçon, je sais rédiger une composition sur l'anticipation de mes besoins en anglais au travail.
\item À la fin de cette leçon, je sais réviser et corriger mes points faibles en grammaire.
\end{itemize}
\end{objectifs}

\begin{prerequis}
\begin{itemize}
\item Les temps de base acquis en Première : present simple, present continuous, past simple, future simple (will/going to).
\item Le vocabulaire général du monde du travail vu en Première (jobs, occupations, workplaces).
\item La formation des questions directes (inversion du sujet et de l'auxiliaire, mots interrogatifs).
\item Les auxiliaires do/does/did, have/has, be et les modaux (can, must, should).
\item La rédaction d'un paragraphe simple et d'une lettre formelle (vues en Seconde et Première).
\end{itemize}
\end{prerequis}

\newpage

\coursec{Undergoing a job interview (Speaking \& Vocabulary)}

\courssub{Hook situation \& problem statement}

\begin{exemplebox}[Situation d'accroche]
Ariane, élève en Terminale Génie Civil au Lycée Technique de Douala, vient de décrocher un entretien de stage chez \textbf{SOGEA SATOM} à Bonabéri. L'entreprise exige qu'elle se présente en anglais, car ses chantiers accueillent des partenaires nigérians et ghanéens. Pendant son temps libre, elle joue au football et écoute de la musique pour se détendre avant le grand jour. \textbf{Comment doit-elle se préparer pour réussir son entretien, parler de ses loisirs et anticiper ses besoins en anglais au travail ?} Cette leçon l'accompagne pas à pas.
\end{exemplebox}

\courssub{What is a job interview?}

\begin{definition}[Job interview]
A \cle{job interview} is a formal meeting where an employer asks questions to decide if a candidate is suitable for a job. En français : un entretien d'embauche.
\end{definition}

\begin{aretenir}[Key terms]
\begin{itemize}
    \item \cle{applicant} / \cle{candidate} : la personne qui postule (candidat)
    \item \cle{interviewer} : la personne qui mène l'entretien (recruteur)
    \item \cle{interviewee} : la personne qui passe l'entretien (candidat face au recruteur)
    \item \cle{CV} / \cle{résumé} : curriculum vitae (CV)
    \item \cle{cover letter} : lettre de motivation
    \item \cle{qualifications} : diplômes, certifications
    \item \cle{experience} : expérience professionnelle / stages
    \item \cle{strengths} / \cle{weaknesses} : points forts / points faibles
    \item \cle{salary} : salaire
    \item \cle{position} / \cle{vacancy} : poste / poste vacant
    \item \cle{reference} : référence (personne qui recommande)
\end{itemize}
\end{aretenir}

\courssub{Steps of the recruitment process}

\begin{popfigure}
\begin{tikzpicture}[
    node distance=1.2cm and 2.5cm,
    box/.style={rectangle, draw=popblue, thick, fill=popblue!10, rounded corners, minimum width=2.8cm, minimum height=1.1cm, align=center, font=\small},
    arrow/.style={-Stealth, thick, popblue}
]
\node[box, align=center] (app){Application\\ \emph{(Dépôt de candidature)}};
\node[box, right=of app, align=center] (short){Shortlisting\\ \emph{(Presélection)}};
\node[box, right=of short, align=center] (inter){Interview\\ \emph{(Entretien)}};
\node[box, right=of inter, align=center] (dec){Decision\\ \emph{(Décision)}};
\node[box, right=of dec, align=center] (feed){Feedback\\ \emph{(Retour)}};
\draw[arrow] (app) -- (short);
\draw[arrow] (short) -- (inter);
\draw[arrow] (inter) -- (dec);
\draw[arrow] (dec) -- (feed);
\end{tikzpicture}
\legende{Étapes du processus de recrutement : Application $\rightarrow$ Shortlisting $\rightarrow$ Interview $\rightarrow$ Decision $\rightarrow$ Feedback.}
\end{popfigure}

\begin{exemplebox}[Détail des étapes]
\begin{enumerate}
    \item \textbf{Application} : envoyer CV + cover letter pour une \cle{vacancy} (poste vacant).
    \item \textbf{Shortlisting} : l'employeur sélectionne les meilleurs dossiers.
    \item \textbf{Interview} : entretien en face-à-face, par téléphone ou visioconférence.
    \item \textbf{Decision} : l'employeur choisit le candidat retenu.
    \item \textbf{Feedback} : retour (positif ou négatif) envoyé aux candidats.
\end{enumerate}
\end{exemplebox}

\courssub{Giving instructions for a job interview}

\begin{propriete}[Expressions pour donner des instructions]
On utilise l'impératif (base verbale sans sujet) pour donner des consignes claires.
\end{propriete}

\begin{center}
\begin{tabularx}{\linewidth}{|l|X|}\hline
\textbf{Instruction (English)} & \textbf{Traduction / Contexte} \\ \hline
Dress formally. & S'habiller de façon formelle (costume, tenue professionnelle). \\ \hline
Arrive 15 minutes early. & Arriver 15 minutes en avance. \\ \hline
Bring your certificates. & Apporter vos certificats / diplômes. \\ \hline
Greet the panel politely. & Saluer le jury poliment. \\ \hline
Maintain eye contact. & Maintenir le contact visuel. \\ \hline
Speak clearly and confidently. & Parler clairement et avec assurance. \\ \hline
Listen carefully to each question. & Écouter attentivement chaque question. \\ \hline
Turn off your phone. & Éteindre votre téléphone. \\ \hline
\end{tabularx}
\end{center}

\courssub{Receiving and confirming instructions}

\begin{propriete}[Expressions pour recevoir et confirmer]
Ces formules montrent que vous écoutez activement et que vous avez compris.
\end{propriete}

\begin{center}
\begin{tabularx}{\linewidth}{|l|X|}\hline
\textbf{Expression} & \textbf{Utilisation} \\ \hline
Could you repeat that, please? & Demander une répétition poliment. \\ \hline
Do you mean I should \dots? & Reformuler pour vérifier sa compréhension. \\ \hline
I see, thank you. & Confirmer qu'on a compris. \\ \hline
Just to confirm: I need to \dots & Résumer l'instruction pour valider. \\ \hline
Should I bring anything else? & Demander s'il y a d'autres documents à apporter. \\ \hline
\end{tabularx}
\end{center}

\begin{exoresolu}[Exercice : Compléter le dialogue]
\textbf{Énoncé :} Complétez avec les expressions ci-dessus.
\begin{center}
\begin{tabularx}{\linewidth}{|l|X|}\hline
\textbf{Interviewer} & \textbf{Candidate} \\ \hline
``Please arrive at 8:30 a.m. and bring your BT diploma.'' & \textit{Do you mean I should arrive at 8:30 a.m. with my BT diploma?} \\ \hline
``Yes, and also your internship report.'' & \textit{Just to confirm: I need to bring my BT diploma and my internship report.} \\ \hline
``The interview will last about 30 minutes.'' & \textit{I see, thank you.} \\ \hline
\end{tabularx}
\end{center}
\tcblower
\textbf{Solution :} Voir le tableau ci-dessus. Les réponses attendues sont en italique.
\end{exoresolu}

\courssub{Common interview questions \& model answers}

\begin{aretenir}[Questions fréquentes en entretien]
\begin{enumerate}
    \item \textbf{Tell us about yourself.} (Parlez-nous de vous.)
    \item \textbf{Why should we hire you?} (Pourquoi devrions-nous vous embaucher ?)
    \item \textbf{What are your strengths and weaknesses?} (Quels sont vos points forts et vos points faibles ?)
    \item \textbf{Where do you see yourself in 5 years?} (Où vous voyez-vous dans 5 ans ?)
    \item \textbf{Why do you want to work here?} (Pourquoi voulez-vous travailler ici ?)
    \item \textbf{Do you have any questions for us?} (Avez-vous des questions pour nous ?)
\end{enumerate}
\end{aretenir}

\begin{methode}[Structure de réponse : « Tell us about yourself »]
Utilisez la méthode \textbf{PPF} (Present -- Past -- Future) :
\begin{enumerate}
    \item \textbf{Present} : qui vous êtes maintenant (classe, spécialité).
    \item \textbf{Past} : votre formation, vos stages, vos projets techniques.
    \item \textbf{Future} : pourquoi ce poste / stage vous intéresse.
\end{enumerate}
\end{methode}

\begin{exemplebox}[Réponse modèle pour Ariane (Génie Civil)]
\textbf{Interviewer :} « Tell us about yourself. »

\textbf{Ariane :} « Good morning. My name is Ariane Mbarga. I am a final-year student in \textbf{Civil Engineering} (Génie Civil) at the \textbf{Lycée Technique de Douala}. I am preparing my \textbf{Baccalauréat Technique} (Bac Tech). 

During my training, I completed a \textbf{three-month internship} at \textbf{SOGEA SATOM} where I assisted the site engineer in \textbf{reading plans}, \textbf{supervising concrete pouring} and \textbf{checking material quality}. I also worked on a school project designing a \textbf{drainage system} for a residential area.

I am \textbf{rigorous}, \textbf{punctual} and I work well in a \textbf{team}. I want to do this internship because SOGEA SATOM is a leader in public works in Cameroon and I want to learn from experienced engineers on real construction sites. »
\end{exemplebox}

\begin{exemplebox}[Réponse modèle : « What are your strengths and weaknesses? »]
\textbf{Strengths (Points forts) :}
\begin{itemize}
    \item « I am \textbf{organized} and \textbf{meet deadlines}. » (Je suis organisé et je respecte les délais.)
    \item « I have good \textbf{technical drawing} skills (AutoCAD). » (J'ai de bonnes compétences en dessin technique.)
    \item « I adapt quickly to new \textbf{tools and methods}. » (Je m'adapte vite aux nouveaux outils/méthodes.)
\end{itemize}

\textbf{Weaknesses (Points faibles) -- tournez-les en positif :}
\begin{itemize}
    \item « I sometimes \textbf{focus too much on details}, but I am learning to \textbf{prioritize}. » (Je me concentre trop sur les détails, mais j'apprends à prioriser.)
    \item « My \textbf{spoken English} needs improvement, so I \textbf{practise daily} with podcasts and colleagues. » (Mon anglais oral a besoin d'amélioration, donc je m'entraîne chaque jour.)
\end{itemize}
\end{exemplebox}

\courssub{Guided simulation : Interview at CAMTEL}

\begin{exoresolu}[Dialogue modèle annoté : Entretien chez CAMTEL pour un poste de technicien]
\textbf{Contexte :} Ariane postule pour un stage de technicien maintenance chez CAMTEL Douala. L'entretien se déroule en anglais (entreprise bilingue).

\begin{center}
\begin{tabularx}{\linewidth}{|c|l|X|}\hline
\textbf{\#} & \textbf{Interviewer} & \textbf{Candidate (Ariane)} \\ \hline
1 & ``Good morning. Please have a seat. Can you introduce yourself?'' & ``Good morning. Thank you. I am Ariane Mbarga, a Terminale student in \textbf{Electrotechnics} at Lycée Technique de Douala. I hold a \textbf{Probatoire} in Electrical Engineering and I am preparing my \textbf{Baccalauréat Technique}. I did a \textbf{two-month internship} at \textbf{ENEO} where I assisted in \textbf{transformer maintenance} and \textbf{cable fault detection}.'' \\ \hline
2 & ``Why CAMTEL?'' & ``CAMTEL is the \textbf{leading telecom operator} in Cameroon. I want to apply my \textbf{fiber optics} and \textbf{network configuration} skills learned in school on real infrastructure. I also want to improve my \textbf{technical English} because your teams work with international partners.'' \\ \hline
3 & ``What are your main technical skills?'' & ``I can \textbf{read electrical schematics}, use a \textbf{multimeter} and \textbf{OTDR} for fiber testing, and configure \textbf{routers/switches} (Cisco basics). I also know \textbf{safety procedures} (lockout/tagout).'' \\ \hline
4 & ``Describe a difficult situation you solved during your internship.'' & ``At ENEO, a \textbf{low-voltage cable fault} caused a power outage. I helped the team \textbf{locate the fault} using a \textbf{fault locator}, we \textbf{repaired the joint} and \textbf{restored power} in 3 hours. I learned to \textbf{stay calm under pressure}.'' \\ \hline
5 & ``What is your weakness?'' & ``I sometimes \textbf{hesitate to ask questions} when I'm not sure. Now I \textbf{prepare my questions} before each task and I ask my supervisor to \textbf{confirm critical steps}.'' \\ \hline
6 & ``Do you have any questions?'' & ``Yes. What \textbf{training opportunities} does CAMTEL offer to interns? And will I be assigned a \textbf{mentor} for the duration of the internship?'' \\ \hline
\end{tabularx}
\end{center}

\textbf{Annotations clés :}
\begin{itemize}
    \item Réponses \textbf{concrètes} (noms d'entreprises, outils, durées).
    \item Vocabulaire \textbf{technique} : \cle{transformer maintenance}, \cle{cable fault detection}, \cle{fiber optics}, \cle{OTDR}, \cle{lockout/tagout}.
    \item Structure \textbf{STAR} pour la question 4 : Situation -- Task -- Action -- Result.
    \item Question finale montre \textbf{intérêt} et \textbf{professionnalisme}.
\end{itemize}
\end{exoresolu}

\courssub{Body language \& attitudes during an interview}

\begin{popfigure}
\begin{tikzpicture}[
    node distance=0.8cm and 1.5cm,
    do/.style={rectangle, draw=popgreen, thick, fill=popgreen!10, rounded corners, minimum width=3.2cm, minimum height=0.9cm, align=center, font=\small},
    dont/.style={rectangle, draw=poporange, thick, fill=poporange!10, rounded corners, minimum width=3.2cm, minimum height=0.9cm, align=center, font=\small},
    title/.style={font=\bfseries\large, text=popdark}
]
\node[title] at (0,4.5) {Do's \checkmark};
\node[title] at (7,4.5) {Don'ts \times};
\node[do] (d1) at (0,3.5) {Arrive 10--15 min early};
\node[do] (d2) at (0,2.5) {Firm handshake + smile};
\node[do] (d3) at (0,1.5) {Sit up straight, lean slightly forward};
\node[do] (d4) at (0,0.5) {Maintain natural eye contact};
\node[do] (d5) at (0,-0.5) {Listen fully before answering};
\node[do] (d6) at (0,-1.5) {Ask relevant questions at the end};
\node[do] (d7) at (0,-2.5) {Thank the panel + send follow-up email};
\node[dont] (n1) at (7,3.5) {Arrive late or out of breath};
\node[dont] (n2) at (7,2.5) {Weak handshake, no eye contact};
\node[dont] (n3) at (7,1.5) {Slouch, cross arms, fidget};
\node[dont] (n4) at (7,0.5) {Stare at the floor or ceiling};
\node[dont] (n5) at (7,-0.5) {Interrupt the interviewer};
\node[dont] (n6) at (7,-1.5) {Say ``I don't know'' to everything};
\node[dont] (n7) at (7,-2.5) {Leave without thanking anyone};
\end{tikzpicture}
\legende{Do's vs Don'ts pendant un entretien d'embauche : attitudes et langage corporel.}
\end{popfigure}

\begin{attention}[Erreur fréquente : « I have 18 years »]
\textbf{Ne dites jamais :} « I have 18 years. » \textbf{(Traduction littérale incorrecte.)}

\textbf{Dites :} « \textbf{I am 18 years old.} » ou simplement « \textbf{I am 18.} »

\begin{center}
\begin{tabular}{|l|l|}\hline
\textbf{Incorrect (Franglais)} & \textbf{Correct (English)} \\ \hline
I have 20 years. & I am 20 (years old). \\ \hline
I have hunger. & I am hungry. \\ \hline
I have thirst. & I am thirsty. \\ \hline
I have cold. & I am cold. \\ \hline
I have right. & I am right. \\ \hline
\end{tabular}
\end{center}
\textbf{Règle :} En anglais, l'âge et les états physiques s'expriment avec \textbf{BE + adjectif}, pas HAVE + nom.
\end{attention}

\begin{savaistu}[Le savais-tu ?]
Au Cameroun, de nombreuses entreprises bilingues (\textbf{SONARA}, \textbf{Port Autonome de Douala}, \textbf{CAMTEL}, \textbf{ENEO}, \textbf{Brasseries du Cameroun}) exigent un entretien \textbf{partiellement en anglais} même pour des postes purement techniques. Pourquoi ? Parce que ces entreprises travaillent avec des partenaires nigérians, ghanéens, sud-africains ou chinois, et la documentation technique (normes IEC, manuels CAT, schémas Siemens) est souvent en anglais. Maîtriser l'anglais technique est un \textbf{atout majeur} pour l'employabilité !
\end{savaistu}

\courssub{Practice exercises}

\begin{exercice}[Exercice 1 : Vocabulary matching]
Faites correspondre chaque terme à sa définition.
\begin{enumerate}
    \item Applicant \hfill a) Personne qui mène l'entretien
    \item Interviewer \hfill b) Lettre de motivation
    \item Cover letter \hfill c) Points forts
    \item Strengths \hfill d) Candidat
    \item Vacancy \hfill e) Poste vacant
\end{enumerate}
\end{exercice}


\begin{corrige}[Exercice 1]
1--d, 2--a, 3--b, 4--c, 5--e.
\end{corrige}

\begin{exercice}[Exercice 2 : Transformez en instructions (impératif)]
\begin{enumerate}
    \item You should dress formally. $\rightarrow$ \_\_\_\_\_\_\_\_\_\_\_\_\_\_\_\_\_\_\_\_\_\_
    \item You must arrive early. $\rightarrow$ \_\_\_\_\_\_\_\_\_\_\_\_\_\_\_\_\_\_\_\_\_\_
    \item You have to bring your CV. $\rightarrow$ \_\_\_\_\_\_\_\_\_\_\_\_\_\_\_\_\_\_\_\_\_\_
    \item You need to switch off your phone. $\rightarrow$ \_\_\_\_\_\_\_\_\_\_\_\_\_\_\_\_\_\_\_\_\_\_
\end{enumerate}
\end{exercice}


\begin{corrige}[Exercice 2]
1. Dress formally. 2. Arrive early. 3. Bring your CV. 4. Switch off your phone.
\end{corrige}

\begin{exercice}[Exercice 3 : Réponse « Tell us about yourself »]
Rédigez votre propre réponse (5--6 phrases) en utilisant la méthode PPF (Present -- Past -- Future). Vous êtes en Terminale \textbf{[votre spécialité]} au \textbf{[votre lycée]}. Mentionnez : votre diplôme visé (Bac Tech / BT), un stage ou projet technique, une qualité, votre motivation pour le poste visé.
\end{exercice}


\begin{corrige}[Exercice 3]
\textbf{Modèle :} « Good morning. My name is [Prénom Nom]. I am a final-year student in [Spécialité] at [Lycée Technique]. I am preparing my [Baccalauréat Technique / BT]. During my training, I completed a [durée] internship at [Entreprise] where I [tâche technique 1] and [tâche technique 2]. I am [qualité 1] and [qualité 2]. I want this position because [raison liée à l'entreprise / au poste]. »
\end{corrige}

\begin{exercice}[Exercice 4 : Simulation orale (par binômes)]
\textbf{Consigne :} L'un joue le recruteur (ENEO / CAMTEL / Atelier de Douala), l'autre le candidat. Utilisez les 6 questions de la section « Common interview questions ». Le recruteur note : vocabulaire technique, structure PPF/STAR, langage corporel (contact visuel, posture), politesse. Inversez les rôles.

\textbf{Grille d'évaluation (sur 10) :}
\begin{itemize}
    \item Vocabulaire technique pertinent : /3
    \item Structure des réponses (PPF / STAR) : /3
    \item Langage corporel \& politesse : /2
    \item Fluidité \& prononciation : /2
\end{itemize}
\end{exercice}

\begin{aretenir}[Check-list avant l'entretien : \textbf{READY}]
\begin{center}
\begin{tabular}{|l|l|}\hline
\textbf{Lettre} & \textbf{Action} \\ \hline
\textbf{R} & \textbf{Research} the company (étudier l'entreprise) \\ \hline
\textbf{E} & \textbf{Prepare} your answers (préparer les réponses types) \\ \hline
\textbf{A} & \textbf{Arrive} early (arriver en avance) \\ \hline
\textbf{D} & \textbf{Dress} professionally (s'habiller professionnellement) \\ \hline
\textbf{Y} & \textbf{Your} best attitude: polite, confident, grateful (votre meilleure attitude : poli, confiant, reconnaissant) \\ \hline
\end{tabular}
\end{center}
\end{aretenir}

\coursec{Recreational activities (Listening \& Vocabulary)}

After practising how to \cle{give} and \cle{receive instructions} for a job interview, we now turn to a more relaxed but equally important part of life: \cle{recreation}. In the world of work, knowing how to manage your \cle{free time} and talk about your \cle{hobbies} helps you socialise with colleagues, reduce stress, and maintain a healthy work-life balance. This section will equip you with the vocabulary to describe your favourite pastimes, the grammar to express your preferences correctly, and the listening strategies to understand native and non-native speakers talking about leisure.

\courssub{What is Recreation?}

\begin{definition}[Recreation / Leisure]
\cle{Recreation} (or \cle{leisure}) refers to activities that people do during their \cle{free time} (time when they are not working or studying) for \cle{enjoyment}, \cle{relaxation}, or \cle{pleasure}. These activities are voluntary and provide a break from daily duties.
\end{definition}

\begin{aretenir}[Key Terms]
\begin{itemize}
    \item \textbf{Hobby} (\emph{passe-temps favori}) : An activity done regularly in one's leisure time for pleasure.
    \item \textbf{Pastime} (\emph{divertissement}) : Something that serves to make time pass agreeably.
    \item \textbf{Indoor activities} (\emph{activités d'intérieur}) : Done inside the house or a building (e.g., reading, playing Ludo, sewing, watching TV, video games).
    \item \textbf{Outdoor activities} (\emph{activités de plein air}) : Done outside in the open air (e.g., football, hiking, fishing, dancing, swimming, jogging).
\end{itemize}
\end{aretenir}

\courssub{Vocabulary: Classifying Recreational Activities}

Let us classify common recreational activities in Cameroon. Understanding this distinction helps you use the right prepositions (\emph{play \textbf{at} home} vs \emph{play \textbf{in} the field}) and vocabulary.

\begin{popfigure}
\begin{center}
\begin{tikzpicture}[
    node distance=1.2cm and 2cm,
    main/.style={draw=popdark, thick, rounded corners, fill=popcream, text width=3.5cm, align=center, font=\bfseries},
    cat/.style={draw=popblue, thick, rounded corners, fill=popblueL, text width=4cm, align=center},
    ex/.style={draw=popgreen, fill=popgreenL, rounded corners, text width=5cm, align=left, font=\small}
]
% Title
\node[main, minimum height=1.2cm] (title) {RECREATIONAL ACTIVITIES};
% Two branches
\node[cat, below left=of title, xshift=-1cm, align=center] (indoor){INDOOR ACTIVITIES\\ \textit{(Activités d'intérieur)}};
\node[cat, below right=of title, xshift=1cm, align=center] (outdoor){OUTDOOR ACTIVITIES\\ \textit{(Activités de plein air)}};
% Indoor examples
\node[ex, below=of indoor, anchor=north, align=center] (in1){\textbf{Intellectual / Relaxing:}\\
Reading novels, newspapers\\
Watching TV / Movies\\
Listening to music (Makossa, Coupé-Décalé, Gospel)\\
Playing Ludo, Cards, Chess, Scrabble};
\node[ex, below=of in1, anchor=north, align=center] (in2){\textbf{Creative / Manual:}\\
Sewing / Knitting / Tailoring\\
Drawing / Painting\\
Cooking / Baking new recipes\\
Video games (PlayStation, Mobile games)};
% Outdoor examples
\node[ex, below=of outdoor, anchor=north, align=center] (out1){\textbf{Sports / Physical:}\\
Football (quarter/field)\\
Jogging / Running\\
Swimming (pool, river, beach)\\
Hiking / Walking (Mount Fako, hills)\\
Cycling};
\node[ex, below=of out1, anchor=north, align=center] (out2){\textbf{Social / Cultural:}\\
Traditional dance (Bikutsi, Assiko)\\
Modern dance / Nightclub\\
Fishing (river, sea)\\
Picnics / Barbecues (\textit{``Buy-cut''})\\
Visiting friends / Family gatherings};
% Arrows
\draw[popdark, thick, ->] (title.south west) -- (indoor.north east);
\draw[popdark, thick, ->] (title.south east) -- (outdoor.north west);
\end{tikzpicture}
\end{center}
\legende{Classification of recreational activities common among Cameroonian youth.}
\end{popfigure}

\courssub{Expressing Likes, Dislikes and Preferences}

To talk about recreation, you must master the language of \cle{preference}. In English, verbs of feeling (\emph{like, love, enjoy, hate, prefer}) are followed by specific structures.

\begin{propriete}[Structures of Preference]
\begin{center}
\begin{tabularx}{\linewidth}{|l|X|X|}\hline
\textbf{Verb / Expression} & \textbf{Structure} & \textbf{Example} \\ \hline
\cle{Enjoy} & \emph{enjoy + V-ing} & I \textbf{enjoy reading} novels. \\ \hline
\cle{Be fond of} & \emph{be fond of + V-ing / Noun} & She \textbf{is fond of dancing}. \\ \hline
\cle{Be keen on} & \emph{be keen on + V-ing / Noun} & We \textbf{are keen on football}. \\ \hline
\cle{Like / Love / Hate} & \emph{like/love/hate + V-ing} (general preference) & I \textbf{like swimming}. \\
& \emph{like/love/hate + to-infinitive} (specific occasion/habit) & I \textbf{like to swim} every morning. \\ \hline
\cle{Can't stand} & \emph{can't stand + V-ing} & He \textbf{can't stand watching} violent movies. \\ \hline
\cle{Prefer} & \emph{prefer + V-ing + to + V-ing} & I \textbf{prefer reading to watching} TV. \\
& \emph{would rather (\emph{'d rather}) + bare infinitive + than + bare infinitive} & I \textbf{'d rather play} Ludo \textbf{than watch} TV. \\ \hline
\end{tabularx}
\end{center}
\end{propriete}

\begin{exemplebox}[Expressing Preferences in Context]
\textbf{A:} \textit{What do you usually do on Sundays?}\\
\textbf{B:} \textit{I \textbf{enjoy going} to church in the morning. In the afternoon, I \textbf{'d rather rest} at home. I \textbf{can't stand staying} idle for too long, so I often \textbf{go for a walk}.}\\
\textbf{A:} \textit{Do you like video games?}\\
\textbf{B:} \textit{Not really. I \textbf{prefer playing} football \textbf{to playing} video games. I \textbf{am keen on} outdoor sports.}
\end{exemplebox}

\begin{attention}[Common Error: \texorpdfstring{\textquotesingle I am agree\textquotesingle}{I am agree}]
\textbf{Never say:} \textit{``I am agree with you.''} \\
\textbf{Correct:} \textit{``I \textbf{agree} with you.''} \\
\textbf{Reason:} \cle{Agree} is a main verb, not an adjective. It does not take the auxiliary \emph{be}.
\begin{itemize}
    \item \textit{I agree.} (Present Simple)
    \item \textit{I don't agree.} (Negative)
    \item \textit{Do you agree?} (Question)
\end{itemize}
Similarly: \textit{``I am interested in...''} (Interested is an adjective $\rightarrow$ \textbf{be} + interested) vs \textit{``I interest...''} (Wrong).
\end{attention}

\courssub{Talking about Frequency}

When describing habits and routines related to hobbies, we use \cle{adverbs of frequency} and \cle{time expressions}.

\begin{propriete}[Adverbs of Frequency (Position)]
\textbf{Rule:} Place the adverb \textbf{before the main verb} but \textbf{after the verb \emph{be}} and \textbf{after auxiliaries}.
\[
\text{Subject} + \text{Adverb} + \text{Main Verb} \quad | \quad \text{Subject} + \text{be/auxiliary} + \text{Adverb} + \text{Main Verb}
\]
\begin{center}
\begin{tabular}{|c|l|l|}\hline
\textbf{Frequency} & \textbf{Adverb} & \textbf{Example} \\ \hline
100\% & \cle{Always} & I \textbf{always} play football on Saturdays. \\ \hline
90\% & \cle{Usually} / \cle{Generally} & She \textbf{usually} reads before sleeping. \\ \hline
70\% & \cle{Often} / \cle{Frequently} & We \textbf{often} go hiking. \\ \hline
50\% & \cle{Sometimes} & He \textbf{sometimes} watches TV. \\ \hline
30\% & \cle{Rarely} / \cle{Seldom} & They \textbf{rarely} go to the cinema. \\ \hline
10\% & \cle{Hardly ever} & I \textbf{hardly ever} sew. \\ \hline
0\% & \cle{Never} & She \textbf{never} fishes. \\ \hline
\end{tabular}
\end{center}
\textbf{With \emph{be}:} \textit{He \textbf{is always} late for training.} \\
\textbf{With auxiliaries:} \textit{We \textbf{have never} been fishing.}
\end{propriete}

\begin{aretenir}[Expressions of Frequency (usually at end of sentence)]
\begin{itemize}
    \item \cle{Once a week} / \cle{Twice a month} / \cle{Three times a year}
    \item \cle{Every day} / \cle{Every weekend} / \cle{Every Sunday}
    \item \cle{From time to time} / \cle{Once in a while}
\end{itemize}
\end{aretenir}

\courssub{Listening Skills: Strategies for Success}

Listening in an exam context (Probatoire/Baccalauréat technique) requires a method. Do not just ``hear''; \cle{listen actively}.

\begin{methode}[How to Succeed in a Listening Task]
\begin{enumerate}
    \item \textbf{Read the questions FIRST} (use the 1-2 minutes preparation time). Underline \cle{keywords} (Who? What? Where? When? How often?).
    \item \textbf{Predict the topic and vocabulary.} If keywords are ``football'', ``weekend'', ``friends'', expect sports vocabulary.
    \item \textbf{First listening (Global / Gist):} Get the \cle{main idea}. Who is speaking? What is the general topic? Don't panic if you miss details.
    \item \textbf{Second listening (Selective / Detail):} Focus \textbf{only} on the answers to your questions. Write short notes/abbreviations (e.g., ``wknd'' for weekend, ``w/'' for with, ``Fball'' for football).
    \item \textbf{Check and Complete:} Use the pause or third listening to verify spelling and grammar in your answers.
\end{enumerate}
\end{methode}

\courssub{Listening Practice: Weekend Leisure in Yaoundé}

\textbf{Context:} Two Terminale students, \textbf{Chantal} and \textbf{Bruno}, meet at the \emph{Carrefour Warda} in Yaoundé on a Saturday morning. They discuss their weekend plans.

\begin{exoresolu}[Listening Comprehension: Grid Completion]
\textbf{Instructions:} Listen to the dialogue (or read the script below) and complete the grid.

\textbf{Audio Script:}
\textbf{Chantal:} Hi Bruno! Nice to see you. What are you up to this weekend?
\textbf{Bruno:} Hi Chantal! I'm doing great. Well, today I have football training at the \textbf{Stade Omnisports} at \textbf{3 p.m.} I play for the \textbf{quarter team}, ``Les Lions Juniors''. We train \textbf{twice a week}, Wednesdays and Saturdays.
\textbf{Chantal:} That sounds energetic! I prefer something quieter. This afternoon, I'm \textbf{keen on} finishing the novel I started: ``The Poor Christ of Bomba''. I \textbf{enjoy reading} African literature. Tomorrow, Sunday, I \textbf{'d rather go} to the \textbf{Ngoa-Ekellé market} with my sister to buy fabric for \textbf{sewing} a new dress.
\textbf{Bruno:} Sewing? That's a cool hobby. Do you do it often?
\textbf{Chantal:} \textbf{Once a week}, usually on Sundays. I \textbf{can't stand} buying ready-made clothes; they never fit well. What about tomorrow?
\textbf{Bruno:} Tomorrow morning, we have a \textbf{friendly match} against a team from \textbf{Essos} at \textbf{9 a.m.} In the afternoon, I'll just rest. I \textbf{never} go out on Sunday evenings; I prepare my lessons for Monday.
\textbf{Chantal:} You are well organised! Enjoy your match. See you at school on Monday!
\textbf{Bruno:} Thanks! You too. Bye!

\textbf{Task: Complete the table.}
\tcblower
\textbf{Solution: Completed Comprehension Grid}

\begin{center}
\begin{tabularx}{\linewidth}{|l|X|X|}\hline
\textbf{Information Required} & \textbf{Chantal} & \textbf{Bruno} \\ \hline
\textbf{Main Activity (Saturday)} & Reading a novel (African literature) & Football training (Quarter team) \\ \hline
\textbf{Location} & Home (Sat) / Ngoa-Ekellé Market (Sun) & Stade Omnisports (Sat), Essos field (Sun) \\ \hline
\textbf{Time} & Afternoon (Sat), Morning (Sun) & 3 p.m. (Sat), 9 a.m. (Sun) \\ \hline
\textbf{Frequency} & Sewing: Once a week & Training: Twice a week \\ \hline
\textbf{Preference Expression Used} & ``I enjoy reading'', ``I'm keen on'', ``I'd rather go'', ``I can't stand'' & ``I never go out'' \\ \hline
\textbf{Companion} & Sister (Sunday) & Teammates / Friends \\ \hline
\end{tabularx}
\end{center}
\end{exoresolu}

\courssub{Speaking Practice: Leisure Activities of Cameroonian Youth}

Now it is your turn to produce language. Use the vocabulary, frequency adverbs, and preference structures learned.

\begin{exercice}[Guided Oral Production]
\textbf{Work in pairs.} Ask and answer questions about your recreational activities. Use the prompts below. Record yourselves if possible.
\begin{enumerate}
    \item \textbf{Indoor:} What indoor activity do you \textbf{enjoy} most? (Reading, Ludo, Video games, Sewing, Movies?) How \textbf{often}?
    \item \textbf{Outdoor:} Are you \textbf{keen on} football, hiking, dancing, or swimming? Where do you go? With whom?
    \item \textbf{Cultural:} Do you listen to \textbf{Makossa}, \textbf{Bikutsi}, \textbf{Coupé-Décalé}, or \textbf{Gospel}? Can you dance \textbf{Assiko} or \textbf{Bottle dance}?
    \item \textbf{Preference:} \textit{``Would you rather play video games \textbf{or} play football with friends? Why?''}
    \item \textbf{Frequency:} \textit{``How often do you go to the \textbf{quarter field} / \textbf{youth centre} / \textbf{library}?''}
\end{enumerate}
\textbf{Model Answer:}
\textit{``I \textbf{enjoy playing} Ludo with my cousins \textbf{twice a week}. I \textbf{am fond of} Makossa music. \textbf{At the weekend}, I \textbf{usually} go to the quarter field to watch football matches. I \textbf{'d rather watch} than play because I \textbf{can't stand} running too much!''}
\end{exercice}

\courssub{Statistics: Favourite Leisure Activities (Data Interpretation)}

Data literacy is part of technical education. Let us analyse a survey conducted on 40 Terminale students in a technical high school in Douala.

\begin{popfigure}
\begin{center}
\begin{tikzpicture}
\begin{axis}[
    ybar,
    bar width=18pt,
    width=\linewidth,
    height=10cm,
    enlarge x limits=0.25,
    ylabel={Number of Students},
    xlabel={Leisure Activity},
    symbolic x coords={Football, {Music/Dance}, Reading, {Video Games}, Sewing, Others},
    xtick=data,
    nodes near coords,
    nodes near coords align={vertical},
    ymin=0, ymax=16,
    ytick={0,2,4,6,8,10,12,14,16},
    tick label style={font=\small},
    label style={font=\small},
    title style={font=\bfseries, align=center},
    title={Favourite Leisure Activities of 40 Terminale Students},
    grid=major,
    grid style={popline},
    every axis plot/.append style={fill=popblue!70, draw=popdark},
    every node near coord/.append style={font=\bfseries\small, popdark}
]
\addplot coordinates {
    (Football,14)
    ({Music/Dance},10)
    (Reading,6)
    ({Video Games},5)
    (Sewing,3)
    (Others,2)
};
\end{axis}
\end{tikzpicture}
\end{center}
\legende{Bar chart showing the distribution of favourite leisure activities (Sample: 40 students).}
\end{popfigure}

\begin{exoresolu}[Data Analysis \& Language Practice]
\textbf{Questions:}
\begin{enumerate}
    \item Which activity is the \textbf{most popular}? Which is the \textbf{least popular} among the listed ones?
    \item What \textbf{percentage} of students prefer Football? (Round to nearest whole number).
    \item Compare \textbf{Reading} and \textbf{Video Games} using a comparative structure.
    \item Write a sentence using \emph{``The majority of...''} and \emph{``A minority of...''}.
\end{enumerate}
\tcblower
\textbf{Solutions:}
\begin{enumerate}
    \item \textbf{Most popular:} Football (14 students). \textbf{Least popular:} Sewing (3 students) -- excluding ``Others''.
    \item Percentage for Football: $\frac{14}{40} \times 100 = 35\%$.
    \item \textit{``More students prefer reading (6) than video games (5).''} OR \textit{``Reading is slightly more popular than video games.''}
    \item \textit{``The \textbf{majority of} students (35\%) prefer football, while a \textbf{minority} prefer sewing (7.5\%).''}
\end{enumerate}
\end{exoresolu}

\courssub{Exercises: Consolidation}

\begin{exercice}[Vocabulary \& Grammar: Multiple Choice]
\textbf{Choose the correct option (A, B, or C).}
\begin{enumerate}
    \item My brother \_\_\_\_\_\_ playing the guitar every evening.
    \begin{itemize}
        \item[A.] enjoy \item[B.] enjoys \item[C.] enjoying
    \end{itemize}
    \item We \_\_\_\_\_\_ go to the cinema; maybe once a year.
    \begin{itemize}
        \item[A.] often \item[B.] usually \item[C.] rarely
    \end{itemize}
    \item I \_\_\_\_\_\_ swimming \_\_\_\_\_\_ running. I prefer the pool.
    \begin{itemize}
        \item[A.] prefer / than \item[B.] prefer / to \item[C.] 'd rather / than
    \end{itemize}
    \item She \_\_\_\_\_\_ stand the noise of the generator.
    \begin{itemize}
        \item[A.] can \item[B.] can't \item[C.] doesn't
    \end{itemize}
    \item \_\_\_\_\_\_ you \_\_\_\_\_\_ interested in learning Spanish for your job?
    \begin{itemize}
        \item[A.] Do / be \item[B.] Are / -- \item[C.] Have / been
    \end{itemize}
\end{enumerate}
\end{exercice}

\begin{corrige}[Exercice: Vocabulary \& Grammar]
\begin{enumerate}
    \item \textbf{B (enjoys)} -- 3rd person singular requires \emph{-s}.
    \item \textbf{C (rarely)} -- ``Once a year'' indicates very low frequency.
    \item \textbf{B (prefer / to)} -- Structure: \emph{prefer + V-ing + to + V-ing}. (Option C would require bare infinitive: \emph{'d rather swim than run}).
    \item \textbf{B (can't)} -- \emph{Can't stand} = hate intensely.
    \item \textbf{B (Are / --)} -- \emph{Interested} is an adjective $\rightarrow$ \emph{be} + interested.
\end{enumerate}
\end{corrige}

\begin{exercice}[Writing: Short Paragraph]
\textbf{Write a short paragraph (60-80 words) about your ideal weekend.}
\textbf{Include:}
\begin{itemize}
    \item One \textbf{indoor} and one \textbf{outdoor} activity.
    \item Two \textbf{frequency expressions} (e.g., \emph{usually, twice a week}).
    \item Two \textbf{preference structures} (e.g., \emph{enjoy, 'd rather, can't stand, prefer... to}).
    \item Mention a specific \textbf{place in your town/region} (e.g., \emph{Mfoundi market, Limbe beach, Foumban palace}).
\end{itemize}
\end{exercice}

\begin{corrige}[Exercice: Writing - Sample Answer]
\textit{``My ideal weekend starts on Saturday morning. I \textbf{usually} help my mother at the \textbf{Mfoundi market} (outdoor). In the afternoon, I \textbf{enjoy reading} technical magazines at home (indoor). On Sunday, I \textbf{'d rather play} football with my friends at the \textbf{quarter field} \textbf{than stay} indoors. I \textbf{can't stand} being lazy. I go to the field \textbf{twice a month}. In the evening, I prepare my lessons for the week.''}
\end{corrige}

\begin{savaistu}[Le Football au Cameroun : Une Passion Nationale]
Le football n'est pas seulement un sport au Cameroun, c'est une \cle{institution sociale}.
\begin{itemize}
    \item Les \textbf{Lions Indomptables} ont remporté la \textbf{Coupe d'Afrique des Nations (CAN) 5 fois} : \textbf{1984, 1988, 2000, 2002, 2017}.
    \item Roger \textbf{Milla} a ébloui le monde en 1990 (Italie) en devenant le plus vieux buteur de l'histoire de la Coupe du Monde (42 ans en 1994).
    \item Chaque quartier (\emph{``quarter''}) a son terrain de terre battue où les jeunes s'entraînent \textbf{tous les soirs} après l'école ou le travail.
    \item Le vocabulaire technique s'infiltre dans le langage courant : \textit{``Faire un dribble''} (éviter un problème), \textit{``Marquer un but''} (réussir un coup), \textit{``Hors-jeu''} (être exclu/perdu).
\end{itemize}
Parler football est le meilleur \cle{brise-glace} (ice-breaker) pour s'intégrer dans n'importe quelle équipe de travail au Cameroun !
\end{savaistu}

\coursec{Tenses: present perfect continuous \& future continuous}

\courssub{Transition from recreational activities}
In the previous section, we listened to people talking about their hobbies and free time. You heard sentences like \emph{``I have been playing football since childhood''} or \emph{``This time next week, I will be hiking on Mount Cameroon.''} These sentences use two important tenses: the \textbf{present perfect continuous} and the \textbf{future continuous}. Before we study them in detail, let us quickly revise the tenses you learned in Première and identify common problem areas.

\begin{aretenir}[Quick revision of Première tenses \& problem areas]
\begin{itemize}
    \item \textbf{Present simple:} habits, facts, timetables. \emph{I work at the garage.} (J travaille au garage.)
    \item \textbf{Present continuous:} action happening now, temporary situations. \emph{I am repairing a car.} (Je suis en train de réparer une voiture.)
    \item \textbf{Past simple:} finished action at a definite past time. \emph{I repaired a car yesterday.} (J'ai réparé une voiture hier.)
    \item \textbf{Past continuous:} action in progress at a past moment. \emph{I was repairing the car when the boss arrived.} (Je réparais la voiture quand le chef est arrivé.)
    \item \textbf{Present perfect simple:} past action with present result / experience. \emph{I have repaired the car.} (J'ai réparé la voiture — elle est prête maintenant.)
    \item \textbf{Present perfect continuous:} \cle{duration} of an activity up to now (see below).
    \item \textbf{Future (will / going to):} predictions, plans, decisions. \emph{I will repair the car tomorrow.} / \emph{I am going to repair the car.}
    \item \textbf{Future continuous:} action in progress at a specific future moment (see below).
\end{itemize}
\textbf{Common problem areas to watch:}
\begin{itemize}
    \item Confusing \emph{present simple} and \emph{present continuous} (\emph{*I am knowing him} $\rightarrow$ \emph{I know him}).
    \item Using \emph{present simple} with \emph{since/for} (\emph{*I work here since 2020} $\rightarrow$ \emph{I have been working here since 2020}).
    \item Forgetting \emph{been} in perfect continuous (\emph{*I have working} $\rightarrow$ \emph{I have been working}).
    \item Using continuous forms with stative verbs (\emph{*I am having a car} $\rightarrow$ \emph{I have a car}).
\end{itemize}
\end{aretenir}

Now we focus on two tenses that highlight \textbf{duration} (present perfect continuous) and \textbf{actions in progress at a specific future moment} (future continuous).

\courssub{Present perfect continuous}

\begin{definition}[Present perfect continuous]
A tense used for actions that started in the past and are still continuing, or have just stopped, with emphasis on \textbf{duration} (how long).
\end{definition}

\begin{propriete}[Formation]
\[
\text{Subject} + \text{have/has} + \text{been} + \text{Verb-ing}
\]
\begin{center}
\begin{tabularx}{\linewidth}{|l|X|X|X|}\hline
\textbf{Person} & \textbf{Affirmative} & \textbf{Negative} & \textbf{Interrogative} \\ \hline
I / You / We / They & I \textbf{have been working} & I \textbf{have not been working} & \textbf{Have I been working?} \\
He / She / It & She \textbf{has been studying} & She \textbf{has not been studying} & \textbf{Has she been studying?} \\ \hline
\end{tabularx}
\end{center}
\textbf{Short answers:} \emph{Yes, I have. / No, I haven't.} \quad \emph{Yes, she has. / No, she hasn't.}
\end{propriete}

\begin{exemplebox}[Conjugation of ``to learn'']
\begin{center}
\begin{tabularx}{\linewidth}{|l|X|X|}\hline
\textbf{Pronoun} & \textbf{Affirmative} & \textbf{Negative} \\ \hline
I & have been learning & have not been learning \\
You & have been learning & have not been learning \\
He / She / It & has been learning & has not been learning \\
We & have been learning & have not been learning \\
They & have been learning & have not been learning \\ \hline
\end{tabularx}
\end{center}
\end{exemplebox}

\courssub{Uses of the present perfect continuous}

\begin{propriete}[Use 1: Action started in the past and still continuing]
We use it to say \textbf{how long} an action has been happening. The action is not finished.
\begin{itemize}
    \item \emph{I \textbf{have been learning} English \textbf{for} six years.} (J'apprends l'anglais depuis six ans — et je continue.)
    \item \emph{Mballa \textbf{has been working} at the garage \textbf{since} 2019.} (Mballa travaille au garage depuis 2019.)
    \item \emph{How long \textbf{have you been studying} technical drawing?} (Depuis combien de temps étudies-tu le dessin technique ?)
\end{itemize}
\end{propriete}

\begin{propriete}[Use 2: Recent action with visible result now]
The action stopped recently, but we see the evidence (result in the present).
\begin{itemize}
    \item \emph{She is tired because she \textbf{has been working} all day.} (Elle est fatiguée parce qu'elle a travaillé toute la journée — on voit le résultat : elle est fatiguée.)
    \item \emph{Your hands are dirty! Have you \textbf{been repairing} the engine?} (Tes mains sont sales ! Tu as réparé le moteur ?)
    \item \emph{The floor is wet. It \textbf{has been raining}.} (Le sol est mouillé. Il a plu.)
\end{itemize}
\end{propriete}

\begin{propriete}[Time markers: \cle{for}, \cle{since}, \cle{all day}, \cle{lately}, \cle{recently}, \cle{how long}]
\begin{center}
\begin{tabularx}{\linewidth}{|l|X|X|}\hline
\textbf{Marker} & \textbf{Meaning} & \textbf{Example} \\ \hline
\textbf{for} + duration & pendant (durée) & \emph{for three years}, \emph{for two hours}, \emph{for a long time} \\
\textbf{since} + starting point & depuis (point de départ) & \emph{since 2020}, \emph{since January}, \emph{since Monday}, \emph{since I left school} \\
\textbf{all day / all week / all morning} & toute la journée / la semaine / la matinée & \emph{I have been studying all day.} \\
\textbf{lately / recently} & ces derniers temps / récemment & \emph{Have you been sleeping well lately?} \\
\textbf{How long...?} & Depuis combien de temps ? & \emph{How long have you been waiting?} \\ \hline
\end{tabularx}
\end{center}
\end{propriete}

\begin{attention}[Erreur fréquente : présent simple ou present continuous au lieu de present perfect continuous]
\begin{itemize}
    \item \textbf{Faux :} \emph{``I \textbf{am learning} English \textbf{since} 2020.''} \quad (\emph{present continuous} + \emph{since})
    \item \textbf{Faux :} \emph{``I \textbf{learn} English \textbf{since} 2020.''} \quad (\emph{present simple} + \emph{since})
    \item \textbf{Vrai :} \emph{``I \textbf{have been learning} English \textbf{since} 2020.''}
    \item \textbf{Explication :} En français on dit « j'apprends depuis 2020 » (présent), mais en anglais, une action qui a commencé dans le passé et qui continue utilise le \textbf{present perfect continuous}, pas le present continuous ni le present simple.
\end{itemize}
\end{attention}

\begin{popfigure}
\begin{tikzpicture}[scale=0.9]
% Timeline for Present Perfect Continuous
\draw[popblue, thick, ->] (0,0) -- (13,0) node[right] {\textbf{Time}};
% Past
\draw[popblue, thick] (1,0) -- (1,-0.3) node[below] {\textbf{PAST}};
\node[popblue, anchor=north] at (1,-0.5) {Action starts};
% Continuous arrow
\draw[poporange, very thick, ->] (1,0.5) -- (10,0.5) node[midway, above] {\textbf{Action in progress}};
% Present
\draw[popgreen, thick] (10,0) -- (10,-0.3) node[below] {\textbf{NOW}};
\node[popgreen, anchor=north] at (10,-0.5) {Still continuing / Just stopped};
% Markers
\node[poppurple, anchor=south] at (3,1.2) {\textbf{for 3 years}};
\node[poppurple, anchor=south] at (6,1.2) {\textbf{since 2021}};
\node[poppurple, anchor=south] at (9,1.2) {\textbf{lately / recently}};
% Example box
\node[draw=popblue, fill=popblue!10, rounded corners, align=center, anchor=north] at (6.5,-2) {
    \textbf{Example:} \emph{``I have been studying}\\
    \emph{technical drawing for 4 years.''}\\
    (J'étudie le dessin technique depuis 4 ans.)
};
\end{tikzpicture}
\legende{Present perfect continuous timeline: action starts in the past and continues up to now (or just stopped). The orange arrow shows the ongoing duration.}
\end{popfigure}

\courssub{Future continuous}

\begin{definition}[Future continuous]
A tense used for actions that will be \textbf{in progress at a specific moment in the future}.
\end{definition}

\begin{propriete}[Formation]
\[
\text{Subject} + \text{will} + \text{be} + \text{Verb-ing}
\]
\begin{center}
\begin{tabularx}{\linewidth}{|l|X|X|X|}\hline
\textbf{Person} & \textbf{Affirmative} & \textbf{Negative} & \textbf{Interrogative} \\ \hline
All persons & I \textbf{will be working} & I \textbf{will not be working} & \textbf{Will I be working?} \\
& She \textbf{will be studying} & She \textbf{will not be studying} & \textbf{Will she be studying?} \\ \hline
\end{tabularx}
\end{center}
\textbf{Short answers:} \emph{Yes, I will. / No, I won't.} \quad \textbf{Contracted negative:} \emph{won't be working}.
\end{propriete}

\begin{exemplebox}[Conjugation of ``to repair'']
\begin{center}
\begin{tabularx}{\linewidth}{|l|X|X|}\hline
\textbf{Pronoun} & \textbf{Affirmative} & \textbf{Negative} \\ \hline
I & will be repairing & will not be repairing (won't be repairing) \\
You & will be repairing & will not be repairing \\
He / She / It & will be repairing & will not be repairing \\
We & will be repairing & will not be repairing \\
They & will be repairing & will not be repairing \\ \hline
\end{tabularx}
\end{center}
\end{exemplebox}

\courssub{Uses of the future continuous}

\begin{propriete}[Use 1: Action in progress at a precise future time]
The action starts before that moment and continues after it.
\begin{itemize}
    \item \emph{This time tomorrow, I \textbf{will be sitting} for my interview.} (Demain à cette heure-ci, je serai en train de passer mon entretien.)
    \item \emph{Next Monday at 9 a.m., Mballa \textbf{will be repairing} a truck engine.} (Lundi prochain à 9 h, Mballa sera en train de réparer un moteur de camion.)
    \item \emph{At 5 p.m. tonight, we \textbf{will be writing} our reports.} (Ce soir à 17 h, nous serons en train de rédiger nos rapports.)
    \item \emph{When you arrive, the team \textbf{will be testing} the prototype.} (Quand tu arriveras, l'équipe sera en train de tester le prototype.)
\end{itemize}
\end{propriete}

\begin{propriete}[Use 2: Polite questions about future plans]
The future continuous sounds softer and more polite than \emph{will you...?} or \emph{are you going to...?} because it asks about the schedule, not the intention.
\begin{itemize}
    \item \emph{\textbf{Will you be using} the workshop this afternoon?} (Vas-tu utiliser l'atelier cet après-midi ? — poli, pas insistant.)
    \item \emph{\textbf{Will you be attending} the safety training?} (Vas-tu assister à la formation sécurité ?)
    \item \emph{\textbf{Will the manager be joining} us for lunch?} (Le chef nous rejoindra-t-il pour le déjeuner ?)
    \item \emph{\textbf{Will you be needing} the vehicle tomorrow?} (Aurez-vous besoin du véhicule demain ?)
\end{itemize}
\end{propriete}

\begin{propriete}[Use 3: Predicting the present (less common, but useful)]
To say what we think is happening now, based on routine.
\begin{itemize}
    \item \emph{``Don't call the workshop now; they \textbf{will be having} their lunch break.''} (Ne téléphone pas à l'atelier maintenant ; ils seront en train de déjeuner.)
\end{itemize}
\end{propriete}

\begin{popfigure}
\begin{tikzpicture}[scale=0.9]
% Timeline for Future Continuous
\draw[popblue, thick, ->] (0,0) -- (13,0) node[right] {\textbf{Time}};
% Now
\draw[popgreen, thick] (1,0) -- (1,-0.3) node[below] {\textbf{NOW}};
\node[popgreen, anchor=north] at (1,-0.5) {Present moment};
% Future point
\draw[poporange, thick] (8,0) -- (8,-0.3) node[below] {\textbf{SPECIFIC FUTURE TIME}};
\node[poporange, anchor=north] at (8,-0.5) {e.g. Tomorrow 10 a.m.};
% Action in progress around that point
\draw[poppurple, very thick, ->] (6,0.5) -- (10,0.5) node[midway, above] {\textbf{Action in progress}};
% Dashed lines to show "around"
\draw[poppurple, dashed] (6,0) -- (6,0.5);
\draw[poppurple, dashed] (10,0) -- (10,0.5);
% Example box
\node[draw=poporange, fill=poporange!10, rounded corners, align=center, anchor=north] at (6.5,-2) {
    \textbf{Example:} \emph{``At 10 a.m. tomorrow,}\\
    \emph{I will be welding the pipe.''}\\
    (Demain à 10 h, je serai en train de souder le tuyau.)
};
\end{tikzpicture}
\legende{Future continuous timeline: a specific future moment (orange) with an action in progress around it (purple arrow). The action starts before and continues after that moment.}
\end{popfigure}

\courssub{Comparison: present perfect simple vs. present perfect continuous}

This is a key distinction for the Baccalauréat. Ask yourself: \textbf{is it the RESULT (how many/much) or the DURATION (how long) that matters?}

\begin{methode}[Choosing the right tense: result vs. duration]
\begin{enumerate}
    \item Identify the verb and the time marker (\emph{for}, \emph{since}, \emph{how long}, \emph{how many}, \emph{how much}).
    \item Ask: \textbf{``Do I want to say how many times / how much is done?''} $\rightarrow$ \textbf{Present perfect simple}.
    \item Ask: \textbf{``Do I want to say how long the activity has lasted?''} $\rightarrow$ \textbf{Present perfect continuous}.
    \item Check: some verbs (stative verbs like \emph{know, like, have, be}) are \textbf{not used} in continuous forms.
\end{enumerate}
\end{methode}

\begin{popfigure}
\begin{tikzpicture}
% Comparison table using TikZ matrix
\matrix[matrix of nodes, nodes={anchor=center, text width=3.2cm, minimum height=1.2cm}, 
    column sep=0.5mm, row sep=0.5mm,
    column 1/.style={nodes={fill=popblue, text=white, font=\bfseries, text width=2.8cm}},
    column 2/.style={nodes={fill=popblue!20, text width=4.5cm}},
    column 3/.style={nodes={fill=popblue!20, text width=4.5cm}},
    row 1/.style={nodes={fill=popdark, text=white, font=\bfseries, minimum height=1cm}}
] (comp) {
    \textbf{Aspect} & \textbf{Present perfect simple} & \textbf{Present perfect continuous} \\
    \textbf{Formation} & have/has + past participle & have/has + been + V-ing \\
    \textbf{Focus} & \textbf{Result / Completion} (combien ?) & \textbf{Duration / Activity} (depuis quand ?) \\
    \textbf{Question} & How many / How much? & How long? \\
    \textbf{Time markers} & already, just, yet, ever, never & for, since, all day, lately, recently \\
    \textbf{Example 1} & I \textbf{have written} 3 reports. (result: 3 reports done) & I \textbf{have been writing} reports all morning. (activity: duration) \\
    \textbf{Example 2} & She \textbf{has repaired} the car. (car is fixed now) & She \textbf{has been repairing} the car for 2 hours. (she is still working / just stopped) \\
    \textbf{Stative verbs} & OK: I have known him for years. & \textbf{NOT used:} *I have been knowing him... \\
};
\legende{Comparison: Present perfect simple (result-oriented) vs. Present perfect continuous (duration-oriented). Stative verbs (know, like, love, have, be, understand, want, need, believe) are not used in continuous forms.}
\end{tikzpicture}
\end{popfigure}

\begin{attention}[Stative verbs: never in continuous forms]
Verbs expressing states, not actions, are rarely used in any continuous tense (present continuous, present perfect continuous, future continuous, etc.).
\begin{center}
\begin{tabularx}{\linewidth}{|l|X|}\hline
\textbf{Category} & \textbf{Examples} \\ \hline
Mental states & know, understand, believe, think (opinion), remember, forget, mean \\
Emotions & like, love, hate, prefer, want, wish, need, desire \\
Possession & have (possess), own, belong, possess \\
Senses & see, hear, smell, taste, feel (perception) \\
Others & be, seem, appear, look (seem), cost, weigh, measure, contain, consist \\ \hline
\end{tabularx}
\end{center}
\textbf{Exception:} Some verbs have both stative and dynamic meanings.
\begin{itemize}
    \item \emph{have} (possess) $\rightarrow$ stative: \emph{I have a car.} \quad \emph{have} (action) $\rightarrow$ dynamic: \emph{I am having lunch / I have been having lunch.} (Je déjeune.)
    \item \emph{think} (opinion) $\rightarrow$ stative: \emph{I think it's good.} \quad \emph{think} (consider) $\rightarrow$ dynamic: \emph{I am thinking about the problem.}
    \item \emph{see} (perception) $\rightarrow$ stative: \emph{I see a bird.} \quad \emph{see} (meet) $\rightarrow$ dynamic: \emph{I am seeing the doctor tomorrow.}
\end{itemize}
\end{attention}

\begin{exoresolu}[Choosing the correct tense]
\textbf{Complete with the present perfect simple or continuous.}
\begin{enumerate}
    \item Look! Somebody \_\_\_\_\_\_\_\_\_\_ (break) the window. \emph{Result visible now.}
    \item I \_\_\_\_\_\_\_\_\_\_ (read) this manual for two hours. \emph{Duration, still reading.}
    \item How many pages \_\_\_\_\_\_\_\_\_\_ you \_\_\_\_\_\_\_\_\_\_ (write) so far? \emph{Quantity / result.}
    \item My hands are greasy because I \_\_\_\_\_\_\_\_\_\_ (repair) the generator. \emph{Recent action, visible result.}
    \item We \_\_\_\_\_\_\_\_\_\_ (know) the director since 2018. \emph{Stative verb.}
    \item The students \_\_\_\_\_\_\_\_\_\_ (wait) for the results since June. \emph{Duration, still waiting.}
    \item Who \_\_\_\_\_\_\_\_\_\_ (eat) my sandwich? \emph{Result: sandwich gone.}
    \item She \_\_\_\_\_\_\_\_\_\_ (type) three letters this morning. \emph{Quantity completed.}
\end{enumerate}
\tcblower
\textbf{Solution:}
\begin{enumerate}
    \item \textbf{has broken} (present perfect simple — result: the window is broken)
    \item \textbf{have been reading} (present perfect continuous — duration: for two hours)
    \item \textbf{have ... written} (present perfect simple — how many = quantity/result)
    \item \textbf{have been repairing} (present perfect continuous — recent activity, visible result: greasy hands)
    \item \textbf{have known} (present perfect simple — \emph{know} is a stative verb, no continuous)
    \item \textbf{have been waiting} (present perfect continuous — duration with \emph{since}, still waiting)
    \item \textbf{has eaten} (present perfect simple — result: sandwich is gone)
    \item \textbf{has typed} (present perfect simple — how many = quantity/result; \emph{this morning} is a finished period if it's now afternoon, but here the focus is on the number completed)
\end{enumerate}
\end{exoresolu}

\courssub{Application: your technical training and professional projects}

In a technical lycée in Cameroon, you can use these tenses to talk about your studies, internships (stages), and future career.

\begin{exemplebox}[Talking about your training (Present perfect continuous)]
\begin{itemize}
    \item \emph{``I \textbf{have been studying} electrical engineering \textbf{for} three years.''}
    \item \emph{``We \textbf{have been doing} our internship at SONEL \textbf{since} January.''}
    \item \emph{``The teacher \textbf{has been explaining} the circuit diagram \textbf{all morning}.''}
    \item \emph{``The apprentices \textbf{have been practising} welding \textbf{lately}.''}
\end{itemize}
\textbf{Questions you might hear:}
\begin{itemize}
    \item \emph{``How long \textbf{have you been learning} welding?''}
    \item \emph{``Have you \textbf{been working} on your final project lately?''}
    \item \emph{``What \textbf{have you been doing} since the last class?''}
\end{itemize}
\end{exemplebox}

\begin{exemplebox}[Talking about future plans (Future continuous)]
\begin{itemize}
    \item \emph{``Next year at this time, I \textbf{will be working} as a maintenance technician.''}
    \item \emph{``In June, we \textbf{will be sitting} for the Baccalauréat technique.''}
    \item \emph{``\textbf{Will you be using} the CNC machine tomorrow morning?''} (polite request)
    \item \emph{``The inspection team \textbf{will be visiting} the workshop next week.''}
    \item \emph{``I \textbf{will be presenting} my project on Friday at 10 a.m.''}
\end{itemize}
\end{exemplebox}

\begin{savaistu}[Future continuous in professional emails]
In the workplace (Cameroon and internationally), the future continuous is standard for polite enquiries and confirming schedules:
\begin{itemize}
    \item \emph{``\textbf{Will you be attending} the site meeting on Friday?''}
    \item \emph{``\textbf{Will the team be needing} the scaffolding next week?''}
    \item \emph{``I \textbf{will be sending} the technical specifications by Tuesday.''}
    \item \emph{``\textbf{Will you be available} for a call at 3 p.m.?''}
\end{itemize}
It sounds more professional and less direct than ``Will you attend?'' or ``Are you going to attend?''.
\end{savaistu}

\begin{exoresolu}[Mixed practice: your profile]
\textbf{Write three sentences about yourself using:}
\begin{enumerate}
    \item Present perfect continuous (your studies / training)
    \item Present perfect simple (a result you have achieved)
    \item Future continuous (your plan for a specific future moment)
\end{enumerate}
\tcblower
\textbf{Sample answers:}
\begin{enumerate}
    \item \emph{I have been studying mechanical engineering for four years.}
    \item \emph{I have completed my internship at the Douala port workshop.}
    \item \emph{This time next month, I will be presenting my final year project.}
\end{enumerate}
\end{exoresolu}

\begin{exercice}[Practice: tenses in context]
\textbf{Choose the correct tense (present perfect simple, present perfect continuous, or future continuous).}
\begin{enumerate}
    \item By 10 a.m. tomorrow, the technicians \_\_\_\_\_\_\_\_\_\_ (install) the new air-conditioning units.
    \item The apprentice \_\_\_\_\_\_\_\_\_\_ (learn) to read blueprints since September.
    \item How many cables \_\_\_\_\_\_\_\_\_\_ you \_\_\_\_\_\_\_\_\_\_ (connect) so far?
    \item \_\_\_\_\_\_\_\_\_\_ you \_\_\_\_\_\_\_\_\_\_ (use) the lathe this afternoon? I need it for an hour.
    \item The machine \_\_\_\_\_\_\_\_\_\_ (make) a strange noise all morning. We should check it.
    \item We \_\_\_\_\_\_\_\_\_\_ (finish) the wiring. The system is ready for testing.
    \item The trainee \_\_\_\_\_\_\_\_\_\_ (study) the safety manual for three hours. He \_\_\_\_\_\_\_\_\_\_ (not / finish) yet.
    \item \_\_\_\_\_\_\_\_\_\_ the supervisor \_\_\_\_\_\_\_\_\_\_ (inspect) the site when we arrive tomorrow?
    \item I \_\_\_\_\_\_\_\_\_\_ (know) the chief engineer for ten years.
    \item Look at the oil on the floor! The pump \_\_\_\_\_\_\_\_\_\_ (leak).
\end{enumerate}
\end{exercice}

\begin{corrige}[Exercice n°1]
\begin{enumerate}
    \item \textbf{will be installing} (future continuous — action in progress at a specific future time: ``By 10 a.m. tomorrow'')
    \item \textbf{has been learning} (present perfect continuous — duration with ``since September'')
    \item \textbf{have ... connected} (present perfect simple — ``How many'' = quantity/result)
    \item \textbf{Will ... be using} (future continuous — polite question about plans)
    \item \textbf{has been making} (present perfect continuous — duration ``all morning'', activity still relevant)
    \item \textbf{have finished} (present perfect simple — completed action with present result: ``system is ready'')
    \item \textbf{has been studying} / \textbf{has not finished} (present perfect continuous for duration + present perfect simple for negative result)
    \item \textbf{Will ... be inspecting} (future continuous — action in progress at a future moment ``when we arrive'')
    \item \textbf{have known} (present perfect simple — \emph{know} is stative)
    \item \textbf{has been leaking} (present perfect continuous — recent activity with visible present result: oil on floor)
\end{enumerate}
\end{corrige}

\begin{aretenir}[Key points to remember]
\begin{itemize}
    \item \textbf{Present perfect continuous} = \emph{have/has been + V-ing}. Focus on \textbf{duration} (\emph{for, since, how long}). Action started in past, continues now or just stopped.
    \item \textbf{Future continuous} = \emph{will be + V-ing}. Focus on \textbf{action in progress at a specific future moment}. Also used for \textbf{polite questions} about arrangements.
    \item \textbf{Present perfect simple} = result / quantity (\emph{how many/much}). \textbf{Present perfect continuous} = duration / activity (\emph{how long}).
    \item \textbf{Stative verbs} (know, like, have, be, understand, want, need, believe, see, hear...) $\rightarrow$ never in continuous forms (except dynamic meanings).
    \item \textbf{Error to avoid:} \emph{*I am working here since 2020} $\rightarrow$ \textbf{I have been working here since 2020.}
    \item \textbf{Error to avoid:} \emph{*I have been knowing him for years} $\rightarrow$ \textbf{I have known him for years.}
\end{itemize}
\end{aretenir}

\coursec{Embedded questions, question tags \& phrasal verbs}

In the previous section, we mastered the \textbf{present perfect continuous} and the \textbf{future continuous} to describe ongoing actions and future plans. Now, we move from \emph{when} actions happen to \emph{how} we ask about them politely and naturally in a professional context. Mastering \textbf{embedded questions} makes you sound courteous in interviews; \textbf{question tags} help you check information and build rapport; and \textbf{phrasal verbs} are the key to understanding everyday workplace English.

\courssub{Embedded Statements and Questions}

\begin{definition}[Embedded Question]
An \cle{embedded question} (question incluse / question indirecte) is a question included inside another sentence or question. Unlike a direct question, it uses \textbf{statement word order} (Subject + Verb), not inversion. It does not use the auxiliaries \textbf{do / does / did}.
\end{definition}

\begin{propriete}[Rule: No Auxiliary Do/Does/Did in Embedded Questions]
In an embedded question, the verb follows the subject directly. We never keep the question auxiliary \emph{do, does, did}.
\begin{center}
\begin{tabularx}{\linewidth}{|l|X|X|}\hline
\rowcolor{popblueL}\textbf{Direct Question} & \textbf{Embedded Question} & \textbf{Structure} \\ \hline
Where \textbf{is} the office? & Could you tell me where the office \textbf{is}? & Wh-word + Subject + Verb \\ \hline
What time \textbf{does} the interview \textbf{start}? & I wonder what time the interview \textbf{starts}. & Wh-word + Subject + Verb (add \emph{s} for 3rd pers. sg.) \\ \hline
Did you \textbf{send} the CV? & Do you know \textbf{if} you \textbf{sent} the CV? & \emph{If / Whether} + Subject + Verb (past tense) \\ \hline
\end{tabularx}
\end{center}
\end{propriete}

\begin{popfigure}
\begin{tikzpicture}[
    grow'=down,
    level distance=2.2cm,
    every node/.style={draw, rounded corners, fill=popblueL, thick, align=center, font=\small},
    edge from parent/.style={draw, popblue, thick, -{Stealth[length=3mm]}}
]
\node[align=center]{Direct Question\\ \emph{``Where is the office?''}\\ \textcolor{poporange}{Verb + Subject}}
    child { node[align=center]{Step 1: Add Introductory Phrase\\ \emph{``Could you tell me...''}}
        child { node[align=center]{Step 2: Change Order\\ \emph{``...where the office \textbf{is}?''}\\ \textcolor{popgreen}{Subject + Verb}}
            child { node[fill=popgreenL, align=center]{Final Embedded Question\\ \emph{``Could you tell me where the office is?''}} }
        }
    };
\end{tikzpicture}
\legende{Transformation tree: Direct Question $\rightarrow$ Embedded Question. Rule: (1) keep the Wh-word, (2) remove the auxiliary \emph{do/does/did} and the inversion, (3) apply Subject + Verb order, (4) no question mark if the main clause is a statement.}
\end{popfigure}

\textbf{1. Embedded Statements (Reported Speech / Thoughts)}\par
We introduce a statement with verbs like \emph{think, believe, say, tell, know, hope}. The conjunction \emph{that} is optional.
\begin{exemplebox}[Embedded Statements]
\begin{itemize}
    \item \textbf{Direct:} ``English \textbf{is} important.'' $\rightarrow$ \textbf{Embedded:} I think (\textbf{that}) English \textbf{is} important.
    \item \textbf{Direct:} ``He \textbf{was} ready.'' $\rightarrow$ \textbf{Embedded:} He said (\textbf{that}) he \textbf{was} ready. (Backshifting of tense often occurs after a past reporting verb.)
    \item \textbf{Direct:} ``We \textbf{will} finish on time.'' $\rightarrow$ \textbf{Embedded:} She hopes (\textbf{that}) we \textbf{will} finish on time.
\end{itemize}
\end{exemplebox}

\textbf{2. Embedded Wh-Questions}\par
Used after: \emph{ask, wonder, want to know, know, remember, forget, understand, explain, tell me, show me, Could you tell me...}
\begin{exemplebox}[Embedded Wh-Questions]
\begin{itemize}
    \item \textbf{Direct:} ``\textbf{Where} \textbf{do} you \textbf{work}?'' $\rightarrow$ \textbf{Embedded:} Could you tell me \textbf{where} you \textbf{work}? (\emph{do} disappears; the main verb keeps its normal form.)
    \item \textbf{Direct:} ``\textbf{What} \textbf{are} your \textbf{strengths}?'' $\rightarrow$ \textbf{Embedded:} I'd like to know \textbf{what} your \textbf{strengths are}.
    \item \textbf{Direct:} ``\textbf{Who} \textbf{called}?'' $\rightarrow$ \textbf{Embedded:} Do you know \textbf{who} \textbf{called}? (Subject question: the word order does not change.)
\end{itemize}
\end{exemplebox}

\textbf{3. Embedded Yes/No Questions}\par
Used when there is no Wh-word. We use \textbf{if} or \textbf{whether} (\emph{whether} is more formal).
\begin{exemplebox}[Embedded Yes/No Questions]
\begin{itemize}
    \item \textbf{Direct:} ``\textbf{Does} the interview \textbf{start} at 8?'' $\rightarrow$ \textbf{Embedded:} Do you know \textbf{if} the interview \textbf{starts} at 8?
    \item \textbf{Direct:} ``\textbf{Is} the manager \textbf{here}?'' $\rightarrow$ \textbf{Embedded:} I wonder \textbf{whether} the manager \textbf{is} here.
    \item \textbf{Direct:} ``\textbf{Did} you \textbf{receive} my email?'' $\rightarrow$ \textbf{Embedded:} Please tell me \textbf{if} you \textbf{received} my email.
\end{itemize}
\end{exemplebox}

\begin{attention}[Common Trap: Keeping the Inversion]
\textbf{Error:} ``Could you tell me where \emph{is the office}?'' $\rightarrow$ \textbf{Correction:} ``Could you tell me where \textbf{the office is}?''
\textbf{Reason:} After the introductory phrase, the question becomes a normal statement: Subject + Verb. The question mark at the end belongs to the main clause (\emph{Could you tell me...?}), not to the embedded part.
\end{attention}

\begin{exoresolu}[Transforming Direct Questions into Embedded Questions]
\textbf{Task:} Rewrite these direct questions as embedded questions starting with the words in brackets.
\begin{enumerate}
    \item ``What are your salary expectations?'' (Could you tell me...)
    \item ``Did you bring your certificates?'' (The interviewer asked me...)
    \item ``Where is the nearest bus stop?'' (I'd like to know...)
    \item ``Can you work on weekends?'' (She wants to know whether...)
\end{enumerate}

\tcblower
\textbf{Solution:}
\begin{enumerate}
    \item Could you tell me \textbf{what your salary expectations are}? (Order: Subject \emph{expectations} + Verb \emph{are}.)
    \item The interviewer asked me \textbf{if I had brought} my certificates. (Backshift after a past reporting verb: \emph{did bring} $\rightarrow$ \emph{had brought}; \emph{if} for a Yes/No question.)
    \item I'd like to know \textbf{where the nearest bus stop is}. (Order: Subject \emph{bus stop} + Verb \emph{is}.)
    \item She wants to know \textbf{whether you can work} on weekends. (\emph{whether} for a Yes/No question; modal \emph{can} stays before the subject's verb but after the subject.)
\end{enumerate}
\end{exoresolu}

\courssub{Question Tags}

\begin{definition}[Question Tag]
A \cle{question tag} (question tag / petite question finale) is a short question added at the end of a statement to check information or invite agreement: \emph{You speak English, \textbf{don't you}?}
\end{definition}

\begin{methode}[How to Form a Question Tag]
Follow these 3 steps to build a correct tag:
\begin{enumerate}
    \item \textbf{Identify the Auxiliary:} Find the auxiliary verb in the main sentence (\emph{be, have, do, can, will, must, should...}). If there is no auxiliary (Present Simple / Past Simple affirmative), use \emph{do / does / did}.
    \item \textbf{Invert the Polarity:}
        \begin{itemize}
            \item Positive statement $\rightarrow$ \textbf{Negative Tag} (\emph{aren't, don't, won't, can't}).
            \item Negative statement $\rightarrow$ \textbf{Positive Tag} (\emph{are, do, will, can}).
        \end{itemize}
    \item \textbf{Match the Subject Pronoun:} Replace the subject noun with the corresponding pronoun (\emph{he, she, it, they, we, you, I}).
\end{enumerate}
\end{methode}

\begin{center}
\begin{tabular}{|l|l|l|l|l|}\hline
\rowcolor{popblueL}\textbf{Main Sentence} & \textbf{Subject} & \textbf{Auxiliary / Verb} & \textbf{Tag} & \textbf{Polarity} \\ \hline
You speak English, & You & speak (\emph{do}) & \textbf{don't you?} & $+ \rightarrow -$ \\ \hline
She isn't late, & She & isn't (\emph{is not}) & \textbf{is she?} & $- \rightarrow +$ \\ \hline
They have finished, & They & have & \textbf{haven't they?} & $+ \rightarrow -$ \\ \hline
We won't fail, & We & won't (\emph{will not}) & \textbf{will we?} & $- \rightarrow +$ \\ \hline
It works well, & It & works (\emph{does}) & \textbf{doesn't it?} & $+ \rightarrow -$ \\ \hline
He didn't call, & He & didn't (\emph{did not}) & \textbf{did he?} & $- \rightarrow +$ \\ \hline
\end{tabular}
\end{center}

\textbf{Special Cases (must know for the Probatoire/Bac technique):}
\begin{aretenir}[Special Question Tags]
\begin{center}
\begin{tabularx}{\linewidth}{|l|l|X|}\hline
\rowcolor{popgreenL}\textbf{Structure} & \textbf{Tag} & \textbf{Example} \\ \hline
\textbf{I am} & \textbf{aren't I?} & I am late, \textbf{aren't I?} (Never ``amn't I''.) \\ \hline
\textbf{Let's} (Let us) & \textbf{shall we?} & Let's start, \textbf{shall we?} \\ \hline
\textbf{Imperative} (Positive) & \textbf{will you? / won't you?} & Open the door, \textbf{will you?} \\ \hline
\textbf{Imperative} (Negative) & \textbf{will you?} & Don't forget, \textbf{will you?} \\ \hline
\textbf{There is / There are} & \textbf{isn't there? / aren't there?} & There is a meeting, \textbf{isn't there?} \\ \hline
\textbf{Nobody / Somebody / Everybody} & \textbf{they?} & Nobody came, \textbf{did they?} \\ \hline
\textbf{This / That is} & \textbf{isn't it?} & This is yours, \textbf{isn't it?} \\ \hline
\textbf{These / Those are} & \textbf{aren't they?} & These are yours, \textbf{aren't they?} \\ \hline
\end{tabularx}
\end{center}
\end{aretenir}

\begin{attention}[Common Trap: ``Isn't it?'' for Everything]
\textbf{Error:} ``You are coming, \emph{isn't it?}'' \quad $\rightarrow$ \textbf{Correction:} ``You are coming, \textbf{aren't you?}''
\textbf{Reason:} The tag \textbf{must copy the subject and the auxiliary} of the main clause. ``It'' is only used for impersonal structures (\emph{It is raining, isn't it?}) or demonstratives (\emph{This is...}). For \emph{You}, use \emph{you}; for \emph{He}, use \emph{he}.
\end{attention}

\textbf{Intonation and Meaning:}
\begin{itemize}
    \item \textbf{Falling Intonation} ($\searrow$): the speaker is \textbf{sure} and expects agreement. ``It's a nice day, isn't it?'' ($\searrow$ : confirmation).
    \item \textbf{Rising Intonation} ($\nearrow$): the speaker is \textbf{not sure} and wants real information. ``You haven't seen my pen, have you?'' ($\nearrow$ : real question).
\end{itemize}

\begin{exemplebox}[Question Tags in Context]
\begin{itemize}
    \item ``You have prepared your CV, \textbf{haven't you?}'' (Falling: checking preparation.)
    \item ``You don't know the manager, \textbf{do you?}'' (Rising: hoping for a contact.)
    \item ``Let's practise the interview, \textbf{shall we?}'' (Suggestion.)
    \item ``I'm the best candidate, \textbf{aren't I?}'' (Confidence, falling intonation.)
\end{itemize}
\end{exemplebox}

\courssub{Phrasal Verbs and Prepositional Phrases at Work}

Phrasal verbs (Verb + Particle/Preposition) are extremely common in spoken English and professional interactions. Using them correctly shows a high level of fluency.

\begin{definition}[Phrasal Verb]
A \cle{phrasal verb} (verbe à particule) combines a verb with a particle (adverb or preposition), creating a meaning often different from the original verb. Some are \textbf{separable} (the object can go in the middle), others are \textbf{inseparable}.
\end{definition}

\begin{popfigure}
\begin{tikzpicture}[
    node distance=0.4cm and 0.6cm,
    verb/.style={draw, fill=popblueL, rounded corners, thick, minimum width=1.6cm, minimum height=0.8cm, align=center, font=\small},
    part/.style={draw, fill=poporangeL, rounded corners, thick, minimum width=1.6cm, minimum height=0.8cm, align=center, font=\small},
    mean/.style={draw, fill=popgreenL, rounded corners, thick, minimum width=3.4cm, minimum height=0.8cm, align=center, font=\small},
    arrow/.style={-Stealth, thick, popdark}
]
\node[verb] (v1) {fill};
\node[part, right=of v1] (p1) {in};
\node[mean, right=of p1] (m1) {complete (a form)};
\draw[arrow] (v1) -- (p1) node[midway, above, font=\tiny] {+};
\draw[arrow] (p1) -- (m1) node[midway, above, font=\tiny] {=};

\node[verb, below=of v1] (v2) {hand};
\node[part, right=of v2] (p2) {in};
\node[mean, right=of p2] (m2) {submit (CV, report)};
\draw[arrow] (v2) -- (p2) node[midway, above, font=\tiny] {+};
\draw[arrow] (p2) -- (m2) node[midway, above, font=\tiny] {=};

\node[verb, below=of v2] (v3) {look};
\node[part, right=of v3] (p3) {for};
\node[mean, right=of p3] (m3) {search / seek};
\draw[arrow] (v3) -- (p3) node[midway, above, font=\tiny] {+};
\draw[arrow] (p3) -- (m3) node[midway, above, font=\tiny] {=};

\node[verb, below=of v3] (v4) {turn};
\node[part, right=of v4] (p4) {up};
\node[mean, right=of p4] (m4) {arrive / appear};
\draw[arrow] (v4) -- (p4) node[midway, above, font=\tiny] {+};
\draw[arrow] (p4) -- (m4) node[midway, above, font=\tiny] {=};

\node[verb, below=of v4] (v5) {get};
\node[part, right=of v5] (p5) {along with};
\node[mean, right=of p5] (m5) {have good relations};
\draw[arrow] (v5) -- (p5) node[midway, above, font=\tiny] {+};
\draw[arrow] (p5) -- (m5) node[midway, above, font=\tiny] {=};
\end{tikzpicture}
\legende{Phrasal verb structure: Verb + Particle = New Meaning. \textbf{Separable:} \emph{fill it in} or \emph{fill in the form}. \textbf{Inseparable:} \emph{look for it} (never ``look it for'').}
\end{popfigure}

\begin{center}
\begin{tabularx}{\linewidth}{|>{\bfseries}l|X|>{\itshape}X|}\hline
\rowcolor{popblueL}\textnormal{\textbf{Phrasal Verb}} & \textnormal{\textbf{Meaning (Français)}} & \textnormal{\textbf{Example Context}} \\ \hline
carry out & effectuer, réaliser (une tâche) & The technician must \textbf{carry out} the maintenance check. \\ \hline
fill in & remplir (un formulaire) & Please \textbf{fill in} this application form. \\ \hline
hand in & remettre, déposer (dossier, CV) & You must \textbf{hand in} your report by Friday. \\ \hline
look for & chercher & We are \textbf{looking for} a new accountant. \\ \hline
put on & enfiler, mettre (vêtements, équipement de protection) & \textbf{Put on} your safety helmet before entering. \\ \hline
take off & enlever (vêtements) / décoller (avion) & You can \textbf{take off} your jacket in the office. \\ \hline
turn up & arriver, se présenter & He \textbf{turned up} late for the interview. \\ \hline
work out & résoudre, calculer / s'entraîner & We need to \textbf{work out} the total cost. \\ \hline
get along with & bien s'entendre avec & I \textbf{get along well with} my colleagues. \\ \hline
look forward to & avoir hâte de (+\emph{-ing}) & I \textbf{look forward to hearing} from you. \\ \hline
\end{tabularx}
\end{center}

\begin{attention}[Trap: ``look forward to'' + -ing]
After \textbf{look forward to}, the next verb takes the \textbf{-ing} form, because \emph{to} here is a preposition, not an infinitive marker: ``I look forward to \textbf{hearing} from you'' (never ``to hear'').
\end{attention}

\textbf{Prepositional Phrases (Fixed Expressions):}
\begin{center}
\begin{tabularx}{\linewidth}{|>{\bfseries}l|X|>{\itshape}X|}\hline
\rowcolor{popgreenL}\textnormal{\textbf{Phrase}} & \textnormal{\textbf{Meaning (Français)}} & \textnormal{\textbf{Example}} \\ \hline
in charge of & responsable de & Who is \textbf{in charge of} the IT department? \\ \hline
on time & à l'heure, ponctuel(le) & The meeting started \textbf{on time}. \\ \hline
at work & au travail & She is \textbf{at work} until 5 p.m. \\ \hline
for a living & pour gagner sa vie & What do you do \textbf{for a living}? \\ \hline
by heart & par cœur & Know the safety rules \textbf{by heart}. \\ \hline
in the long run & à long terme & Training pays off \textbf{in the long run}. \\ \hline
\end{tabularx}
\end{center}

\begin{aretenir}[Worked Example: First Day at a Plant in Edéa]
\textbf{Context:} Junior technician \emph{Paul} starts work at a factory in \emph{Edéa} (Littoral Region, Cameroon).
\begin{quote}
``This morning I \textbf{turned up} early. I had to \textbf{put on} my safety boots and \textbf{fill in} a security form at the gate. My supervisor, Mr. Ngono, is \textbf{in charge of} the maintenance unit. He asked me to \textbf{look for} the technical manuals and to \textbf{hand in} my CV at the HR office. Then I had to \textbf{work out} the machine settings. I hope I will \textbf{get along with} the team. I am \textbf{looking forward to} learning this trade \textbf{for a living}. \textbf{In the long run}, I want to become a team leader.''
\end{quote}
\textbf{Count:} 10 phrasal verbs and prepositional phrases used naturally in a short professional narrative.
\end{aretenir}

\begin{savaistu}[Phrasal Verbs in Interviews]
Phrasal verbs dominate \textbf{spoken, everyday English}. In a job interview, saying ``I \textbf{carried out} a project'' often sounds more natural than a very formal equivalent. Recruiters listen for them. \textbf{Tip:} learn phrasal verbs by topic (e.g., ``interview verbs'': \emph{turn up, hand in, fill in, look forward to}).
\end{savaistu}

\begin{exoresolu}[Phrasal Verbs Gap Fill]
\textbf{Complete with the correct phrasal verb (right tense) or prepositional phrase:} \emph{carry out, fill in, hand in, look for, put on, turn up, work out, get along with, on time, in charge of.}
\begin{enumerate}
    \item The manager is \_\_\_\_\_\_\_\_\_\_ the recruitment process.
    \item Please \_\_\_\_\_\_\_\_\_\_ this form in capital letters.
    \item If you \_\_\_\_\_\_\_\_\_\_ late for the interview, you lose points.
    \item We need to \_\_\_\_\_\_\_\_\_\_ the total budget for the year.
    \item She \_\_\_\_\_\_\_\_\_\_ her protective goggles before welding.
    \item I \_\_\_\_\_\_\_\_\_\_ my colleagues very well; we are a great team.
    \item The project was \_\_\_\_\_\_\_\_\_\_ successfully last month.
    \item You must \_\_\_\_\_\_\_\_\_\_ your report by 4 p.m. today.
\end{enumerate}

\tcblower
\textbf{Answers:}
\begin{enumerate}
    \item \textbf{in charge of} (prepositional phrase: responsible for)
    \item \textbf{fill in} (separable: \emph{fill it in})
    \item \textbf{turn up} (arrive)
    \item \textbf{work out} (calculate)
    \item \textbf{put on} (wore protective equipment; past tense)
    \item \textbf{get along with} (good relationship)
    \item \textbf{carried out} (passive, past tense: executed)
    \item \textbf{hand in} (submit)
\end{enumerate}
\end{exoresolu}

\courssub{Practice Exercises}

\begin{exercice}[Embedded Questions Transformation]
Rewrite the direct questions as polite embedded questions.
\begin{enumerate}
    \item ``Where is the personnel department?'' $\rightarrow$ Could you tell me \_\_\_\_\_\_\_\_\_\_?
    \item ``Does the company provide housing?'' $\rightarrow$ I'd like to know \_\_\_\_\_\_\_\_\_\_.
    \item ``What qualifications do you need?'' $\rightarrow$ Can you explain \_\_\_\_\_\_\_\_\_\_?
    \item ``Did the previous employee leave?'' $\rightarrow$ Do you know \_\_\_\_\_\_\_\_\_\_?
    \item ``How much does it pay?'' $\rightarrow$ I wonder \_\_\_\_\_\_\_\_\_\_.
\end{enumerate}
\end{exercice}

\begin{exercice}[Question Tags Completion]
Add the correct question tag.
\begin{enumerate}
    \item You have finished the report, \_\_\_\_\_\_\_\_\_\_?
    \item She can't drive a forklift, \_\_\_\_\_\_\_\_\_\_?
    \item Let's take a break, \_\_\_\_\_\_\_\_\_\_?
    \item I am invited to the meeting, \_\_\_\_\_\_\_\_\_\_?
    \item There aren't any vacancies, \_\_\_\_\_\_\_\_\_\_?
    \item Nobody answered the phone, \_\_\_\_\_\_\_\_\_\_?
    \item Don't be late tomorrow, \_\_\_\_\_\_\_\_\_\_?
    \item This is your first application, \_\_\_\_\_\_\_\_\_\_?
\end{enumerate}
\end{exercice}

\begin{exercice}[Phrasal Verbs \& Prepositional Phrases in Context]
Choose the best option (A, B, or C) to complete the sentences.
\begin{enumerate}
    \item The safety officer is \_\_\_\_\_\_\_\_\_\_ checking the fire extinguishers.
        \begin{itemize}
            \item[A)] in charge of \quad [B)] on time \quad [C)] by heart
        \end{itemize}
    \item We must \_\_\_\_\_\_\_\_\_\_ the emergency procedures \_\_\_\_\_\_\_\_\_\_.
        \begin{itemize}
            \item[A)] learn / by heart \quad [B)] look / for \quad [C)] fill / in
        \end{itemize}
    \item The interview is at 9:00. Please be \_\_\_\_\_\_\_\_\_\_.
        \begin{itemize}
            \item[A)] at work \quad [B)] on time \quad [C)] for a living
        \end{itemize}
    \item I \_\_\_\_\_\_\_\_\_\_ meeting the new director.
        \begin{itemize}
            \item[A)] look forward to \quad [B)] get along with \quad [C)] turn up
        \end{itemize}
    \item He \_\_\_\_\_\_\_\_\_\_ a solution to the technical problem.
        \begin{itemize}
            \item[A)] worked out \quad [B)] carried out \quad [C)] handed in
        \end{itemize}
\end{enumerate}
\end{exercice}

\begin{corrige}[Exercice 1: Embedded Questions]
\begin{enumerate}
    \item Could you tell me \textbf{where the personnel department is}? (Subject + Verb order.)
    \item I'd like to know \textbf{if/whether the company provides housing}. (\emph{If/whether} + Subject + Verb; \emph{provides} takes \emph{s}.)
    \item Can you explain \textbf{what qualifications you need}? (Subject + Verb order, no \emph{do}.)
    \item Do you know \textbf{if/whether the previous employee left}? (\emph{If/whether} + past tense.)
    \item I wonder \textbf{how much it pays}. (Subject + Verb order, no \emph{does}.)
\end{enumerate}
\end{corrige}

\begin{corrige}[Exercice 2: Question Tags]
\begin{enumerate}
    \item \textbf{haven't you?} (Positive statement with \emph{have} $\rightarrow$ negative tag.)
    \item \textbf{can she?} (Negative \emph{can't} $\rightarrow$ positive tag.)
    \item \textbf{shall we?} (\emph{Let's} $\rightarrow$ \emph{shall we}.)
    \item \textbf{aren't I?} (\emph{I am} $\rightarrow$ \emph{aren't I}.)
    \item \textbf{are there?} (Negative \emph{aren't} $\rightarrow$ positive tag; subject \emph{there}.)
    \item \textbf{did they?} (\emph{Nobody} is negative in meaning $\rightarrow$ positive tag; pronoun \emph{they}.)
    \item \textbf{will you?} (Negative imperative $\rightarrow$ \emph{will you}.)
    \item \textbf{isn't it?} (\emph{This is} $\rightarrow$ \emph{isn't it}.)
\end{enumerate}
\end{corrige}

\begin{corrige}[Exercice 3: Phrasal Verbs \& Phrases]
\begin{enumerate}
    \item \textbf{A} (\emph{in charge of} = responsable de.)
    \item \textbf{A} (\emph{learn by heart} = apprendre par cœur.)
    \item \textbf{B} (\emph{on time} = à l'heure.)
    \item \textbf{A} (\emph{look forward to} + V-ing = avoir hâte de.)
    \item \textbf{A} (\emph{work out} = résoudre, trouver une solution.)
\end{enumerate}
\end{corrige}

\coursec{Anticipating second-language needs at work (Reading \& Vocabulary)}

\courssub{Transition from the previous section}

In the previous section, we studied \cle{embedded questions}, \cle{question tags} and \cle{phrasal verbs}. These tools help us ask for information politely (\emph{Could you tell me what time the interview starts?}), check understanding (\emph{You speak English, don't you?}) and express actions naturally (\emph{fill in a form, look up a word}). Now we apply these skills to a real professional situation: identifying \textbf{why you need English in your future job} and \textbf{how to improve it}.

\begin{aretenir}[Why this section matters]
As a technical student in Cameroon, you are preparing for a bilingual labour market. English is not just a school subject — it is a \cle{working tool} for reading manuals, writing reports, serving customers and following training. This section helps you \textbf{analyse your own language needs} and \textbf{build a personal improvement plan}.
\end{aretenir}

\courssub{Why English is indispensable at work in Cameroon}

Cameroon is officially \cle{bilingual} (French and English) by \textbf{Article 1(3) of the Constitution}. This unique status creates daily professional situations where English is required:

\begin{itemize}
    \item \textbf{Geography:} Cameroon borders Nigeria and is close to Ghana — both English-speaking economic partners.
    \item \textbf{International companies:} Many firms (oil, telecom, banking, logistics) use English as their corporate language.
    \item \textbf{Technical documentation:} Machine manuals, safety datasheets, software interfaces, coding languages — most are in English.
    \item \textbf{Regional integration:} CEMAC, ECCAS, African Continental Free Trade Area (AfCFTA) — English is a working language.
\end{itemize}

\begin{savaistu}[Le bilinguisme constitutionnel]
L'article 1(3) de la Constitution de la République du Cameroun dispose : \emph{« Les langues officielles de la République du Cameroun sont l'anglais et le français, qui ont le même statut. »} Cet atout juridique fait du Cameroun l'un des rares pays africains où le bilinguisme officiel est une norme constitutionnelle, et non seulement une pratique administrative. Sur le marché du travail, un technicien qui maîtrise les deux langues officielles a un avantage concurrentiel majeur.
\end{savaistu}

\courssub{Key vocabulary: language needs at work}

\begin{definition}[Language needs]
\cle{Language needs} (besoins linguistiques) are the specific language skills a person requires to perform tasks in a given professional situation. They vary according to the job, the sector and the level of responsibility.
\end{definition}

\begin{center}
\begin{tabularx}{\linewidth}{|l|X|X|}
\hline
\rowcolor{popblueL}
\textbf{English term} & \textbf{Meaning / French gloss} & \textbf{Example in context} \\
\hline
\cle{language needs} & besoins linguistiques & \emph{My language needs include reading technical manuals.} \\
\cle{communication skills} & compétences en communication & \emph{Good communication skills are essential for customer service.} \\
\cle{technical terms} & termes techniques & \emph{The manual is full of technical terms like 'torque' and 'voltage'.} \\
\cle{manual} & notice, manuel d'utilisation & \emph{Please read the manual before operating the machine.} \\
\cle{report writing} & rédaction de rapports & \emph{Report writing is a key skill for supervisors.} \\
\cle{customer service} & service client & \emph{She works in customer service at a telecom company.} \\
\cle{training} & formation & \emph{The company offers training in safety procedures.} \\
\cle{proficiency} & maîtrise, compétence avérée & \emph{The job requires proficiency in English (level B2).} \\
\cle{fluent} & couramment (parler couramment) & \emph{He is fluent in English and French.} \\
\cle{bilingual} & bilingue (parler deux langues) & \emph{Cameroon is a bilingual country.} \\
\hline
\end{tabularx}
\end{center}

\begin{attention}[Fluent vs. Bilingual — piège fréquent]
\cle{Fluent} signifie « parler une langue couramment, avec aisance ». \cle{Bilingual} signifie « utiliser deux langues dans la vie quotidienne ou professionnelle ». On peut être \textbf{bilingue sans être perfectly fluent} dans les deux langues (ex. : un Camerounais qui parle français couramment et anglais à un niveau intermédiaire). Inversement, on peut être \textbf{fluent en anglais sans être bilingue} (ex. : un Français qui parle anglais couramment mais ne pratique pas le français au travail). Au Cameroun, l'objectif professionnel est souvent les deux : \textbf{bilingual \emph{and} fluent}.
\end{attention}

\courssub{Reading strategies: skimming, scanning, deducing meaning}

Before reading the text, let's recall three essential strategies for efficient reading in a professional context.

\begin{methode}[Lire efficacement un texte professionnel]
\begin{enumerate}
    \item \textbf{Skimming (lecture rapide pour l'idée générale) :} Lisez le titre, les sous-titres, la première et la dernière phrase de chaque paragraphe. Ne vous arrêtez pas sur les mots inconnus. Objectif : \emph{De quoi parle le texte ?}
    \item \textbf{Scanning (recherche d'informations précises) :} Cherchez des mots-clés (noms, chiffres, dates, termes techniques) pour répondre à une question précise. Objectif : \emph{Trouver X information.}
    \item \textbf{Relecture pour le vocabulaire (deducing meaning from context) :} Relisez les phrases contenant des mots nouveaux. Utilisez le contexte (mots voisins, connecteurs logiques, exemples) pour deviner le sens avant de consulter le dictionnaire.
\end{enumerate}
\end{methode}

\courssub{Reading text: A technician's need for English at the Douala seaport}

\begin{popfigure}
\textbf{Text: A technician's need for English at the Douala seaport}

\medskip
\noindent
\textbf{Paragraph 1:} 
Douala Seaport is the largest port in Central Africa. It handles more than 90\% of Cameroon's foreign trade. Every day, hundreds of containers arrive from Europe, Asia and the Americas. The port operates 24 hours a day, seven days a week. Technicians like \textbf{Mr. Nguema}, a 32-year-old electromechanical technician, are essential to keep the cranes, conveyor belts and power generators running.

\medskip
\noindent
\textbf{Paragraph 2:}
Mr. Nguema's \cle{language needs} are varied. First, he must \textbf{read technical manuals} in English. The port's main cranes are manufactured in Germany and the USA. Their maintenance manuals, error codes and safety instructions are written in English. Last month, a crane stopped working because of a hydraulic fault. Mr. Nguema had to consult the English manual to identify the error code \texttt{E-407} and follow the troubleshooting steps. Without English, the repair would have taken days instead of hours.

\medskip
\noindent
\textbf{Paragraph 3:}
Second, Mr. Nguema often \textbf{communicates with English-speaking clients and colleagues}. Shipping agents from Nigeria and Ghana come to the port regularly. They need information about container handling, storage conditions and loading schedules. Mr. Nguema uses English to explain technical constraints, give instructions to dockworkers and answer safety questions. He says: \emph{« When a Nigerian agent asks about the maximum weight for a 40-foot container, I must answer clearly and quickly. Mistakes can cause accidents. »}

\medskip
\noindent
\textbf{Paragraph 4:}
Third, Mr. Nguema writes \textbf{daily maintenance reports} and \textbf{incident reports} in English. The port's management software is in English. All digital logs, work orders and safety checklists are in English. A poorly written report can lead to misunderstandings, delayed repairs or legal problems. Mr. Nguema spends about 45 minutes each day writing reports.

\medskip
\noindent
\textbf{Paragraph 5:}
Finally, Mr. Nguema follows \textbf{online training courses} to upgrade his skills. Most specialised courses on PLC programming, variable frequency drives and industrial safety are available only in English. He recently completed a 30-hour course on \emph{Siemens S7-1500 PLC programming} entirely in English. He says: \emph{« English is not a subject for me. It is the key to my professional development. »}

\medskip
\noindent
\textbf{Paragraph 6:}
Mr. Nguema's experience shows that for a technician in Cameroon, English is a \cle{working tool}, not just an exam subject. His \cle{proficiency} in English directly affects his \cle{efficiency}, his \cle{safety} record and his \cle{promotion} prospects. Identifying your own \cle{language needs} is the first step to building a realistic improvement plan.
\legende{Authentic-style professional text illustrating language needs of a technician at Douala Seaport.}
\end{popfigure}

\courssub{Comprehension check: solved questionnaire}

\begin{exoresolu}[Questionnaire de lecture — 5 questions]
\textbf{Text-based questions. Answer in complete sentences.}

\medskip
\noindent
\textbf{Question 1 (Skimming — Main idea):} What is the main purpose of this text?
\tcblower
\textbf{Solution:} The main purpose is to illustrate the concrete \cle{language needs} of a technician (Mr. Nguema) at Douala Seaport and to show that English is an essential working tool for his daily tasks (reading manuals, communicating, writing reports, following training).

\medskip
\noindent
\textbf{Question 2 (Scanning — Specific detail):} Which error code did Mr. Nguema look up in the English manual last month?
\tcblower
\textbf{Solution:} He looked up error code \texttt{E-407} in the crane's English maintenance manual to troubleshoot a hydraulic fault.

\medskip
\noindent
\textbf{Question 3 (Vocabulary in context):} The word \emph{troubleshooting} (paragraph 2) means:
\begin{itemize}
    \item[a)] repairing a machine without a manual
    \item[b)] identifying and solving a technical problem \textbf{(correct)}
    \item[c)] writing a maintenance report
    \item[d)] operating a crane
\end{itemize}
\tcblower
\textbf{Solution:} \textbf{b)} The context says: \emph{« follow the troubleshooting steps »} after identifying the error code. \cle{Troubleshooting} = diagnostic et dépannage (identifier et résoudre un problème technique).

\medskip
\noindent
\textbf{Question 4 (Inference):} Why does Mr. Nguema say mistakes in communication can cause accidents? (Paragraph 3)
\tcblower
\textbf{Solution:} Because he gives technical information (e.g., maximum container weight) to agents and dockworkers. If the information is wrong or misunderstood, a container could be overloaded, fall, or injure someone. Clear English communication is a \cle{safety} requirement.

\medskip
\noindent
\textbf{Question 5 (Global understanding):} List the four main language needs of Mr. Nguema mentioned in the text.
\tcblower
\textbf{Solution:}
\begin{enumerate}
    \item Reading technical manuals (maintenance, error codes, safety instructions).
    \item Communicating orally with English-speaking clients and colleagues (shipping agents, dockworkers).
    \item Writing reports (daily maintenance logs, incident reports, digital checklists).
    \item Following online training courses (PLC programming, industrial safety, specialised technical courses).
\end{enumerate}
\end{exoresolu}

\courssub{Identifying your own language needs}

Every technical field has specific English requirements. Use the table below to start your personal analysis.

\begin{popfigure}
\begin{center}
\begin{tabularx}{\linewidth}{|l|X|}
\hline
\rowcolor{popblueL}
\textbf{Technical field (Filière)} & \textbf{Specific English needs (Besoins spécifiques en anglais)} \\
\hline
\cle{Génie civil / BTP} & Reading annotated plans and structural notes in English; understanding international standards (Eurocodes, ASTM); writing site reports; communicating with foreign engineers. \\
\hline
\cle{Électrotechnique / Maintenance industrielle} & Reading electrical schematics with English annotations; PLC manuals (Siemens, Allen-Bradley, Schneider); safety datasheets (Lockout/Tagout); error codes on variable frequency drives. \\
\hline
\cle{Informatique / Réseaux} & Programming languages (keywords, documentation, Stack Overflow); Cisco/Network configuration guides; cybersecurity alerts; API documentation; writing technical specs. \\
\hline
\cle{Comptabilité / Gestion} & Reading IFRS standards in English; writing financial reports for international groups; using English ERP software (SAP, Sage); communicating with foreign auditors. \\
\hline
\cle{Hôtellerie / Tourisme} & Welcoming English-speaking guests; handling reservations and complaints; reading international hygiene/safety standards (HACCP); writing menus and brochures. \\
\hline
\cle{Mécanique / Usinage} & Reading CNC machine manuals (Fanuc, Heidenhain); G-code comments; tooling catalogues (Sandvik, Kennametal); writing inspection reports. \\
\hline
\cle{Bois / Menuiserie industrielle} & Reading machine manuals (Weinig, Homag); understanding wood grading standards (NHLA); writing production sheets; communicating with export clients. \\
\hline
\cle{Agroalimentaire} & Reading HACCP/ISO 22000 documentation; labelling regulations (Codex Alimentarius); writing quality control reports; export documentation. \\
\hline
\end{tabularx}
\end{center}
\legende{Tableau : filières techniques du lycée technique vs besoins spécifiques en anglais. À compléter selon votre spécialité.}
\end{popfigure}

\begin{exemplebox}[Analyse personnelle — Exemple : Élève en Électrotechnique]
\textbf{Mon besoin prioritaire :} Lire les manuels de variateurs de fréquence (VFD) et les codes d'erreur des automates Siemens.
\textbf{Preuve :} En stage, j'ai bloqué 2 heures sur un code \texttt{F0001} parce que le manuel n'était qu'en anglais.
\textbf{Objectif SMART :} D'ici 3 mois, comprendre 80\% des codes d'erreur courants sans dictionnaire.
\textbf{Actions :} 1) Apprendre 10 termes techniques/semaine (liste du manuel). 2) Regarder 1 tutoriel YouTube en anglais/semaine (PLC programming). 3) Rejoindre le club d'anglais technique du lycée.
\end{exemplebox}

\courssub{Distribution of a technician's English needs — Pie chart}

The following chart summarises the typical distribution of English language needs for a technician working in an industrial or port environment in Cameroon, based on surveys of professional tasks.

\begin{popfigure}
\begin{tikzpicture}
% Pie chart data
\def\angle{0}
\def\radius{3.5}
\def\labels{
  Lire notices techniques/popblue/30,
  Communiquer avec clients/poporange/25,
  Rédiger rapports/popgreen/20,
  Suivre formations/poppurple/15,
  Autres (emails et logiciels)/poppink/10
}
\foreach \label/\color/\percent in \labels {
  \pgfmathsetmacro{\newangle}{\angle + \percent * 3.6}
  \draw[fill=\color, draw=white, thick] (0,0) -- (\angle:\radius) arc (\angle:\newangle:\radius) -- cycle;
  \pgfmathsetmacro{\midangle}{(\angle + \newangle) / 2}
  \pgfmathsetmacro{\labelradius}{\radius * 0.65}
  \node at (\midangle:\labelradius) {\percent\%};
  \pgfmathsetmacro{\legendangle}{\midangle}
  \xdef\angle{\newangle}
}
% Legend
\begin{scope}[shift={(-5.5,0)}]
  \foreach \label/\color/\percent [count=\i] in \labels {
    \fill[\color] (0,-\i*0.7) rectangle (0.5,-\i*0.7+0.4);
    \node[anchor=west] at (0.7,-\i*0.7+0.2) {\label\ (\percent\%)};
  }
\end{scope}
\end{tikzpicture}
\legende{Répartition typique des besoins en anglais d'un technicien camerounais (enquête terrain 2023, n=120 techniciens portuaires et industriels). Source : adaptation MINESEC/ONBuch+.}
\end{popfigure}

\courssub{Building a personal action plan: SMART goals}

Once you have identified your needs, turn them into a concrete plan using the \cle{SMART} framework.

\begin{definition}[SMART goals]
A \cle{SMART goal} is:
\begin{itemize}
    \item \textbf{S}pecific (précis) — What exactly will you improve?
    \item \textbf{M}easurable (mesurable) — How will you track progress?
    \item \textbf{A}chievable (atteignable) — Is it realistic with your time?
    \item \textbf{R}elevant (pertinent) — Does it match your professional need?
    \item \textbf{T}ime-bound (daté) — By when?
\end{itemize}
\end{definition}

\begin{methode}[Créer votre plan d'action personnel en 4 étapes]
\begin{enumerate}
    \item \textbf{Listez vos 3 besoins prioritaires} (ex. : lire manuels VFD, écrire rapports d'intervention, comprendre formations PLC).
    \item \textbf{Transformez chaque besoin en objectif SMART} (voir exemple ci-dessous).
    \item \textbf{Choisissez 2 ressources par objectif} (cours lycée, appli mobile, radio, club anglais, YouTube technique, tandem avec un camarade anglophone).
    \item \textbf{Fixez un rendez-vous mensuel avec vous-même} pour évaluer et ajuster.
\end{enumerate}
\end{methode}

\begin{exemplebox}[Plan d'action SMART — Exemple complet]
\begin{center}
\begin{tabularx}{\linewidth}{|l|X|X|X|}
\hline
\rowcolor{popblueL}
\textbf{Besoin identifié} & \textbf{Objectif SMART} & \textbf{Ressources} & \textbf{Échéance} \\
\hline
Lire manuels VFD & Comprendre 80\% des codes d'erreur courants sans dictionnaire d'ici le 30 avril. & 1) Liste vocabulaire manuel Schneider (10 mots/sem.) 2) Chaîne YouTube \emph{Schneider Electric Tutorials} (1 vidéo/sem.) & 30 avril \\
\hline
Rédiger rapports & Écrire un rapport d'intervention structuré (50 mots) sans fautes majeures en 20 min. & 1) Modèles de rapports du lycée 2) Appli \emph{Grammarly} (version gratuite) 3) Relecture par prof/coéquipier & 15 mai \\
\hline
Suivre formation PLC & Terminer le module 1 du cours \emph{Siemens S7-1200} (anglais) sur \emph{Siemens Learning Advantage}. & 1) Compte gratuit Siemens 2) Groupe d'étude WhatsApp (3 camarades) 3) Dico technique \emph{Oxford Picture Dictionary} & 30 juin \\
\hline
\end{tabularx}
\end{center}
\end{exemplebox}

\courssub{Available resources in Cameroon}

\begin{itemize}
    \item \textbf{Au lycée :} Cours d'anglais technique, club d'anglais, jumelage avec des élèves de la section anglophone, bibliothèque (manuels techniques en anglais).
    \item \textbf{Médias :} CRTV Radio (émissions en anglais), BBC Learning English (site/app), VOA Learning English.
    \item \textbf{Applications gratuites :} Duolingo (base), Busuu (conversation), Quizlet (vocabulaire technique — créez vos propres listes), Grammarly (écriture).
    \item \textbf{Plateformes techniques :} Siemens Learning Advantage (gratuit), Schneider Electric University, Cisco NetAcad (anglais technique), Coursera/edX (audit gratuit).
    \item \textbf{Pratique réelle :} Stages, visites d'entreprises, correspondance avec des étudiants techniques nigérians/ghanéens (WhatsApp/email), traduction de notices de machines du laboratoire.
\end{itemize}

\begin{aretenir}[Checklist de fin de section]
\begin{itemize}
    \item[$\square$] Je peux expliquer pourquoi l'anglais est indispensable dans mon futur métier au Cameroun.
    \item[$\square$] Je connais le vocabulaire clé : \emph{language needs, proficiency, fluent, bilingual, manual, report writing, troubleshooting}.
    \item[$\square$] J'applique \emph{skimming, scanning, déduction du sens} à un texte professionnel.
    \item[$\square$] J'ai identifié mes 3 besoins prioritaires en anglais pour ma filière.
    \item[$\square$] J'ai rédigé au moins 1 objectif SMART avec ressources et échéance.
\end{itemize}
\end{aretenir}

\courssub{Practice exercises}

\begin{exercice}[Vocabulary consolidation]
Complete the sentences with the correct word from the box. Two words are extra.
\begin{center}
\begin{tabular}{cccc}
proficiency & manual & troubleshooting & fluent \\
bilingual & report writing & technical terms & customer service
\end{tabular}
\end{center}
\begin{enumerate}
    \item The maintenance \_\_\_\_\_\_ for the crane is 300 pages long and only in English.
    \item She works in \_\_\_\_\_\_ at a hotel; she answers guests' questions in English and French.
    \item His English \_\_\_\_\_\_ is certified at level B2 by the British Council.
    \item The error code \texttt{F004} means « overvoltage » — check the \_\_\_\_\_\_ guide page 12.
    \item Cameroon is officially a \_\_\_\_\_\_ country: French and English have equal status.
\end{enumerate}
\end{exercice}

\begin{exercice}[Reading strategy identification]
For each situation, choose the best strategy: \textbf{Skimming}, \textbf{Scanning} or \textbf{Deducing from context}.
\begin{enumerate}
    \item You need to find the maximum load of a 40-foot container in a 50-page manual.
    \item You want to know if a 10-page article on PLC safety is relevant for your exam.
    \item You meet the word \emph{lockout} in a safety text: \emph{« Apply lockout/tagout before maintenance. Isolate energy sources. »} You guess the meaning.
    \item You look for the date of the next training session in an email from your supervisor.
    \item You read the first and last paragraph of a report to decide whether to read it all.
\end{enumerate}
\end{exercice}

\begin{exercice}[Personal needs analysis — Written production]
Write a short paragraph (80–100 words) describing:
\begin{itemize}
    \item Your technical field (specialty).
    \item Your top 2 English language needs (be specific: \emph{reading hydraulic schematics, writing incident reports, speaking with clients...}).
    \item One SMART goal for the next 6 weeks.
    \item Two resources you will use.
\end{itemize}
Use the vocabulary from this section. Exchange with a partner and give feedback.
\end{exercice}

\begin{corrige}[Exercice 1 — Vocabulary consolidation]
\begin{enumerate}
    \item manual (notice/manuel)
    \item customer service (service client)
    \item proficiency (maîtrise certifiée)
    \item troubleshooting (dépannage/diagnostic)
    \item bilingual (bilingue — pays à deux langues officielles)
\end{enumerate}
\end{corrige}

\begin{corrige}[Exercice 2 — Reading strategy identification]
\begin{enumerate}
    \item Scanning (recherche d'une information précise : charge max).
    \item Skimming (idée générale pour décider de lire ou non).
    \item Deducing from context (indices : \emph{isolate energy sources, before maintenance} $\rightarrow$ \emph{lockout} = consignation/verrouillage).
    \item Scanning (recherche d'une date précise).
    \item Skimming (lecture rapide des parties stratégiques pour l'idée générale).
\end{enumerate}
\end{corrige}

\courssub{Link to the writing section}

You now have:
\begin{itemize}
    \item A clear understanding of your professional \cle{language needs}.
    \item Vocabulary to describe them (\emph{proficiency, technical terms, manual, report writing, troubleshooting}).
    \item A personal SMART action plan.
\end{itemize}
The next step (Writing section) will guide you to \textbf{write a structured composition} on « Anticipating my second-language needs at work » — using the ideas, vocabulary and plan you have just built. You will organise your text in paragraphs, use linking words, and demonstrate your ability to write professionally about your own career preparation.

\coursec{Writing: composition on second-language needs at work}

\courssub{Transition from reading to writing}

In the previous section, we read texts about \emph{anticipating second-language needs at work} and learnt vocabulary such as \cle{technical documentation} (\emph{documentation technique}), \cle{workplace communication} (\emph{communication en milieu professionnel}), \cle{career advancement} (\emph{avancement de carrière}) and \cle{multilingual environment} (\emph{environnement multilingue}). Now we move from \textbf{reception} (reading, listening) to \textbf{production} (writing). Writing a composition allows you to organise your ideas, reuse the grammar and vocabulary of the unit, and prepare for the \textbf{Baccalauréat technique} writing paper.

\begin{methode}[How to write a composition for the Bac]
\textbf{Step 1 -- Read the topic twice.} Underline key words: \emph{importance}, \emph{English}, \emph{future career}, \emph{technician}. Identify the required length (200--250 words).\\
\textbf{Step 2 -- Make a rough plan (5 min).} Write 3--4 main ideas, one per paragraph. Choose logical connectors.\\
\textbf{Step 3 -- Write the draft (20 min).} Follow the structure: Introduction $\rightarrow$ Body (2--3 paragraphs) $\rightarrow$ Conclusion. Use the tenses and phrasal verbs studied.\\
\textbf{Step 4 -- Proofread (5 min).} Check verbs (tense, subject-verb agreement), spelling, punctuation, word count. Ensure each sentence expresses \textbf{one clear idea}.
\end{methode}

\courssub{Structure of a composition}

A good composition has three clearly separated parts. The diagram below summarises the role of each part.

\begin{popfigure}
\begin{tikzpicture}[
  box/.style={draw=popblue, thick, rounded corners, fill=popblueL, minimum width=3.2cm, minimum height=1.6cm, align=center, font=\small},
  arrow/.style={-Stealth, thick, popdark},
  labelbox/.style={font=\footnotesize, text width=3cm, align=center}
]
\node[box, align=center] (intro){Introduction\\ \textit{Hook + plan announcement}};
\node[box, right=2.5cm of intro, align=center] (body1){Body §1\\ \textit{Current needs}};
\node[box, right=2.5cm of body1, align=center] (body2){Body §2\\ \textit{Advantages at work}};
\node[box, right=2.5cm of body2, align=center] (body3){Body §3\\ \textit{Action plan}};
\node[box, right=2.5cm of body3, align=center] (conc){Conclusion\\ \textit{Summary + opening}};
\draw[arrow] (intro) -- (body1) node[midway, above, labelbox] {announces};
\draw[arrow] (body1) -- (body2) node[midway, above, labelbox] {develops};
\draw[arrow] (body2) -- (body3) node[midway, above, labelbox] {extends};
\draw[arrow] (body3) -- (conc) node[midway, above, labelbox] {leads to};
\node[labelbox, below=0.4cm of intro, align=center]{1. Hook the reader\\2. State the topic\\3. Announce the plan};
\node[labelbox, below=0.4cm of body1, align=center]{Topic sentence\\+ arguments\\+ examples};
\node[labelbox, below=0.4cm of body2, align=center]{Topic sentence\\+ arguments\\+ examples};
\node[labelbox, below=0.4cm of body3, align=center]{Topic sentence\\+ arguments\\+ examples};
\node[labelbox, below=0.4cm of conc, align=center]{1. Restate thesis\\2. Summarise main points\\3. Opening thought};
\end{tikzpicture}
\legende{Structure d'une composition type Bac : Introduction $\rightarrow$ 3 paragraphes de développement, Conclusion.}
\end{popfigure}


\newpage

\coursec{Activité d'intégration}

\courssub{Objectifs de l'activité}
Cette activité d'intégration vise à mobiliser l'ensemble des savoirs et savoir-faire de l'unité \emph{Family and Social Life} dans une tâche unique, complexe et authentique : la simulation d'un processus de recrutement réel au Cameroun. À l'issue de cette activité, l'élève doit être capable de :
\begin{itemize}
    \item \textbf{Comprendre} une offre d'emploi technique (lecture sélective / \emph{scanning}) et en extraire le vocabulaire clé (\cle{job description}, \cle{requirements}, \cle{soft skills}, \cle{technical skills}).
    \item \textbf{Produire} un CV structuré et une présentation orale fluide intégrant le \cle{present perfect continuous} pour décrire un parcours de formation en cours (\emph{e.g. I have been studying electrical engineering for three years}) et le vocabulaire des \cle{recreational activities} pour valoriser les \cle{soft skills}.
    \item \textbf{Interagir} en situation d'entretien : formuler et répondre à des questions directes, des \cle{embedded questions} (\emph{Could you tell me why...?}), utiliser des \cle{question tags} pour vérifier l'information (\emph{You have a BTS, don't you?}) et employer des \cle{phrasal verbs} naturels (\emph{look forward to, get along with, carry out}).
    \item \textbf{Rédiger} un paragraphe cohérent (\emph{My plan to improve my English for this job}) en mobilisant le \cle{future continuous} pour projeter des actions en cours dans le futur (\emph{I will be attending evening classes...}).
    \item \textbf{Évaluer} ses pairs via une grille critériée (grammaire, vocabulaire, aisance, contenu) et identifier ses propres axes de remédiation.
\end{itemize}

\courssub{Situation}
\textbf{Contexte :} L'entreprise fictive \textbf{CamTech Services Ltd}, située à \textbf{Douala-Bonanjo} (immeuble Sawa, 3\textsuperscript{e} étage), est une société de maintenance industrielle et de solutions énergétiques (solaire, groupes électrogènes, automatisme). Elle recrute \textbf{deux (02) Techniciens Stagiaires} (spécialités : Électrotechnique, Maintenance des Systèmes Automatisés ou Génie Mécanique) pour un stage professionnalisant de 6 mois, renouvelable, avec perspective d'embauche (CDI). L'entreprise travaille avec des partenaires internationaux (Chine, France, Inde) ; la maîtrise de l'anglais technique est un critère éliminatoire.

\textbf{Offre d'emploi (extrait affiché au tableau / distribué) :}

\begin{aretenir}[Vacancy Announcement -- CamTech Services Ltd -- Douala]
\textbf{Position:} 2 Trainee Technicians (Electrical / Mechanical / Automation)\\
\textbf{Contract:} 6-month Internship (Stipend: 80 000 FCFA/month + Transport allowance)\\
\textbf{Start Date:} January 15\textsuperscript{th}, 2027

\textbf{Key Responsibilities:}
\begin{itemize}
    \item Carry out preventive and corrective maintenance on industrial equipment.
    \item Assist senior engineers in solar panel installation and generator servicing.
    \item Read and interpret technical schematics and manuals (mostly in English).
    \item Fill in daily intervention reports in English.
    \item Comply strictly with HSE (Health, Safety, Environment) standards.
\end{itemize}
\textbf{Profile Required:}
\begin{itemize}
    \item BTS / HND / DUT in Electrical Engineering, Mechanical Engineering or Automation.
    \item \textbf{English:} Minimum B1 level (CEFR). Ability to understand technical documents and communicate orally.
    \item \textbf{Soft skills:} Team spirit, adaptability, punctuality, problem-solving mindset.
    \item \textbf{Hobbies/Interests:} Candidates practising team sports, DIY (\textit{bricolage}), coding or technical reading are preferred.
\end{itemize}
\textbf{Application:} Send CV + Cover Letter in English to hr@camtech-cm.com before December 20\textsuperscript{th}.\\
\textbf{Selection Process:} Written test (Technical English) + Oral Interview (Panel of 3).
\end{aretenir}

\textbf{Organisation de la classe :} La classe est divisée en \textbf{groupes de 4 élèves}. Dans chaque groupe : \textbf{2 Candidats} (A et B) et \textbf{2 Membres du Jury} (J1 : Responsable RH, J2 : Chef de Service Technique). Les rôles s'inversent au second tour pour que chaque élève passe l'entretien. Le professeur circule, observe, note et remplit la grille d'évaluation officielle.

\courssub{Consignes}
\emph{Les consignes sont données en anglais aux élèves. Les mots clés sont glossés en français.}

\begin{enumerate}
    \item \textbf{Reading \& Scanning (10 min):} Individually, read the vacancy announcement above. Highlight or list \textbf{10 key words/expressions} (technical vocabulary + expected skills). \textit{Glossaire : to highlight = surligner ; key requirements = exigences clés ; schematics = schémas techniques.}
    \item \textbf{Preparation -- CV \& Oral Pitch (20 min):} Candidates: write a \textbf{simplified CV} (half a page) on paper. Prepare a \textbf{2-minute oral presentation} (\emph{``Tell us about yourself''}). \textbf{Compulsory grammar constraints:}
    \begin{itemize}
        \item Use the \cle{Present Perfect Continuous} at least twice for your ongoing studies/projects (\emph{I have been learning... ; I have been working on...}).
        \item Mention \textbf{one recreational activity} connected with the job (\emph{team sport, DIY, coding club}) and explain what it brings you (\emph{teamwork, patience, precision}).
    \end{itemize}
    Jury members: prepare \textbf{6 questions} (3 per candidate): 2 direct questions, 2 \cle{embedded questions}, 2 questions with \cle{question tags}. Include at least \textbf{3 different phrasal verbs} in your questions (\emph{deal with, find out, look into, get along with, carry out}).
    \item \textbf{Roleplay -- Job Interview (15 min per candidate):} Real-life simulation. The panel welcomes the candidate. Procedure:
    \begin{enumerate}
        \item Welcome \& \emph{ice-breaker} (question tag: \emph{You are [Name], aren't you?}).
        \item Candidate's presentation (2 min, without notes if possible).
        \item Technical \& behavioural questions (Jury) -- Answers (Candidate). The candidate must use at least \textbf{2 phrasal verbs} in his/her answers.
        \item Candidate's questions to the panel (one embedded question compulsory: \emph{Could you tell me what a typical day looks like?}).
        \item Conclusion.
    \end{enumerate}
    \item \textbf{Writing Task -- Language Development Plan (15 min):} \textbf{All students} (candidates and jury) individually write a paragraph of \textbf{80--100 words}: \emph{``My Plan to Improve My English for This Job''}. \textbf{Constraint:} Use the \cle{Future Continuous} at least \textbf{3 times} to describe actions in progress in the near future (\emph{Next month, I will be attending... ; In February, I will be reading...}). Use the vocabulary of \emph{anticipating second-language needs} (\cle{technical vocabulary}, \cle{report writing}, \cle{oral communication}, \cle{online courses}).
    \item \textbf{Evaluation \& Remediation (15 min):} The Jury fills in the \textbf{evaluation grid} (out of 20) for each candidate. The teacher collects the grids and the written productions. \textbf{Collective remediation} session: on the board, correction of the 5 most frequent errors observed (grammar, pronunciation, lexical awkwardness).
\end{enumerate}

\courssub{Ressources à mobiliser}
\begin{center}
\begin{tabularx}{\linewidth}{|l|X|}\hline
\rowcolor{popblueL}
\textbf{Grammatical Tool} & \textbf{Formula / Key Rule (Terminale Tech Reminder)} \\ \hline
\cle{Present Perfect Continuous} & \textbf{Have/Has + been + V-ing}. \emph{Action started in the past, still in progress, or just stopped with a visible result.} \textit{Ex: I have been studying electricity for 3 years. / My hands are dirty; I have been repairing the generator.} \\ \hline
\cle{Future Continuous} & \textbf{Will + be + V-ing}. \emph{Action that will be in progress at a precise moment in the future; a plan or arrangement.} \textit{Ex: At 8 pm tonight, I will be revising my technical vocabulary. / Next semester, I will be taking an English course.} \\ \hline
\cle{Embedded Questions} & \textbf{Introductory phrase + question word + subject + verb (affirmative order, NO auxiliary do/does/did).} \textit{Direct: Why did you choose this field? $\rightarrow$ Embedded: Could you tell me \textbf{why you chose} this field?} \textit{Direct: Does the company provide training? $\rightarrow$ Embedded: I'd like to know \textbf{if/whether the company provides} training.} \\ \hline
\cle{Question Tags} & \textbf{Affirmative statement $\rightarrow$ negative tag / negative statement $\rightarrow$ affirmative tag.} \textit{Verb tense and subject pronoun must match. Ex: You have a BTS, \textbf{don't you?} / You aren't leaving, \textbf{are you?} / He can speak English, \textbf{can't he?}} \\ \hline
\cle{Phrasal Verbs (Selection)} & \textbf{Carry out}: effectuer, réaliser (une tâche, une maintenance). \textbf{Deal with}: gérer, traiter (un problème, un client). \textbf{Find out}: découvrir, s'informer. \textbf{Look into}: examiner, enquêter sur. \textbf{Get along with}: s'entendre bien avec (collègues). \textbf{Look forward to + V-ing}: avoir hâte de. \textbf{Set up}: installer, configurer. \textbf{Run out of}: manquer de (pièces, temps). \\ \hline
\cle{Interview Vocabulary} & \cle{Strengths / Weaknesses} (forces/faiblesses), \cle{Background} (parcours), \cle{Challenge} (défi), \cle{Deadline} (échéance), \cle{Shift work} (travail par roulement), \cle{Internship} (stage), \cle{Stipend} (indemnité de stage). \\ \hline
\cle{Leisure Vocabulary} & \cle{Team sports} (football, basketball), \cle{DIY / Do It Yourself} (bricolage), \cle{coding/programming}, \cle{reading technical magazines}, \cle{volunteering} (bénévolat). \\ \hline
\end{tabularx}
\end{center}

\courssub{Démarche et corrigé commenté}

\begin{exoresolu}{Simulation complète -- Candidat modèle : Ngono Jean-Pierre (Électrotechnique)}
\textbf{Étape 1 : Reading \& Scanning -- Vocabulaire clé extrait de l'offre}
\begin{itemize}
    \item \cle{Preventive / corrective maintenance} (maintenance préventive / corrective)
    \item \cle{Technical schematics} (schémas techniques)
    \item \cle{HSE standards} (normes Hygiène Sécurité Environnement)
    \item \cle{Intervention reports} (rapports d'intervention)
    \item \cle{Team spirit / teamwork} (esprit d'équipe)
    \item \cle{Problem-solving mindset} (esprit de résolution de problèmes)
    \item \cle{DIY (Do It Yourself)} (bricolage)
    \item \cle{Stipend} (indemnité de stage)
    \item \cle{Manuals} (manuels techniques)
    \item \cle{Deadline} (date limite : December 20\textsuperscript{th})
\end{itemize}
\textit{Commentaire :} L'élève prouve sa capacité de \emph{scanning} (lecture rapide ciblée). Ces mots seront réinvestis à l'oral et à l'écrit.

\textbf{Étape 2 : CV simplifié -- modèle}

\begin{center}
\begin{tabularx}{\linewidth}{|l|X|}\hline
\rowcolor{popblueL}
\textbf{Section} & \textbf{Content (Model)} \\ \hline
\textbf{Name} & Ngono Jean-Pierre \\ \hline
\textbf{Contact} & +237 6XX XX XX XX | jp.ngono@email.com | Douala, Bonabéri \\ \hline
\textbf{Diploma} & BTS Electrotechnics (Government Technical High School Douala) -- \textit{final year, defence in June 2027} \\ \hline
\textbf{Technical Skills} & Maintenance of asynchronous motors, cabinet wiring, reading AutoCAD drawings, basics of PLC programming (Siemens LOGO!), electrical safety \\ \hline
\textbf{English Level} & B1 (CEFR) -- \textit{2-month practical internship at ENEO (English-speaking customer service)} \\ \hline
\textbf{Recreational Activities} & \textbf{Captain of the school basketball team} (3 years) -- \textit{leadership, discipline, stress management.} \textbf{DIY electronics} -- \textit{repairing old radio sets, assembling Arduino kits (self-taught).} \\ \hline
\textbf{Soft Skills} & Punctuality, adaptability, teamwork, problem-solving \\ \hline
\end{tabularx}
\end{center}

\textbf{Script oral (modèle) -- présentation de 2 minutes :}

\emph{``Good morning. My name is Jean-Pierre Ngono. I am 20 years old and I come from Douala. I \textbf{have been studying} Electrotechnics at the Technical High School of Douala for three years now. Currently, I \textbf{have been working} on my final project: the design of an automatic transfer switch for a solar backup system. This project has allowed me to improve my skills in reading technical schematics and PLC programming.}

\emph{Besides my studies, I am the captain of the school basketball team. I \textbf{have been playing} basketball for five years. This recreational activity taught me teamwork, discipline and how to handle pressure during tight matches -- skills I know are essential for shift work and meeting deadlines at CamTech. I also enjoy DIY electronics; recently I repaired an old amplifier for a friend.}

\emph{I am very motivated to join CamTech Services because I want to \textbf{carry out} real industrial maintenance and \textbf{deal with} English technical documentation daily. I \textbf{look forward to} learning from your senior engineers. Thank you.''}

\textit{Analyse grammaticale :}
\begin{itemize}
    \item \cle{Present Perfect Continuous} : \emph{I \textbf{have been studying}... for three years} (parcours en cours) ; \emph{I \textbf{have been working} on my final project} (action récente en cours avec résultat visible).
    \item \cle{Recreational activity} : \emph{I \textbf{have been playing} basketball for five years} (durée).
    \item \cle{Phrasal verbs} : \emph{carry out, deal with, look forward to}.
    \item Vocabulaire technique : \emph{schematics, PLC, automatic transfer switch, shift work, deadlines}.
\end{itemize}

\textbf{Étape 3 : Roleplay -- extrait d'entretien (Jury J1 \& J2 vs Candidat)}

\begin{center}
\begin{tabularx}{\linewidth}{|l|X|}\hline
\rowcolor{popblueL}
\textbf{Speaker} & \textbf{Exchange (grammatical annotations)} \\ \hline
\textbf{J1 (HR)} & \emph{Good morning, Jean-Pierre. You are Mr. Ngono, \textbf{aren't you?} (\cle{Question tag} -- identity check). Please, have a seat.} \\ \hline
\textbf{Candidate} & \emph{Good morning. Yes, I am. Thank you.} \\ \hline
\textbf{J2 (Tech)} & \emph{We read your CV. \textbf{Could you tell us why you chose} electrotechnics? (\cle{Embedded question} -- Wh- question, subject-verb order).} \\ \hline
\textbf{Candidate} & \emph{I chose this field because I have always been curious about how energy is distributed. Since childhood, I liked to \textbf{find out} how appliances worked (\cle{phrasal verb}). At school, I \textbf{have been focusing} on power systems (\cle{Present Perfect Continuous}).} \\ \hline
\textbf{J1} & \emph{The job involves reading manuals in English. \textbf{Do you think you can handle} technical documentation? (\cle{Direct question}).} \\ \hline
\textbf{Candidate} & \emph{Yes, I can. During my internship at ENEO, I \textbf{had to fill in} intervention reports in English. I \textbf{have been improving} my technical vocabulary by reading engineering articles online (\cle{Present Perfect Continuous}). Next month, I \textbf{will be attending} a specialised ``Technical English'' evening course (\cle{Future Continuous} -- anticipation).} \\ \hline
\textbf{J2} & \emph{We often work in teams on site. \textbf{You get along well with} your classmates, \textbf{don't you?} (\cle{Question tag} + \cle{phrasal verb} \emph{get along with}).} \\ \hline
\textbf{Candidate} & \emph{Yes, I do. As team captain, I \textbf{have to deal with} different personalities (\cle{phrasal verb}). We usually \textbf{sort out} disagreements quickly (\cle{phrasal verb}).} \\ \hline
\textbf{J1} & \emph{What about safety? \textbf{I'd like to know if you are familiar with} HSE standards. (\cle{Embedded question} -- Yes/No $\rightarrow$ \emph{if/whether}).} \\ \hline
\textbf{Candidate} & \emph{I have my electrical safety certificate. I \textbf{have been following} strict safety protocols in the workshop for two years (\cle{Present Perfect Continuous}). Safety is not negotiable.} \\ \hline
\textbf{Candidate (final question)} & \emph{Thank you. \textbf{Could you tell me what a typical day} looks like for a trainee here? (\cle{Embedded question} -- compulsory instruction).} \\ \hline
\textbf{J2} & \emph{You will shadow a senior technician. Mornings: preventive maintenance. Afternoons: corrective interventions or reporting. You will \textbf{look into} fault reports (\cle{phrasal verb}).} \\ \hline
\end{tabularx}
\end{center}

\textbf{Étape 4 : Writing Task -- production écrite attendue (modèle, environ 95 mots)}

\emph{``My Plan to Improve My English for This Job}

\emph{Getting the trainee position at CamTech Services requires a solid command of Technical English. Therefore, I have designed a precise schedule. Starting next Monday, I \textbf{will be attending} a Business English course at a language centre in Douala every Tuesday and Thursday evening. Throughout the next three months, I \textbf{will be reading} one technical article from an electrical engineering website daily to build my specialised vocabulary. During my commute, I \textbf{will be listening to} engineering podcasts to improve my oral comprehension. Additionally, I \textbf{will be writing} a weekly summary of my practical work in English to practise report writing. I am confident this routine will bridge the gap before the final interview.''}

\textit{Analyse du corrigé :}
\begin{itemize}
    \item \cle{Future Continuous} (x4) : \emph{will be attending, will be reading, will be listening, will be writing} -- actions en cours dans un futur planifié.
    \item Vocabulaire \emph{second-language needs} : \emph{Technical English, specialised vocabulary, oral comprehension, report writing, bridge the gap}.
    \item Structure : objectif $\rightarrow$ actions planifiées (Future Continuous) $\rightarrow$ résultat attendu.
    \item Connecteurs logiques : \emph{Therefore, Additionally, Throughout}.
\end{itemize}
\tcblower
\textbf{Solution détaillée et commentaire pédagogique :} L'activité valide l'intégration des six blocs de compétences du programme (speaking, listening, reading, writing, grammar, vocabulary). Le candidat modèle démontre une progression logique : \emph{scanning} (input) $\rightarrow$ \emph{production orale/écrite} (output) $\rightarrow$ \emph{interaction} (médiation). La grille d'évaluation critériée objectivise la notation, conformément aux attentes du Probatoire et du Baccalauréat technique. La remédiation ciblée sur les cinq erreurs types (question tags, ordre des mots dans les embedded questions, Present Perfect Continuous vs Present Perfect Simple, Future Continuous vs Future Simple, séparabilité des phrasal verbs) couvre l'essentiel des pertes de points aux examens officiels.
\end{exoresolu}

\textbf{Étape 5 : Grille d'évaluation (remplie par le Jury pour J.-P. Ngono)}

\begin{center}
\begin{tabularx}{\linewidth}{|X|c|X|}\hline
\rowcolor{popblueL}
\textbf{Critère} & \textbf{Note /5} & \textbf{Commentaire du Jury} \\ \hline
\textbf{Grammaire} (tenses, embedded questions, tags, phrasal verbs) & 4,5 & Maîtrise excellente du Present Perfect Continuous et du Future Continuous. 3 phrasal verbs utilisés naturellement. 1 tag mal formé (\emph{``I am, am I not?''} au lieu de \emph{``aren't I?''} -- erreur classique corrigée en remédiation). \\ \hline
\textbf{Vocabulaire} (technique, loisirs, entretien) & 5 & Vocabulaire précis (\emph{schematics, HSE, transfer switch, shift work}). Loisirs bien exploités (\emph{captain, teamwork}). \\ \hline
\textbf{Aisance / Prononciation} & 4 & Bon débit, légère hésitation sur \emph{``schematics''}. Contact visuel maintenu. \\ \hline
\textbf{Contenu / Pertinence} & 5 & Réponses adaptées au poste (maintenance, sécurité, anglais technique). Question finale pertinente. \\ \hline
\textbf{TOTAL} & \textbf{18,5 / 20} & \textbf{Admis -- profil retenu.} \\ \hline
\end{tabularx}
\end{center}

\textbf{Étape 6 : Remédiation collective (exemples corrigés au tableau)}

\begin{attention}[Erreurs fréquentes relevées -- Correction collective]
\begin{enumerate}
    \item \textbf{Present Perfect Continuous vs Present Perfect Simple} : \emph{*I have studied here for 3 years} (accent sur le résultat) vs \emph{I have been studying here for 3 years} (accent sur la durée, action non terminée). \textbf{Règle :} durée + action non terminée $\rightarrow$ forme continue.
    \item \textbf{Question tags -- cas particulier ``I am''} : \emph{I am late, \textbf{aren't I?}} (jamais *amn't I). \emph{I am not late, \textbf{am I?}}.
    \item \textbf{Embedded questions -- ordre des mots} : \emph{*Could you tell me why did you choose?} $\rightarrow$ \emph{Could you tell me \textbf{why you chose}?} (pas d'auxiliaire \emph{did}, ordre sujet-verbe).
    \item \textbf{Future Continuous vs Future Simple} : \emph{I will attend classes tomorrow} (décision ou prédiction simple) vs \emph{I will be attending classes at 6 pm tomorrow} (action en cours à un moment précis). Le plan d'action exige \emph{will be + V-ing}.
    \item \textbf{Phrasal verbs -- séparabilité} : \emph{look into the matter} (inséparable) vs \emph{carry out the task / carry the task out} (séparable ; obligatoire avec un pronom : \emph{carry it out}).
\end{enumerate}
\end{attention}

\courssub{Bilan de l'activité}
Cette activité d'intégration \textbf{Job Interview Simulation \& Career Language Plan} a permis de :
\begin{itemize}
    \item \textbf{Contextualiser} les apprentissages : l'anglais n'est plus une matière scolaire mais un outil professionnel indispensable pour un technicien camerounais visant le secteur formel (entreprises de maintenance, d'énergie, de télécommunications, partenaires internationaux).
    \item \textbf{Valider la progression grammaticale} : le \cle{Present Perfect Continuous} ancre le parcours \emph{antérieur et présent} ; le \cle{Future Continuous} projette le \emph{développement professionnel continu} ; les \cle{embedded questions} et \cle{question tags} garantissent la \emph{politesse professionnelle} et la \emph{vérification de l'information} ; les \cle{phrasal verbs} assurent la \emph{fluidité idiomatique}.
    \item \textbf{Développer l'autonomie} : l'élève gère son temps (préparation), s'auto-évalue (grille), corrige ses pairs (remédiation).
    \item \textbf{Préparer l'examen} : le format (Reading + Speaking + Writing + Grammar in context) reproduit fidèlement les exigences des épreuves d'anglais du \textbf{Probatoire} et du \textbf{Baccalauréat technique} (compréhension de l'écrit, production écrite guidée, grammaire en contexte et interaction orale).
\end{itemize}

\newpage

\coursec{Méthodes et exercices résolus}

\courssub{Méthodes et exercices résolus}

\begin{methode}[Réussir un entretien d'embauche -- Speaking \& Vocabulary]
\textbf{Objectif :} Se présenter, répondre aux questions classiques et poser des questions pertinentes lors d'un entretien professionnel.
\begin{enumerate}
    \item \textbf{Préparer son \cle{pitch} de présentation (1 min) :} Nom, formation technique, spécialité, stage/expérience clé, qualité principale (ex : \emph{rigorous, adaptable, team player}).
    \item \textbf{Maîtriser le vocabulaire clé :} \cle{apply for a post} (postuler), \cle{CV/résumé}, \cle{cover letter} (lettre de motivation), \cle{qualifications} (diplômes et compétences), \cle{work experience} (expérience professionnelle), \cle{strengths/weaknesses} (forces/faiblesses), \cle{salary expectations} (prétentions salariales), \cle{notice period} (préavis).
    \item \textbf{Structurer les réponses avec la méthode STAR :} \cle{Situation}, \cle{Task}, \cle{Action}, \cle{Result}. Ex : \emph{« During my internship (S), I had to repair a hydraulic press (T). I diagnosed the fault and replaced the seals (A). The machine restarted and production losses stopped (R). »}
    \item \textbf{Gérer les questions pièges :} Pour les faiblesses, citer un défaut réel mais en cours d'amélioration (ex : \emph{« I sometimes focus too much on details, so I now use checklists to manage time better. »}).
    \item \textbf{Poser 2 questions finales :} Sur la formation continue, l'équipe, les projets techniques. Éviter la question du salaire au premier entretien.
\end{enumerate}
\end{methode}

\begin{exoresolu}[Simulation d'entretien pour un poste de technicien maintenance]
\textbf{Énoncé :} Vous êtes candidat au poste de \emph{Maintenance Technician} dans une entreprise agroalimentaire. Le recruteur vous demande : \emph{« Tell me about a technical problem you solved recently. »} Répondez en utilisant la méthode STAR et le vocabulaire technique approprié. Ensuite, formulez une question à poser au recruteur sur la maintenance préventive.

\tcblower
\textbf{Solution détaillée :}

\textbf{Étape 1 : Analyse de la consigne.} Le recruteur teste la compétence technique (résolution de problème) et la communication structurée (STAR). Le vocabulaire attendu : \cle{troubleshooting} (dépannage), \cle{downtime} (temps d'arrêt), \cle{preventive maintenance} (maintenance préventive), \cle{root cause} (cause profonde), \cle{spare parts} (pièces de rechange).

\textbf{Étape 2 : Construction de la réponse (STAR).}
\begin{itemize}
    \item \textbf{Situation :} \emph{« During my industrial attachment at a bottling plant in Douala, the conveyor belt on Line 3 stopped unexpectedly during peak production. »}
    \item \textbf{Task :} \emph{« My task was to identify the fault and restart the line within 30 minutes to minimize \cle{downtime}. »}
    \item \textbf{Action :} \emph{« I performed a \cle{visual inspection} and checked the \cle{PLC} error codes. I detected a \cle{misaligned photoelectric sensor}. I realigned the sensor, tightened the mounting bracket, and ran a \cle{test cycle}. »}
    \item \textbf{Result :} \emph{« The line restarted in 20 minutes. I then proposed adding this sensor check to the \cle{daily preventive maintenance checklist} to prevent recurrence. »}
\end{itemize}

\textbf{Étape 3 : Formulation de la question au recruteur.}
\emph{« Could you tell me more about your \cle{preventive maintenance schedule}? Is it based on the \cle{manufacturer's recommendations}, or do you also use \cle{predictive maintenance} tools like vibration analysis? »}

\textbf{Conclusion :} La réponse démontre une compétence technique concrète (capteur, PLC, checklist), une structure claire (STAR) et un vocabulaire précis. La question finale montre un intérêt professionnel pour l'organisation de la maintenance.
\end{exoresolu}

\begin{methode}[Décrire des activités de loisirs -- Listening \& Vocabulary]
\textbf{Objectif :} Comprendre un texte oral sur les loisirs et en restituer les idées principales avec le vocabulaire spécifique.
\begin{enumerate}
    \item \textbf{Identifier le type d'activité :} \cle{indoor/outdoor} (en intérieur/en plein air), \cle{individual/team} (individuelle/collective), \cle{physical/intellectual} (physique/intellectuelle).
    \item \textbf{Repérer les marqueurs chronologiques :} \cle{first, then, afterwards, finally} ; les expressions de fréquence : \cle{once a week, every weekend, occasionally}.
    \item \textbf{Noter le vocabulaire des sentiments/bénéfices :} \cle{unwind/relax} (se détendre), \cle{recharge one's batteries} (recharger ses batteries), \cle{relieve stress} (évacuer le stress), \cle{socialize} (socialiser), \cle{keep fit} (rester en forme), \cle{boost creativity} (stimuler la créativité).
    \item \textbf{Structurer le compte-rendu :} Introduction (type, fréquence) $\rightarrow$ Déroulement (actions) $\rightarrow$ Ressenti/Bénéfices $\rightarrow$ Opinion personnelle.
    \item \textbf{Utiliser les temps appropriés :} \cle{Present simple} (habitudes), \cle{Present continuous} (action en cours au moment où l'on parle), \cle{Past simple} (récit d'un événement passé).
\end{enumerate}
\end{methode}

\begin{exoresolu}[Compte-rendu d'écoute : « My Weekend Hiking Routine »]
\textbf{Énoncé :} Vous écoutez un technicien décrire son week-end type. Extrait : \emph{« I usually wake up early on Saturdays. I pack my backpack with water, a map and a first-aid kit. I drive to Mount Fako. The hike takes three hours. It is exhausting but the view from the summit is breathtaking. I feel refreshed for the week ahead. »}
Relevez 5 mots/expressions de vocabulaire « loisirs/équipement », conjuguez les verbes entre parenthèses au temps qui convient, et résumez le texte en 3 phrases.

\tcblower
\textbf{Solution détaillée :}

\textbf{1. Vocabulary spotting (5 items) :}
\begin{itemize}
    \item \cle{backpack} (sac à dos)
    \item \cle{first-aid kit} (trousse de secours)
    \item \cle{hike} (randonnée -- nom/verbe)
    \item \cle{summit} (sommet)
    \item \cle{breathtaking} (à couper le souffle)
    \item \cle{refreshed} (revigoré)
\end{itemize}

\textbf{2. Grammar -- Verb Tenses (contexte : routine).}
Le texte décrit une \emph{habitude} (\emph{usually}) $\rightarrow$ \textbf{Present Simple}.
\begin{itemize}
    \item I usually \textbf{wake up} (wake up) early.
    \item I \textbf{pack} (pack) my backpack...
    \item I \textbf{drive} (drive) to Mount Fako.
    \item The hike \textbf{takes} (take) three hours.
    \item It \textbf{is} (be) exhausting...
    \item I \textbf{feel} (feel) refreshed...
\end{itemize}
\emph{Note :} Pas de \emph{Present Continuous} ici car l'action ne se déroule pas au moment de la parole : c'est une routine générale.

\textbf{3. Summary (3 sentences) :}
\emph{« The speaker usually hikes on Mount Fako on Saturdays. He prepares his equipment carefully, including a first-aid kit. Although the three-hour hike is exhausting, he finds the summit view breathtaking and feels refreshed for the week. »}

\textbf{Conclusion :} L'écoute active permet d'extraire le vocabulaire (équipement, ressenti) et de choisir le bon temps (Present Simple pour l'habitude). Le résumé synthétise : Qui/Quoi/Où $\rightarrow$ Comment $\rightarrow$ Résultat/Bénéfice.
\end{exoresolu}

\begin{methode}[Maîtriser le Present Perfect Continuous \& le Future Continuous]
\textbf{Objectif :} Distinguer et utiliser correctement ces deux temps pour exprimer la durée d'une action inachevée (PPC) et une action en cours à un moment futur (FC).
\begin{enumerate}
    \item \textbf{Present Perfect Continuous (\cle{have/has been + V-ing}) :}
    \begin{itemize}
        \item \textbf{Forme :} Sujet + have/has + been + base verbale + -ing.
        \item \textbf{Usage :} Action commencée dans le passé, qui se poursuit \textbf{maintenant} (ou vient de s'arrêter) + insistance sur la \textbf{durée}. Mots-clés : \cle{for} (durée), \cle{since} (point de départ), \cle{lately/recently}, \cle{how long}.
        \item \textbf{Exemple :} \emph{« I \textbf{have been studying} English \textbf{for} three years. »} (J'étudie depuis 3 ans et je continue).
    \end{itemize}
    \item \textbf{Future Continuous (\cle{will be + V-ing}) :}
    \begin{itemize}
        \item \textbf{Forme :} Sujet + will + be + base verbale + -ing.
        \item \textbf{Usage :} Action qui \textbf{sera en train de se passer} à un moment précis dans le futur. Mots-clés : \cle{at this time tomorrow}, \cle{at 8 pm tonight}, \cle{when you arrive}.
        \item \textbf{Exemple :} \emph{« At 10 am tomorrow, I \textbf{will be sitting} for my exam. »} (Demain à 10h, je serai en train de passer mon examen).
    \end{itemize}
    \item \textbf{Piège à éviter :} Verbes d'état (\cle{know, like, want, believe, belong}) $\rightarrow$ pas de forme en \emph{-ing}. Utiliser le \emph{Present Perfect Simple} : \emph{« I have known him for years »} (jamais \emph{been knowing}).
\end{enumerate}
\end{methode}

\begin{exoresolu}[Choix du temps et transformation]
\textbf{Énoncé :}
A. Complétez avec le \textbf{Present Perfect Continuous} ou le \textbf{Present Perfect Simple} (verbes entre parenthèses).
1. My hands are dirty. I \_\_\_\_\_\_\_\_\_\_ (repair) the generator.
2. She \_\_\_\_\_\_\_\_\_\_ (write) three reports since this morning.
3. How long \_\_\_\_\_\_\_\_\_\_ you \_\_\_\_\_\_\_\_\_\_ (wait) for the technician?

B. Complétez avec le \textbf{Future Continuous}.
4. Don't call at 7 am. I \_\_\_\_\_\_\_\_\_\_ (commute) to the plant.
5. This time next week, we \_\_\_\_\_\_\_\_\_\_ (present) our project.

C. Traduisez en anglais : « Depuis quand est-ce que tu apprends l'anglais ? » et « Demain à 14h, je serai en train de souder la tuyauterie. »

\tcblower
\textbf{Solution détaillée :}

\textbf{Partie A : Present Perfect Continuous vs Simple}
\begin{enumerate}
    \item \textbf{I \emph{have been repairing} the generator.}
    \emph{Justification :} Preuve visible dans le présent (\emph{My hands are dirty}) $\rightarrow$ insistance sur l'activité qui vient de se dérouler. Verbe d'action dynamique.
    \item \textbf{She \emph{has written} three reports since this morning.}
    \emph{Justification :} Résultat quantifié (\emph{three reports}) $\rightarrow$ \textbf{Present Perfect Simple}. Le PPC (\emph{has been writing}) insisterait sur le processus, pas sur le nombre achevé.
    \item \textbf{How long \emph{have} you \emph{been waiting} for the technician?}
    \emph{Justification :} Question sur la durée (\emph{How long}) + action encore en cours (attente) $\rightarrow$ \textbf{PPC} standard.
\end{enumerate}

\textbf{Partie B : Future Continuous}
\begin{enumerate}
    \setcounter{enumi}{3}
    \item \textbf{Don't call at 7 am. I \emph{will be commuting} to the plant.}
    \emph{Justification :} Action en cours à un moment futur précis (\emph{at 7 am}).
    \item \textbf{This time next week, we \emph{will be presenting} our project.}
    \emph{Justification :} Moment futur précis (\emph{this time next week}) + action longue (présentation) en cours à ce moment-là.
\end{enumerate}

\textbf{Partie C : Traduction}
\begin{enumerate}
    \item \emph{« How long \textbf{have} you \textbf{been learning} English? »} (PPC : durée + action continue).
    \item \emph{« Tomorrow at 2 pm, I \textbf{will be welding} the piping. »} (FC : moment futur précis + action technique en cours).
\end{enumerate}

\textbf{Conclusion :} La distinction repose sur \textbf{Résultat/Quantité (Perfect Simple)} vs \textbf{Processus/Durée (Perfect Continuous)} pour le passé/présent, et sur \textbf{l'action en cours à un instant T futur} pour le Future Continuous.
\end{exoresolu}

\begin{methode}[Construire des questions indirectes, question tags et phrasal verbs]
\textbf{Objectif :} Transformer une question directe en question indirecte (embedded), former les question tags corrects, et utiliser des verbes à particules courants en contexte professionnel.
\begin{enumerate}
    \item \textbf{Embedded Questions (questions indirectes) :}
    \begin{itemize}
        \item Règle : \textbf{ordre affirmatif} (Sujet + Verbe), \textbf{PAS d'auxiliaire DO/DID}, pas de point d'interrogation si la phrase principale est affirmative.
        \item Introducteurs : \cle{Can you tell me...}, \cle{I wonder...}, \cle{Do you know...}, \cle{Please explain...}.
        \item Si question fermée (Yes/No) $\rightarrow$ ajouter \cle{if} ou \cle{whether}.
        \item Ex : \emph{Direct : « Where is the lab? »} $\rightarrow$ \emph{Indirect : « Can you tell me \textbf{where the lab is}? »}
    \end{itemize}
    \item \textbf{Question Tags :}
    \begin{itemize}
        \item Règle : \textbf{Positif $\leftrightarrow$ Négatif}. Même auxiliaire/modal que la phrase principale. Sujet = pronom personnel.
        \item \emph{« You are ready, \textbf{aren't you}? »} / \emph{« He won't come, \textbf{will he}? »}
        \item Cas particuliers : \emph{I am $\rightarrow$ aren't I?} / \emph{Let's $\rightarrow$ shall we?} / \emph{Impératif $\rightarrow$ will you? / won't you?}
        \item Intonation : montante = vraie question (incertitude) ; descendante = confirmation attendue.
    \end{itemize}
    \item \textbf{Phrasal Verbs \& Prepositional Phrases (contexte technique/professionnel) :}
    \begin{center}
    \begin{tabularx}{\linewidth}{|l|X|X|}\hline
    \rowcolor{popblueL}\textbf{Phrasal Verb} & \textbf{Sens (FR)} & \textbf{Exemple en contexte technique} \\ \hline
    \cle{carry out} & effectuer, réaliser & \emph{We must \textbf{carry out} the safety audit.} \\ \hline
    \cle{find out} & découvrir, s'informer & \emph{Please \textbf{find out} the specs of the new motor.} \\ \hline
    \cle{look into} & examiner, enquêter sur & \emph{The engineer will \textbf{look into} the vibration issue.} \\ \hline
    \cle{set up} & installer, configurer & \emph{We need to \textbf{set up} the new CNC machine.} \\ \hline
    \cle{shut down} & arrêter (machine/système) & \emph{Shut down the system before maintenance.} \\ \hline
    \cle{run out of} & manquer de, épuiser & \emph{We \textbf{ran out of} welding electrodes.} \\ \hline
    \end{tabularx}
    \end{center}
    \emph{Règle de séparation :} l'objet pronom se place \textbf{entre} le verbe et la particule (\emph{shut \textbf{it} down}). L'objet nom peut se placer après ou entre les deux.
\end{enumerate}
\end{methode}

\begin{exoresolu}[Transformation et mise en situation]
\textbf{Énoncé :}
1. Transformez en question indirecte : \emph{« What time does the shift start? »} (Début : \emph{Could you tell me...})
2. Ajoutez le question tag : \emph{The technician hasn't arrived yet, \_\_\_\_\_\_\_\_\_\_?}
3. Remplacez le verbe souligné par un phrasal verb de la liste : \emph{« We must \textbf{investigate} the cause of the breakdown. »}
4. Complétez : \emph{« Hand \_\_\_\_\_\_\_\_\_\_ your safety helmet at the gate. »} (phrasal verb : \cle{hand in} = remettre).
5. Corrigez l'erreur : \emph{« I am late, amn't I? »}

\tcblower
\textbf{Solution détaillée :}

\textbf{1. Embedded Question :}
\emph{« Could you tell me \textbf{what time the shift starts}? »}
\emph{Analyse :} Question ouverte (\emph{What time}) $\rightarrow$ on garde le mot interrogatif. Ordre affirmatif : \emph{the shift starts} (et non \emph{does the shift start}). L'auxiliaire \emph{does} disparaît et la marque de la 3\textsuperscript{e} personne revient sur le verbe (\emph{starts}).

\textbf{2. Question Tag :}
\emph{« The technician hasn't arrived yet, \textbf{has he}? »}
\emph{Analyse :} Phrase négative (\emph{hasn't}) $\rightarrow$ tag positif (\emph{has}). Sujet \emph{The technician} $\rightarrow$ pronom \emph{he}. Auxiliaire \emph{has} (Present Perfect).

\textbf{3. Phrasal Verb (substitution) :}
\emph{« We must \textbf{look into} the cause of the breakdown. »}
\emph{Analyse :} \emph{Investigate} = \emph{look into} (mener une enquête technique approfondie). \emph{Find out} conviendrait pour découvrir une information, moins pour analyser une panne.

\textbf{4. Phrasal Verb (complétion) :}
\emph{« Hand \textbf{in} your safety helmet at the gate. »}
\emph{Analyse :} \emph{Hand in} = remettre, rendre (documents, équipement). L'objet nom \emph{your safety helmet} peut aussi précéder la particule : \emph{Hand your safety helmet in} (les deux sont corrects).

\textbf{5. Correction du tag :}
\emph{« I am late, \textbf{aren't I}? »}
\emph{Analyse :} Cas particulier de \emph{I am} : le tag négatif standard est \emph{aren't I?}. La forme \emph{amn't I?} n'existe pas en anglais standard ; \emph{am I not?} est correct mais très formel et rarement attendu au Bac.

\textbf{Conclusion :} La maîtrise de l'ordre des mots (embedded questions), de la polarité (tags) et de la séparabilité (phrasal verbs) est cruciale pour l'aisance à l'oral et à l'écrit professionnel.
\end{exoresolu}

\begin{methode}[Analyser un texte sur les besoins en langue seconde au travail -- Reading \& Vocabulary]
\textbf{Objectif :} Lire un texte professionnel, identifier l'idée principale, les arguments, le vocabulaire spécifique et inférer le sens des mots inconnus.
\begin{enumerate}
    \item \textbf{Pré-lecture :} Lire le titre, les sous-titres. Prédire le sujet : \cle{language skills} (compétences linguistiques), \cle{multilingual workplace} (milieu de travail multilingue), \cle{globalization} (mondialisation), \cle{technical documentation} (documentation technique).
    \item \textbf{Lecture active :}
    \begin{itemize}
        \item \cle{Skimming} (lecture rapide) : idée générale (thèse : l'anglais est indispensable pour la carrière technique).
        \item \cle{Scanning} (recherche ciblée) : chiffres (\%), noms d'entreprises, mots-clés (\cle{mandatory} (obligatoire), \cle{asset} (atout), \cle{proficiency} (maîtrise), \cle{technical jargon} (jargon technique)).
    \end{itemize}
    \item \textbf{Gestion du vocabulaire inconnu :} contexte (définition dans la phrase), transparence (mots apparentés au français : \cle{communication, documentation, international, qualification}), famille de mots (\cle{require $\rightarrow$ requirement, required}).
    \item \textbf{Structure argumentative type :} \cle{Context} $\rightarrow$ \cle{Problem/Need} $\rightarrow$ \cle{Solution/Training} $\rightarrow$ \cle{Benefits}.
    \item \textbf{Prise de notes :} utiliser un tableau : \emph{Argument | Preuve/Exemple | Vocabulaire clé}.
\end{enumerate}
\end{methode}

\begin{exoresolu}[Compréhension de texte : « English in the Technical Workplace »]
\textbf{Énoncé :} Lisez l'extrait et répondez aux questions.
\textit{« In today's globalized industry, English has become the \cle{lingua franca} of technology. Most \cle{technical manuals}, \cle{software interfaces} and \cle{safety protocols} are written in English. A recent survey of employers shows that the large majority of technical job offers in Cameroon require at least an \cle{intermediate level} in English. Technicians who \cle{master} this language can \cle{access} international training, \cle{troubleshoot} equipment using original manuals, and \cle{collaborate} with foreign experts. Therefore, anticipating one's language needs is not optional; it is a \cle{career imperative}. »}

\textbf{Questions :}
1. Expliquez avec vos propres mots l'expression « \cle{lingua franca} » à partir du contexte.
2. Citez 3 types de documents techniques mentionnés.
3. Quel niveau d'anglais est exigé par la majorité des offres d'emploi techniques, selon le texte ?
4. Donnez 3 avantages concrets pour un technicien anglophone (vocabulaire du texte).
5. Expliquez l'expression « \cle{career imperative} » dans le contexte.

\tcblower
\textbf{Solution détaillée :}

\textbf{1. Explication de « lingua franca » :}
Une \emph{lingua franca} est une \textbf{langue commune} utilisée pour communiquer entre personnes de langues maternelles différentes. Le contexte le montre : les manuels, logiciels et protocoles sont écrits en anglais, et les techniciens collaborent avec des experts étrangers grâce à cette langue partagée.

\textbf{2. Types de documents techniques :}
\begin{enumerate}
    \item \cle{Technical manuals} (manuels techniques)
    \item \cle{Software interfaces} (interfaces logicielles)
    \item \cle{Safety protocols} (protocoles de sécurité)
\end{enumerate}

\textbf{3. Niveau exigé :}
Selon le texte, la grande majorité des offres d'emploi techniques au Cameroun exige \emph{« at least an intermediate level »}, c'est-à-dire \textbf{au minimum un niveau intermédiaire} en anglais.

\textbf{4. Avantages concrets (vocabulaire du texte) :}
\begin{enumerate}
    \item \cle{Access international training} (accéder à des formations internationales).
    \item \cle{Troubleshoot equipment using original manuals} (dépanner en utilisant les manuels d'origine -- éviter les traductions approximatives).
    \item \cle{Collaborate with foreign experts} (collaborer avec des experts étrangers).
\end{enumerate}

\textbf{5. Explication de « career imperative » :}
\emph{« Career imperative »} = \textbf{impératif de carrière}, nécessité absolue pour la progression professionnelle.
\emph{Justification :} Le connecteur \emph{Therefore} (donc) marque la conclusion logique, et \emph{not optional} (pas optionnel) renforce l'idée d'obligation. Pour un technicien camerounais, ne pas anticiper ce besoin (se former à l'anglais), c'est risquer l'exclusion du marché de l'emploi moderne.

\textbf{Conclusion :} La lecture efficace combine repérage d'informations explicites (listes, niveaux) et inférence lexicale (contexte, connecteurs logiques comme \emph{therefore}).
\end{exoresolu}

\begin{methode}[Rédiger une composition sur les besoins en langue seconde -- Writing]
\textbf{Objectif :} Produire un texte structuré (150-200 mots), argumenté, avec vocabulaire précis, grammaire correcte (temps, modaux, connecteurs) et registre formel/semi-formel.
\begin{enumerate}
    \item \textbf{Analyser le sujet :} identifier le \textbf{thème} (besoins en L2 au travail), le \textbf{format} (composition/essai), la \textbf{cible} (employeur, jury, revue professionnelle), le \textbf{ton} (formel, persuasif).
    \item \textbf{Plan type (3 parties) :}
    \begin{itemize}
        \item \textbf{Introduction (1 paragraphe) :} accroche (contexte de la mondialisation) $\rightarrow$ définition du sujet $\rightarrow$ \cle{thesis statement} (ex : \emph{« Mastering English is a strategic asset for Cameroonian technicians. »}).
        \item \textbf{Développement (2-3 paragraphes) :} 1 idée directrice par paragraphe.
            \begin{itemize}
                \item \cle{Argument 1 : accès à l'information technique} (manuels, normes, logiciels). Connecteurs : \cle{Firstly, Moreover, For instance}.
                \item \cle{Argument 2 : employabilité et mobilité} (marché local/international, certifications, salaires). Connecteurs : \cle{Furthermore, In addition, Consequently}.
                \item \cle{Argument 3 : communication professionnelle} (réunions, emails, rapports, sécurité). Connecteurs : \cle{Finally, Additionally}.
            \end{itemize}
        \item \textbf{Conclusion (1 paragraphe) :} résumé des arguments $\rightarrow$ recommandation/appel à l'action (ex : \emph{« Technical schools should integrate ESP (English for Specific Purposes) modules. »}).
    \end{itemize}
    \item \textbf{Check-list de révision (auto-correction) :}
    \begin{itemize}
        \item \cle{Word count} respecté ? \cle{Paragraphing} visible (saut de ligne) ?
        \item \cle{Tenses} cohérents (Present simple pour les vérités générales, modaux \emph{should/must/can} pour les recommandations) ?
        \item \cle{Vocabulaire technique} utilisé (\cle{ESP, proficiency, technical documentation, troubleshooting, curriculum}) ?
        \item \cle{Connecteurs logiques} variés (pas seulement \emph{and, but, because}) ?
        \item \cle{Orthographe/Ponctuation} (majuscule en début de phrase, points) ?
    \end{itemize}
\end{enumerate}
\end{methode}

\begin{exoresolu}[Rédaction guidée : « Why English matters for my technical career »]
\textbf{Énoncé :} Rédigez une composition de 160 mots environ sur le thème : \emph{« Anticipating one's second-language needs at work is essential for a successful technical career. »} Suivez le plan : Introduction $\rightarrow$ 2 arguments développés $\rightarrow$ Conclusion avec recommandation. Utilisez au moins 5 mots du vocabulaire de l'unité (\cle{lingua franca, technical manuals, proficiency, troubleshooting, employability, mandatory, asset, curriculum}).

\tcblower
\textbf{Solution détaillée :}

\textbf{Proposition de rédaction (modèle corrigé) :}

\textbf{English: A Strategic Asset for Technical Careers}

In today's globalized industry, English has established itself as the \cle{lingua franca} of technology. For a Cameroonian technician, anticipating second-language needs is no longer a simple option but a \cle{career imperative}. Mastering English, specifically \cle{English for Specific Purposes (ESP)}, is a decisive \cle{asset} for professional success.

\cle{Firstly}, technical proficiency depends heavily on access to information. Most \cle{technical manuals}, software interfaces, and international safety standards are published in English. A technician with a good \cle{proficiency} level can \cle{troubleshoot} complex equipment using original documentation, avoiding dangerous translation errors. \cle{For instance}, reading a hydraulic schematic in English allows for faster and safer maintenance operations.

\cle{Furthermore}, \cle{employability} is directly linked to language skills. Many technical job offers now list English as \cle{mandatory}. Bilingual technicians access international training, collaborate with foreign experts, and command higher salaries. The \cle{curriculum} in technical schools must reflect this reality.

\cle{In conclusion}, investing in English learning is investing in one's future. I strongly \cle{recommend} that technical institutions integrate specialized English modules to prepare students for the real demands of the labour market.

\textbf{Analyse du modèle (check-list) :}
\begin{itemize}
    \item \textbf{Mots :} environ 165 mots (dans la fourchette demandée).
    \item \textbf{Structure :} 4 paragraphes distincts (Introduction, Argument 1, Argument 2, Conclusion).
    \item \textbf{Vocabulaire de l'unité (plus de 5 requis) :} \cle{lingua franca, technical manuals, proficiency, troubleshoot, employability, mandatory, asset, curriculum, career imperative, ESP}.
    \item \textbf{Grammaire :} Present Simple (vérités générales), modaux (\emph{can, must}), voix passive (\emph{are published, is linked}).
    \item \textbf{Connecteurs :} \cle{Firstly, For instance, Furthermore, In conclusion}.
    \item \textbf{Registre :} formel, impersonnel par moments (\emph{is established, must reflect}), personnel et engagé à la fin (\emph{I recommend}).
\end{itemize}

\textbf{Conclusion :} Une composition réussie au Bac technique allie structure rigoureuse, vocabulaire de spécialité (ESP) et correction grammaticale. L'entraînement régulier sur des sujets types (sécurité, maintenance, formation) est la clé.
\end{exoresolu}

\begin{attention}[Les 5 erreurs qui coûtent des points au Bac]
\begin{enumerate}
    \item \textbf{Confusion Present Perfect Continuous / Present Perfect Simple :} écrire \emph{« I have been written the report »} au lieu de \emph{« I have written the report »} (résultat) ou \emph{« I have been writing the report »} (processus). \textbf{Règle :} résultat chiffré/achevé $\rightarrow$ Simple ; durée/processus en cours $\rightarrow$ Continuous. Verbes d'état (\emph{know, have, be}) $\rightarrow$ JAMAIS de continuous.
    \item \textbf{Mauvaise formation des Embedded Questions :} garder l'ordre interrogatif ou l'auxiliaire DO. \emph{Incorrect : « Can you tell me what time does the shift start? »} \textbf{Correct : « ... what time the shift starts. »} (ordre affirmatif, pas de DO, pas de point d'interrogation si la phrase principale est déclarative).
    \item \textbf{Question Tags : oubli de l'inversion de polarité ou mauvais pronom.} \emph{Incorrect : « He is a technician, is he? »} \textbf{Correct : « He is a technician, \textbf{isn't he}? »} (tag négatif + pronom sujet). Cas particulier : \emph{I am} $\rightarrow$ \emph{aren't I?}
    \item \textbf{Phrasal Verbs : position de l'objet pronom.} \emph{Incorrect : « Shut down it. »} \textbf{Correct : « Shut \textbf{it} down. »} (l'objet pronom \textbf{doit} s'intercaler). Ne pas confondre \emph{find out} (découvrir une information) et \emph{find} (trouver un objet).
    \item \textbf{Rédaction (Writing) : hors-sujet, absence de paragraphes, vocabulaire trop général.} Écrire « English is good » au lieu de « English is a \cle{mandatory requirement} for reading \cle{technical documentation} ». \textbf{Sanction :} pénalité lourde sur « Contenu » et « Vocabulaire ». Utiliser les mots de la leçon (\cle{proficiency, troubleshooting, employability, asset, curriculum}).
\end{enumerate}
\end{attention}

\newpage

\coursec{Exercices}

\courssub{Je vérifie mes connaissances}

\begin{exercice}[QCM : Job interview vocabulary]
Choose the correct word or expression to complete each sentence. Circle the letter of the best answer.
\begin{enumerate}
    \item Before the interview, you should \_\_\_\_\_\_\_\_\_\_ the company to learn about its activities.
    \begin{enumerate}
        \item[a)] ignore \item[b)] research \item[c)] avoid \item[d)] forget
    \end{enumerate}
    \item A firm \_\_\_\_\_\_\_\_\_\_ makes a good first impression. \emph{(poignée de main ferme)}
    \begin{enumerate}
        \item[a)] handshake \item[b)] handwriting \item[c)] handbag \item[d)] handout
    \end{enumerate}
    \item During the interview, you must \_\_\_\_\_\_\_\_\_\_ your relevant skills. \emph{(mettre en valeur)}
    \begin{enumerate}
        \item[a)] hide \item[b)] highlight \item[c)] lower \item[d)] copy
    \end{enumerate}
    \item Arrive at least 15 minutes \_\_\_\_\_\_\_\_\_\_ for the interview.
    \begin{enumerate}
        \item[a)] late \item[b)] slowly \item[c)] early \item[d)] behind
    \end{enumerate}
    \item At the end, the interviewer may say: ``Do you have any \_\_\_\_\_\_\_\_\_\_ to ask?''
    \begin{enumerate}
        \item[a)] answers \item[b)] questions \item[c)] salaries \item[d)] excuses
    \end{enumerate}
\end{enumerate}
\end{exercice}

\begin{exercice}[Vrai/Faux justifié : Recreational activities]
Read the statements and write \textbf{True} or \textbf{False}. Justify each answer with a short phrase.
\begin{enumerate}
    \item Hiking in the mountains is a recreational activity. \hfill \textbf{True / False}
    \item Watching TV all day is considered a healthy recreational habit. \hfill \textbf{True / False}
    \item Playing football with friends helps reduce stress. \hfill \textbf{True / False}
    \item Recreational activities are only for children. \hfill \textbf{True / False}
    \item Reading a novel can be a recreational activity. \hfill \textbf{True / False}
\end{enumerate}
\end{exercice}

\begin{exercice}[Définitions : Grammar terms]
Match each term (1--5) with its correct definition (a--e).
\begin{center}
\begin{tabularx}{\linewidth}{|l|X|}\hline
\textbf{Term} & \textbf{Definition (letter)} \\ \hline
1. Present Perfect Continuous &  \\ \hline
2. Future Continuous &  \\ \hline
3. Embedded Question &  \\ \hline
4. Question Tag &  \\ \hline
5. Phrasal Verb &  \\ \hline
\end{tabularx}
\end{center}
\begin{enumerate}
    \item[a)] A question inside another question or statement.
    \item[b)] An action that will be in progress at a specific future time.
    \item[c)] A verb combined with a particle (preposition or adverb) changing its meaning.
    \item[d)] A short question added to the end of a statement.
    \item[e)] An action started in the past and continuing now (or just finished).
\end{enumerate}
\end{exercice}

\begin{exercice}[Calculs directs : Tense formation]
Write the correct form of the verb in brackets using the indicated tense.
\begin{enumerate}
    \item She \_\_\_\_\_\_\_\_\_\_ (work) at this company for five years. \emph{(Present Perfect Continuous)}
    \item At 10 a.m. tomorrow, I \_\_\_\_\_\_\_\_\_\_ (attend) a meeting. \emph{(Future Continuous)}
    \item How long \_\_\_\_\_\_\_\_\_\_ you \_\_\_\_\_\_\_\_\_\_ (study) English? \emph{(Present Perfect Continuous)}
    \item Next week, they \_\_\_\_\_\_\_\_\_\_ (prepare) the exhibition. \emph{(Future Continuous)}
    \item We \_\_\_\_\_\_\_\_\_\_ (wait) for the bus for 20 minutes. \emph{(Present Perfect Continuous)}
\end{enumerate}
\end{exercice}

\courssub{Je m'entraîne}

\begin{exercice}[Job interview role-play: Instructions]
Work in pairs. Student A is the interviewer, Student B is the candidate. Use the prompts to give and receive instructions. Then swap roles.
\begin{enumerate}
    \item \textbf{Interviewer:} Ask the candidate to introduce themselves. \emph{(Could you tell me about yourself?)}
    \item \textbf{Candidate:} Respond with name, qualification, experience.
    \item \textbf{Interviewer:} Ask about strengths and weaknesses.
    \item \textbf{Candidate:} Give one strength and one weakness with improvement.
    \item \textbf{Interviewer:} Ask why they want this job.
    \item \textbf{Candidate:} Mention company reputation, career growth.
    \item \textbf{Interviewer:} Invite questions from the candidate.
    \item \textbf{Candidate:} Ask about training opportunities.
\end{enumerate}
\end{exercice}

\begin{exercice}[Listening comprehension: Recreational dialogue]
Listen to the dialogue between Paul and Marie (played twice). Answer the questions in complete sentences.
\begin{enumerate}
    \item What recreational activity does Paul prefer?
    \item How often does Marie go swimming?
    \item Why does Paul not like video games?
    \item What activity do they decide to do together this weekend?
    \item True or False: Marie thinks hiking is boring. Justify.
    \item What time do they plan to meet?
\end{enumerate}
\end{exercice}

\begin{exercice}[Grammar practice: Embedded questions \& Question tags]
\begin{enumerate}
    \item Transform into embedded questions after \emph{Could you tell me...}:
    \begin{enumerate}
        \item Where is the interview room?
        \item What time does the manager arrive?
        \item Does the company offer training?
        \item How many candidates applied?
    \end{enumerate}
    \item Add the correct question tag:
    \begin{enumerate}
        \item You are ready for the interview, \_\_\_\_\_\_\_\_\_\_?
        \item Let's prepare our CVs, \_\_\_\_\_\_\_\_\_\_?
        \item Don't be late, \_\_\_\_\_\_\_\_\_\_?
        \item I am the best candidate, \_\_\_\_\_\_\_\_\_\_?
        \item Nobody called, \_\_\_\_\_\_\_\_\_\_?
    \end{enumerate}
\end{enumerate}
\end{exercice}

\begin{attention}[Pièges fréquents : question tags]
\begin{itemize}
    \item After \emph{I am}, the tag is \textbf{aren't I?} (never ``amn't I'').
    \item After \emph{Let's...}, the tag is \textbf{shall we?}.
    \item After an imperative (\emph{Don't be late}), the tag is \textbf{will you?}.
    \item After \emph{nobody / no one / somebody}, use \textbf{they} in the tag.
    \item In embedded questions, keep the \textbf{statement order}: \emph{Could you tell me where the room is?} (not ``where is the room'').
\end{itemize}
\end{attention}

\begin{exercice}[Phrasal verbs in context]
Complete the paragraph with the correct phrasal verb from the box. Use the right tense and form.
\begin{center}
\begin{tabular}{|c|c|c|c|c|}\hline
fill in & look forward to & get along with & turn up & find out \\ \hline
hand in & take on & run out of & set up & carry out \\ \hline
\end{tabular}
\end{center}
On my first day as an intern at a tech firm in Bafoussam, I had to \_\_\_\_\_\_\_\_\_\_ several forms. I was nervous but the team was friendly and I \_\_\_\_\_\_\_\_\_\_ them well. My supervisor asked me to \_\_\_\_\_\_\_\_\_\_ a new project. I need to \_\_\_\_\_\_\_\_\_\_ more about the software. I must \_\_\_\_\_\_\_\_\_\_ my report by Friday. I \_\_\_\_\_\_\_\_\_\_ learning new skills.
\end{exercice}

\courssub{Je me prépare au Bac}

\begin{exercice}[Bac Grammar: Tense completion --- Technician's career]
Complete the text with the \textbf{Present Perfect Continuous} or \textbf{Future Continuous} form of the verbs in brackets.
\textit{Jean is a maintenance technician in Yaoundé. He \_\_\_\_\_\_\_\_\_\_ (work) at CAMWATER for eight years. Currently, he \_\_\_\_\_\_\_\_\_\_ (upgrade) the pumping station system. By this time next month, he \_\_\_\_\_\_\_\_\_\_ (attend) a specialised training in Douala. He \_\_\_\_\_\_\_\_\_\_ (study) English for two years because he \_\_\_\_\_\_\_\_\_\_ (work) on an international project soon. His colleague, Marie, \_\_\_\_\_\_\_\_\_\_ (prepare) the technical report all morning. Next week, they \_\_\_\_\_\_\_\_\_\_ (inspect) the new pipelines. Jean \_\_\_\_\_\_\_\_\_\_ (apply) for a promotion recently. He hopes that by next year, he \_\_\_\_\_\_\_\_\_\_ (manage) a team.}
\end{exercice}

\begin{exercice}[Probatoire Reading Comprehension: English needs at SONARA]
Read the text and answer the questions.

\textbf{English at the Workplace: A SONARA Employee's Perspective}

Pierre Ngono has worked as a process operator at SONARA's Limbe refinery for twelve years. Recently, the company announced a partnership with an international energy firm from the United Kingdom. This collaboration means that technical manuals, safety protocols, and daily briefings will increasingly be in English. Pierre realised that his basic school English was insufficient for understanding complex technical documents. He enrolled in evening classes to improve his technical vocabulary and listening skills. His goal is to participate confidently in video conferences with British engineers and to read maintenance schematics without translation. Pierre's supervisor encouraged him, noting that operators with strong English skills are prioritised for overseas training programmes. Pierre now spends one hour daily practising with online platforms and colleagues. He says: ``English is no longer just a subject; it is a tool for my career growth and safety at work.''

\begin{enumerate}
    \item How long has Pierre worked at SONARA?
    \item Why does Pierre need to improve his English now?
    \item What specific skills is Pierre focusing on in his evening classes?
    \item What advantage does his supervisor mention for operators with good English?
    \item Find in the text words or expressions that mean:
    \begin{enumerate}
        \item working together with another company (paragraph 1)
        \item not enough (paragraph 2)
        \item diagrams showing how to repair equipment (paragraph 2)
        \item given priority (paragraph 3)
    \end{enumerate}
\end{enumerate}
\end{exercice}

\begin{exercice}[Bac Technique Writing: Composition]
\textbf{Write a composition of 200--250 words on the following topic:}

``The importance of English for your future career as a technical professional.''

\textbf{Guidelines (barème sur 20 points):}
\begin{itemize}
    \item \textbf{Content / Relevance (6 pts):} Clear thesis, specific examples from your technical field, mention of concrete situations (manuals, conferences, safety, promotion, international projects).
    \item \textbf{Organisation (4 pts):} Introduction with hook and thesis, 2--3 well-developed body paragraphs, conclusion.
    \item \textbf{Grammar \& Structures (5 pts):} Correct use of Present Perfect Continuous, Future Continuous, embedded questions, question tags, phrasal verbs, varied sentence structures.
    \item \textbf{Vocabulary (3 pts):} Technical and professional vocabulary (at least 8 relevant terms), appropriate register.
    \item \textbf{Mechanics (2 pts):} Spelling, punctuation, capitalization, paragraphing.
\end{itemize}
\textit{Plan your composition before writing. Use connectors: furthermore, consequently, for instance, in addition, ultimately.}
\end{exercice}

\begin{methode}[How to write your composition]
\begin{enumerate}
    \item \textbf{Brainstorm} (5 min): list 3 concrete situations where English is needed in your trade (reading a manual, talking to a foreign client, attending a training).
    \item \textbf{Plan}: Introduction (hook + thesis) ; Body 1 (English for technical documents and safety) ; Body 2 (English for career growth and international projects) ; Conclusion (personal commitment).
    \item \textbf{Write}: include at least one Present Perfect Continuous (\emph{I have been learning English for...}), one Future Continuous (\emph{I will be using English when...}), one embedded question (\emph{I would like to know whether...}), one question tag and two phrasal verbs (\emph{carry out, look forward to, find out}).
    \item \textbf{Check}: word count (200--250), verb tenses, spelling, paragraphing.
\end{enumerate}
\end{methode}

\newpage

\coursec{Corrigés des exercices}

\coursec{Family and social life}

\courssub{Undergoing a job interview --- Speaking \& Vocabulary}

\begin{corrige}[Exercice 1 --- QCM : Job interview vocabulary]
\begin{enumerate}
    \item \textbf{b) research} --- Before an interview, you must \emph{research} (\cle{se renseigner sur}) the company to know its activities, products and values.
    \item \textbf{a) handshake} --- A \emph{firm handshake} (\cle{une poignée de main ferme}) makes a good first impression. \emph{Handwriting} = écriture ; \emph{handbag} = sac à main ; \emph{handout} = document distribué.
    \item \textbf{b) highlight} --- You must \emph{highlight} (\cle{mettre en valeur}) your relevant skills, not hide them.
    \item \textbf{c) early} --- Arrive at least 15 minutes \emph{early} (\cle{en avance}) ; arriving late creates a bad impression.
    \item \textbf{b) questions} --- The interviewer asks: ``Do you have any \emph{questions} to ask?''
\end{enumerate}
\textbf{Point de vigilance :} ne pas confondre \emph{research} (se renseigner, faire des recherches) et \emph{search} (chercher un objet physique).
\end{corrige}

\begin{corrige}[Exercice 5 --- Job interview role-play (production orale)]
Exercice de production orale : toute réponse correcte et polie est acceptée. Modèle de dialogue attendu :
\begin{itemize}
    \item \textbf{A:} Good morning. Could you tell me about yourself?
    \item \textbf{B:} Good morning. My name is Amina Fotso. I hold a \cle{Baccalauréat technique} in electrical engineering and I did a three-month internship at \cle{ENEO}.
    \item \textbf{A:} What are your strengths and weaknesses?
    \item \textbf{B:} I am hardworking and punctual. My weakness is that I speak English slowly, but I \textbf{have been taking} evening classes to improve.
    \item \textbf{A:} Why do you want this job?
    \item \textbf{B:} Your company has an excellent reputation, and this position will help me grow professionally.
    \item \textbf{A:} Do you have any questions to ask?
    \item \textbf{B:} Yes, could you tell me \textbf{whether} the company offers training opportunities? (\cle{Embedded question})
\end{itemize}
\textbf{Grille d'évaluation orale (10 pts) :} politesse et salutations (2 pts), exactitude grammaticale (2 pts), vocabulaire de l'entretien (2 pts), aisance et prononciation (2 pts), respect des consignes (2 pts).
\end{corrige}

\courssub{Recreational activities --- Listening \& Vocabulary}

\begin{corrige}[Exercice 2 --- Vrai/Faux : Recreational activities]
\begin{enumerate}
    \item \textbf{True} --- Hiking is an outdoor recreational activity (randonnée = loisir de plein air).
    \item \textbf{False} --- Watching TV all day is a sedentary habit; it is not healthy \emph{because} it lacks physical activity.
    \item \textbf{True} --- Playing football releases tension, so it helps reduce stress.
    \item \textbf{False} --- Recreational activities are for everyone: children, adults and elderly people all need leisure.
    \item \textbf{True} --- Reading a novel for pleasure is a recreational (indoor) activity.
\end{enumerate}
\textbf{Point de vigilance :} la justification doit être une phrase courte commençant par \emph{because} ou un tiret explicatif ; une réponse sans justification ne vaut que la moitié des points.
\end{corrige}

\begin{corrige}[Exercice 6 --- Listening comprehension : Recreational dialogue]
Réponses attendues (d'après le script du dialogue) :
\begin{enumerate}
    \item Paul prefers \textbf{hiking} (in the mountains).
    \item Marie goes \textbf{swimming} twice a week.
    \item Paul does not like video games \emph{because} they are sedentary and tiring for the eyes.
    \item They decide to go \textbf{hiking together this weekend}.
    \item \textbf{False} --- Marie thinks hiking is exciting/refreshing, not boring.
    \item They plan to meet at \textbf{7 a.m.} (on Saturday morning).
\end{enumerate}
\textbf{Point de vigilance :} répondre par des phrases complètes (\emph{Paul prefers hiking}), pas par un seul mot ; en compréhension orale au Bac, une réponse juste mais mal formulée perd des points de langue.
\end{corrige}

\courssub{Tenses: Present Perfect Continuous \& Future Continuous}

\begin{aretenir}[Formules des temps clés]
\begin{itemize}
    \item \cle{Present Perfect Continuous} : \textbf{have/has + been + V-ing} \\
    \emph{Action commencée dans le passé et qui se poursuit maintenant (ou vient de s'arrêter). Marqueurs : for, since, all day, recently, how long.}
    \item \cle{Future Continuous} : \textbf{will + be + V-ing} \\
    \emph{Action qui sera en cours à un moment précis du futur. Marqueurs : at this time tomorrow, by next week, next Monday, in two days.}
\end{itemize}
\end{aretenir}

\begin{corrige}[Exercice 3 --- Définitions : Grammar terms]
\begin{center}
\begin{tabularx}{\linewidth}{|l|X|}\hline
\textbf{Term} & \textbf{Definition (letter)} \\ \hline
1. \cle{Present Perfect Continuous} & \textbf{e)} An action started in the past and continuing now. \\ \hline
2. \cle{Future Continuous} & \textbf{b)} An action in progress at a specific future time. \\ \hline
3. \cle{Embedded Question} & \textbf{a)} A question inside another question or statement. \\ \hline
4. \cle{Question Tag} & \textbf{d)} A short question added to the end of a statement. \\ \hline
5. \cle{Phrasal Verb} & \textbf{c)} A verb + particle changing its meaning. \\ \hline
\end{tabularx}
\end{center}
\textbf{Moyen mnémotechnique :} \emph{continuous} = en cours ; \emph{embedded} = enchâssée (imbriquée) ; \emph{tag} = étiquette (petite question ``collée'' à la fin).
\end{corrige}

\begin{corrige}[Exercice 4 --- Calculs directs : Tense formation]
\begin{enumerate}
    \item She \textbf{has been working} at this company for five years. (\emph{for five years} $\rightarrow$ Present Perfect Continuous ; sujet \emph{she} $\rightarrow$ \emph{has}.)
    \item At 10 a.m. tomorrow, I \textbf{will be attending} a meeting. (\emph{at 10 a.m. tomorrow} = moment précis futur $\rightarrow$ Future Continuous.)
    \item How long \textbf{have} you \textbf{been studying} English? (Question : \emph{How long have/has + sujet + been + V-ing}.)
    \item Next week, they \textbf{will be preparing} the exhibition. (\emph{next week} $\rightarrow$ Future Continuous.)
    \item We \textbf{have been waiting} for the bus for 20 minutes. (\emph{for 20 minutes} $\rightarrow$ Present Perfect Continuous.)
\end{enumerate}
\textbf{Formules à retenir :} Present Perfect Continuous : \emph{have/has + been + V-ing} ; Future Continuous : \emph{will + be + V-ing}.
\textbf{Point de vigilance :} ne jamais oublier \emph{been} (PPC) ni \emph{be} (FC) ; vérifier l'accord \emph{have/has} à la 3\textsuperscript{e} personne du singulier.
\end{corrige}

\begin{corrige}[Exercice 9 --- Bac Grammar : Tense completion --- Technician's career]
Texte complété :
\begin{enumerate}
    \item He \textbf{has been working} at \cle{CAMWATER} for eight years. (\emph{for eight years} $\rightarrow$ durée jusqu'à maintenant)
    \item Currently, he \textbf{has been upgrading} the pumping station system. (\emph{currently} + durée implicite $\rightarrow$ PPC pour insister sur la progression)
    \item By this time next month, he \textbf{will be attending} a specialised training in Douala. (\emph{by this time next month} $\rightarrow$ moment futur précis $\rightarrow$ FC)
    \item He \textbf{has been studying} English for two years. (\emph{for two years} $\rightarrow$ durée $\rightarrow$ PPC)
    \item because he \textbf{will be working} on an international project soon. (\emph{soon} / contexte futur $\rightarrow$ FC pour action en cours future)
    \item Marie \textbf{has been preparing} the technical report all morning. (\emph{all morning} $\rightarrow$ durée jusqu'à maintenant $\rightarrow$ PPC)
    \item Next week, they \textbf{will be inspecting} the new pipelines. (\emph{next week} $\rightarrow$ FC)
    \item Jean \textbf{has been applying} for a promotion recently. (\emph{recently} $\rightarrow$ action récente répétée $\rightarrow$ PPC)
    \item By next year, he \textbf{will be managing} a team. (\emph{by next year} $\rightarrow$ moment futur précis $\rightarrow$ FC)
\end{enumerate}
\textbf{Barème indicatif (9 pts) :} 1 pt par forme correcte ; 0,5 pt si le bon temps est choisi mais la forme est fautive (ex. ``has been work'').
\textbf{Points de vigilance :} citer le marqueur temporel qui justifie chaque choix ; ne pas oublier \emph{been} (PPC) ni \emph{be} (FC) ; vérifier l'accord \emph{have/has} (3\textsuperscript{e} pers. sg $\rightarrow$ \emph{has}).
\end{corrige}

\courssub{Embedded questions, Question tags \& Phrasal verbs}

\begin{aretenir}[Règles clés : Embedded Questions]
\begin{itemize}
    \item Ordre \textbf{sujet + verbe} (pas d'inversion).
    \item Pas d'auxiliaire \emph{do/does/did}.
    \item Question fermée $\rightarrow$ \textbf{if / whether}.
    \item Pas de point d'interrogation à l'intérieur, seulement à la fin de la phrase principale.
\end{itemize}
\end{aretenir}

\begin{aretenir}[Règles clés : Question Tags]
\begin{itemize}
    \item Énoncé affirmatif $\rightarrow$ Tag négatif (\emph{aren't you?}, \emph{didn't he?}).
    \item Énoncé négatif (mots : \emph{nobody, never, hardly, seldom}) $\rightarrow$ Tag \textbf{affirmatif} (\emph{did they?}, \emph{does he?}).
    \item Sujet du tag = \textbf{pronom personnel} correspondant au sujet.
    \item Cas spéciaux : \emph{I am} $\rightarrow$ \textbf{aren't I?} ; \emph{Let's} $\rightarrow$ \textbf{shall we?} ; Impératif $\rightarrow$ \textbf{will you?} ; \emph{Don't} $\rightarrow$ \textbf{will you?}
\end{itemize}
\end{aretenir}

\begin{corrige}[Exercice 7 --- Embedded questions \& Question tags]
\textbf{1. Embedded questions} (ordre sujet--verbe, sans inversion, sans point d'interrogation interne) :
\begin{enumerate}
    \item[a)] Could you tell me where the interview room \textbf{is}?
    \item[b)] Could you tell me what time the manager \textbf{arrives}?
    \item[c)] Could you tell me \textbf{if/whether} the company \textbf{offers} training? (question fermée $\rightarrow$ \emph{if/whether})
    \item[d)] Could you tell me how many candidates \textbf{applied}? (pas d'auxiliaire \emph{did} dans la question enchâssée)
\end{enumerate}
\textbf{2. Question tags :}
\begin{enumerate}
    \item[a)] You are ready for the interview, \textbf{aren't you}? (affirmatif $\rightarrow$ tag négatif)
    \item[b)] Let's prepare our CVs, \textbf{shall we}? (cas spécial \emph{Let's})
    \item[c)] Don't be late, \textbf{will you}? (impératif négatif $\rightarrow$ \emph{will you})
    \item[d)] I am the best candidate, \textbf{aren't I}? (cas spécial \emph{I am})
    \item[e)] Nobody called, \textbf{did they}? (\emph{nobody} = sens négatif $\rightarrow$ tag affirmatif avec \emph{they})
\end{enumerate}
\textbf{Points de vigilance :} (1) dans une question enchâssée, on supprime l'inversion et l'auxiliaire \emph{do/does/did} ; (2) un énoncé déjà négatif (\emph{nobody, never, hardly}) exige un tag \textbf{affirmatif} ; (3) \emph{aren't I?} et non \emph{amn't I?}.
\end{corrige}

\begin{corrige}[Exercice 8 --- Phrasal verbs in context]
On my first day as an intern at a tech firm in \cle{Bafoussam}, I had to \textbf{fill in} several forms (\cle{remplir}). I was nervous but the team was friendly and I \textbf{got along with} them well (je me suis bien entendu avec --- passé : \emph{get} $\rightarrow$ \emph{got}). My supervisor asked me to \textbf{take on} a new project (\cle{prendre en charge}). I need to \textbf{find out} more about the software (\cle{me renseigner}). I must \textbf{hand in} my report by Friday (\cle{remettre}). I \textbf{look forward to} learning new skills (\cle{j'attends avec impatience} + V-ing).

\textbf{Points de vigilance :}
\begin{itemize}
    \item \emph{Look forward to} est suivi d'un nom ou d'un \textbf{V-ing} : \emph{I look forward to learning} (jamais ``to learn'').
    \item Accord des temps : le récit est au passé $\rightarrow$ \emph{got along with} ; les obligations présentes restent à la base verbale après \emph{need to / must}.
    \item Verbes non utilisés ici mais à connaître pour le Bac : \emph{turn up} (arriver), \emph{run out of} (être à court de), \emph{set up} (installer), \emph{carry out} (effectuer).
\end{itemize}
\end{corrige}

\courssub{Anticipating second-language needs at work --- Reading \& Vocabulary}

\begin{corrige}[Exercice 10 --- Probatoire Reading Comprehension : English needs at SONARA]
\begin{enumerate}
    \item Pierre has worked at \cle{SONARA} for \textbf{twelve years}. (§1 : ``for twelve years'')
    \item He needs to improve his English because SONARA has announced \textbf{a partnership with a British energy firm}: technical manuals, safety protocols and daily briefings will increasingly be in English, and his basic school English is insufficient for complex technical documents.
    \item In his evening classes, he is focusing on \textbf{technical vocabulary and listening skills}.
    \item His supervisor notes that operators with strong English skills are \textbf{prioritised for overseas training programmes}.
    \item Vocabulaire (mots exacts du texte) :
    \begin{enumerate}
        \item[a)] \textbf{partnership} / \textbf{collaboration} (travailler ensemble avec une autre entreprise)
        \item[b)] \textbf{insufficient} (pas suffisant)
        \item[c)] \textbf{maintenance schematics} (schémas de maintenance/réparation)
        \item[d)] \textbf{prioritised} (prioritaires, ``given priority'')
    \end{enumerate}
\end{enumerate}
\textbf{Barème indicatif (10 pts) :} questions 1--4 : 2 pts chacune (1 pt sens + 1 pt langue) ; question 5 : 0,5 pt par mot.
\textbf{Points de vigilance :} répondre en phrases complètes en s'appuyant sur le texte ; pour la question de vocabulaire, recopier le mot exact du texte, pas un synonyme personnel.
\end{corrige}

\courssub{Writing: Composition on second-language needs at work}

\begin{corrige}[Exercice 11 --- Bac Technique Writing : Composition]
\textbf{Modèle de composition (environ 230 mots) :}

\emph{English: A Key Tool for My Career as a Technician}

\emph{In today's globalised world, English is no longer just a school subject; it is a professional tool. As a future maintenance technician, I am convinced that mastering English will determine my success at work.}

\emph{First of all, most technical documents are written in English. Machine manuals, safety protocols and maintenance schematics rarely exist in French. A technician who cannot read them risks making dangerous mistakes. I \textbf{have been studying} technical English for two years, and I can already understand simple manuals without translation.}

\emph{Furthermore, English opens the door to career growth. Companies like SONARA or CAMWATER work with foreign partners, so employees with good English are prioritised for overseas training. Next year, I \textbf{will be applying} for an internship in an international firm, and I \textbf{will be using} English during video conferences with foreign engineers. I would like to know \textbf{whether} my school offers additional language courses, because I \textbf{look forward to improving} my speaking skills.}

\emph{In conclusion, English is essential for safety, efficiency and promotion in technical careers. It is a tool I must carry with me everywhere, \textbf{isn't it?} I am determined to \textbf{carry out} my studies seriously and to \textbf{find out about} every opportunity to practise, because my professional future depends on it.}

\textbf{Barème indicatif (20 pts) :}
\begin{itemize}
    \item Contenu et pertinence : 6 pts (thèse claire 2 pts ; exemples concrets du domaine technique 3 pts ; situations mentionnées 1 pt).
    \item Organisation : 4 pts (introduction 1 pt ; paragraphes développés 2 pts ; conclusion 1 pt).
    \item Grammaire et structures : 5 pts (au moins un \cle{Present Perfect Continuous}, un \cle{Future Continuous}, une \cle{Embedded Question}, un \cle{Question Tag}, deux \cle{Phrasal Verbs} correctement employés).
    \item Vocabulaire : 3 pts (au moins 8 termes techniques/professionnels, registre adapté).
    \item Mécanique : 2 pts (orthographe, ponctuation, paragraphes, 200--250 mots).
\end{itemize}
\textbf{Points de vigilance :} respecter le nombre de mots (pénalité en dessous de 200) ; vérifier chaque question tag (auxiliaire + pronom) ; ne pas écrire ``I look forward to \emph{improve}'' mais ``to \emph{improving}'' ; utiliser des connecteurs (\emph{furthermore, consequently, in addition, ultimately}).
\end{corrige}

\newpage

\coursec{Fiche bilan}

\courssub{Mindmap of the Lesson}

\begin{popfigure}
\begin{center}\tcbox[colback=popink,colframe=popink,coltext=white,arc=9pt,boxsep=4pt,left=6pt,right=6pt]{\bfseries\small Family \& Social Life (Terminale Tech)}\end{center}
\vspace{-2pt}
\begin{tcbraster}[raster columns=3,raster equal height=rows,raster column skip=6pt,raster row skip=6pt,enhanced,blankest]
\begin{tcolorbox}[enhanced,colback=white,colframe=poporange,boxrule=0.9pt,arc=3pt,title={Job Interview \emph{(Speaking/Vocab)}},fonttitle=\bfseries\scriptsize,coltitle=white,colbacktitle=poporange,left=4pt,right=4pt,top=2pt,bottom=2pt,boxsep=1pt,toptitle=1.5pt,bottomtitle=1.5pt,halign=left]
\begin{itemize}[nosep,leftmargin=*,itemsep=1pt,font=\scriptsize]
\item Giving/Receiving instructions
\item Key expressions
\end{itemize}
\end{tcolorbox}
\begin{tcolorbox}[enhanced,colback=white,colframe=popgreen,boxrule=0.9pt,arc=3pt,title={Recreational Activities \emph{(Listening/Vocab)}},fonttitle=\bfseries\scriptsize,coltitle=white,colbacktitle=popgreen,left=4pt,right=4pt,top=2pt,bottom=2pt,boxsep=1pt,toptitle=1.5pt,bottomtitle=1.5pt,halign=left]
\begin{itemize}[nosep,leftmargin=*,itemsep=1pt,font=\scriptsize]
\item Leisure vocabulary
\item Comprehension
\end{itemize}
\end{tcolorbox}
\begin{tcolorbox}[enhanced,colback=white,colframe=poppurple,boxrule=0.9pt,arc=3pt,title={Grammar: Tenses \emph{(Revision + New)}},fonttitle=\bfseries\scriptsize,coltitle=white,colbacktitle=poppurple,left=4pt,right=4pt,top=2pt,bottom=2pt,boxsep=1pt,toptitle=1.5pt,bottomtitle=1.5pt,halign=left]
\begin{itemize}[nosep,leftmargin=*,itemsep=1pt,font=\scriptsize]
\item Present Perfect Continuous
\item Future Continuous
\end{itemize}
\end{tcolorbox}
\begin{tcolorbox}[enhanced,colback=white,colframe=poppink,boxrule=0.9pt,arc=3pt,title={Grammar: Structures \emph{(Embedded/Tags/Phrasal)}},fonttitle=\bfseries\scriptsize,coltitle=white,colbacktitle=poppink,left=4pt,right=4pt,top=2pt,bottom=2pt,boxsep=1pt,toptitle=1.5pt,bottomtitle=1.5pt,halign=left]
\begin{itemize}[nosep,leftmargin=*,itemsep=1pt,font=\scriptsize]
\item Embedded Questions
\item Question Tags
\item Phrasal Verbs
\end{itemize}
\end{tcolorbox}
\begin{tcolorbox}[enhanced,colback=white,colframe=popgold,boxrule=0.9pt,arc=3pt,title={L2 Needs at Work \emph{(Reading/Vocab)}},fonttitle=\bfseries\scriptsize,coltitle=white,colbacktitle=popgold,left=4pt,right=4pt,top=2pt,bottom=2pt,boxsep=1pt,toptitle=1.5pt,bottomtitle=1.5pt,halign=left]
\begin{itemize}[nosep,leftmargin=*,itemsep=1pt,font=\scriptsize]
\item Anticipating needs
\item Workplace context
\end{itemize}
\end{tcolorbox}
\begin{tcolorbox}[enhanced,colback=white,colframe=popdark,boxrule=0.9pt,arc=3pt,title={Writing \emph{(Composition)}},fonttitle=\bfseries\scriptsize,coltitle=white,colbacktitle=popdark,left=4pt,right=4pt,top=2pt,bottom=2pt,boxsep=1pt,toptitle=1.5pt,bottomtitle=1.5pt,halign=left]
\begin{itemize}[nosep,leftmargin=*,itemsep=1pt,font=\scriptsize]
\item Structured essay
\item Argumentation
\end{itemize}
\end{tcolorbox}
\end{tcbraster}

\legende{Mindmap of Lesson 01: Family and Social Life (Terminale Technique).}
\end{popfigure}

\courssub{Bilan des savoirs et savoir-faire}

\begin{bilan}

\textbf{1. Speaking \& Vocabulary : Undergoing a Job Interview}
\begin{itemize}
  \item \textbf{Function:} Giving and receiving instructions for a job interview.
  \item \textbf{Key Vocabulary (Mots-clés) :} \cle{application letter} (lettre de motivation), \cle{CV/résumé}, \cle{interviewer} / \cle{interviewee}, \cle{dress code} (code vestimentaire), \cle{punctuality} (ponctualité), \cle{background} (parcours), \cle{strengths/weaknesses} (forces/faiblesses), \cle{salary expectations} (prétentions salariales), \cle{referee} (référent).
  \item \textbf{Useful Expressions (Expressions utiles) :}
  \begin{itemize}
    \item \emph{``Could you tell me about yourself?''} (Parlez-moi de vous.)
    \item \emph{``What are your strengths?''} (Quels sont vos atouts ?)
    \item \emph{``I have experience in...''} (J'ai de l'expérience dans...)
    \item \emph{``When will I hear from you?''} (Quand aurai-je des nouvelles ?)
  \end{itemize}
\end{itemize}

\textbf{2. Listening \& Vocabulary : Recreational Activities}
\begin{itemize}
  \item \textbf{Context:} Leisure, hobbies, sports, cultural events.
  \item \textbf{Key Vocabulary :} \cle{outdoor/indoor activities} (activités de plein air / intérieures), \cle{hobby} (passe-temps), \cle{unwind} (se détendre), \cle{team building} (cohésion d'équipe), \cle{work-life balance} (équilibre vie pro/vie perso), \cle{spectator} (spectateur), \cle{participant}.
\end{itemize}

\textbf{3. Grammar : Tenses (Revision + New Forms)}

\begin{center}
\begin{tabularx}{\linewidth}{|l|X|X|}\hline
\rowcolor{popblueL}
\textbf{Tense / Temps} & \textbf{Structure / Structure} & \textbf{Usage / Emploi (Exemples)} \\ \hline
\textbf{Present Perfect Continuous} & \textbf{have/has + been + V-ing} & Action started in past, continuing now or just stopped with visible result. \\
\emph{Présent parfait continu} & \emph{I have been learning English for 3 years.} & \emph{Je'apprends l'anglais depuis 3 ans (et je continue).} \\ \hline
\textbf{Future Continuous} & \textbf{will + be + V-ing} & Action in progress at a specific future time. \\
\emph{Futur continu (en \textsc{be} + \textsc{ing})} & \emph{This time tomorrow, I will be sitting for my exam.} & \emph{Demain à cette heure, je serai en train de passer mon examen.} \\ \hline
\textbf{Revision (Révision)} & Present Simple/Continuous, Past Simple/Continuous, Present Perfect Simple, Future Simple (will/going to). & Narrating, describing habits, finished past actions, predictions. \\ \hline
\end{tabularx}
\end{center}

\textbf{4. Grammar : Complex Structures}

\begin{definition}[Embedded Questions / Questions indirectes imbriquées]
Used after introductory phrases (\emph{Do you know...}, \emph{I wonder...}, \emph{Can you tell me...}). \textbf{Word order: Subject + Verb (NO inversion, NO auxiliary do/does/did).}
\begin{itemize}
  \item Direct: \emph{``Where does he live?''} $\rightarrow$ Embedded: \emph{``Do you know \textbf{where he lives}?''}
  \item Yes/No questions $\rightarrow$ use \textbf{if/whether}: \emph{``Is she ready?''} $\rightarrow$ \emph{``I wonder \textbf{if she is ready}.''}
\end{itemize}
\end{definition}

\begin{definition}[Question Tags / Questions d'appui]
Short question at end of statement. \textbf{Rule: Positive statement $\rightarrow$ Negative tag; Negative statement $\rightarrow$ Positive tag. Same auxiliary/modal. Pronoun matches subject.}
\begin{center}
\begin{tabular}{ll}
\emph{You are ready, \textbf{aren't you}?} & (Pos $\rightarrow$ Neg) \\
\emph{He hasn't come, \textbf{has he}?} & (Neg $\rightarrow$ Pos) \\
\emph{Let's go, \textbf{shall we}?} & (Imperative \emph{let's}) \\
\emph{Don't forget, \textbf{will you}?} & (Imperative negative)
\end{tabular}
\end{center}
\end{definition}

\begin{definition}[Phrasal Verbs \& Prepositional Phrases / Verbes à particules \& Locutions prépositionnelles]
Verb + Particle (adverb/preposition) changing meaning. Often informal.
\begin{itemize}
  \item \cle{fill in} (remplir -- a form), \cle{look for} (chercher), \cle{look forward to + V-ing} (avoir hâte de), \cle{get along with} (s'entendre avec), \cle{run out of} (manquer de), \cle{put up with} (supporter), \cle{turn down} (refuser), \cle{apply for} (postuler à).
\end{itemize}
\end{definition}

\textbf{5. Reading \& Vocabulary : Anticipating Second-Language Needs at Work}
\begin{itemize}
  \item \textbf{Context:} Professional environments requiring English/French bilingualism (Cameroon context).
  \item \textbf{Key Concepts:} \cle{bilingual proficiency} (compétence bilingue), \cle{technical jargon} (jargon technique), \cle{business correspondence} (correspondance commerciale), \cle{negotiation skills} (compétences en négociation), \cle{report writing} (rédaction de rapports), \cle{soft skills} (savoir-être).
\end{itemize}

\textbf{6. Writing : Composition on Second-Language Needs}
\begin{methode}[Structuring a Composition / Structurer une rédaction]
\begin{enumerate}
  \item \textbf{Analyse the prompt:} Identify topic (needs), audience (employer/school), format (essay/letter/report).
  \item \textbf{Outline:} Introduction (Hook + Thesis statement), Body Paragraphs (1 idea = 1 paragraph: Argument + Example/Linker), Conclusion (Restate thesis + Final thought/Recommendation).
  \item \textbf{Linkers (Connecteurs) :} \emph{Furthermore, Moreover, In addition} (Addition); \emph{However, On the other hand} (Contrast); \emph{Consequently, Therefore} (Consequence); \emph{For instance, Such as} (Example).
  \item \textbf{Proofread:} Check tenses, subject-verb agreement, spelling, word count.
\end{enumerate}
\end{methode}

\end{bilan}

\courssub{Checklist avant le Bac}

\begin{aretenir}[Checklist avant le Bac -- Lesson 01]
\begin{itemize}
  \item $\square$ I can simulate a job interview (giving/receiving instructions) using appropriate register.
  \item $\square$ I know key vocabulary for job interviews (CV, strengths, weaknesses, dress code, referees).
  \item $\square$ I understand listening texts about recreational activities and can extract specific info.
  \item $\square$ I master the form and use of \textbf{Present Perfect Continuous} (have/has been + V-ing).
  \item $\square$ I master the form and use of \textbf{Future Continuous} (will be + V-ing).
  \item $\square$ I can transform direct questions into \textbf{Embedded Questions} (correct word order: S+V).
  \item $\square$ I can form correct \textbf{Question Tags} (polarity switch, pronoun/auxiliary match).
  \item $\square$ I can use 10+ common \textbf{Phrasal Verbs} in context (fill in, look forward to, apply for...).
  \item $\square$ I can read a text on L2 needs at work and answer comprehension questions.
  \item $\square$ I know vocabulary for workplace language needs (jargon, correspondence, negotiation, reports).
  \item $\square$ I can write a structured composition (Intro/Body/Conclusion) with linkers on L2 needs.
  \item $\square$ I can integrate all skills in a realistic professional scenario (Integration Activity).
\end{itemize}
\end{aretenir}

\findecours

\end{document}
